% Application cumulée de l’éthique et de l’empathie — Informatique 3ème — Cours signé OnBuch+
% Programme officiel MINESEC. Compiler avec : tectonic N11-application-cumulee-ethique-empathie.tex
\newcommand{\DOCMATIERE}{Informatique}
\newcommand{\DOCNIVEAU}{3ème}
\newcommand{\DOCMODULE}{Informatique – classe de 3ème}
\newcommand{\DOCLECON}{N11}
\newcommand{\DOCTITRE}{Application cumulée de l’éthique et de l’empathie}
% OnBuch+ — préambule « cours signés » ENRICHI (compilé avec tectonic / XeLaTeX)
% Système visuel « Pop » d'OnBuch+ : fond crème (jamais blanc), bordures encre
% épaisses, ombres « dures » décalées, titres Archivo Black en capitales.
% Version MATHÉMATIQUES : graphes (pgfplots/tikz), boîtes pédagogiques
% (définition, propriété, méthode, attention, le savais-tu, exercice résolu,
% exercice, TP, bilan), page de garde et signature OnBuch+ ancrée en bas.
%
% Macros attendues dans main.tex AVANT \input{preamble} :
%   \DOCMATIERE  \DOCNIVEAU  \DOCMODULE  \DOCLECON (numéro)  \DOCTITRE
\documentclass[11pt]{article}

\usepackage{fontspec}
\usepackage[french]{babel}
\usepackage[a4paper,margin=2cm,top=3.4cm,bottom=2.8cm,headheight=1.4cm,footskip=1.3cm]{geometry}
\usepackage{amsmath,amssymb,amsfonts}
\usepackage{mathtools} % \xmapsto, \coloneqq... souvent utilisés par les modèles
\usepackage{cancel} % simplifications barrées (bilans de forces)
% \abs{z} (module, valeur absolue) et \norm{u} (norme), utilisés par les modèles
\providecommand{\abs}{}\let\abs\relax\DeclarePairedDelimiter{\abs}{\lvert}{\rvert}
\providecommand{\norm}{}\let\norm\relax\DeclarePairedDelimiter{\norm}{\lVert}{\rVert}
\usepackage[table]{xcolor} % [table] : fournit aussi \rowcolors (alternance de lignes)
\usepackage{pagecolor}
\usepackage{tikz}
\usetikzlibrary{arrows.meta,positioning,calc,shapes.geometric,shapes.misc,shapes.arrows,shapes.symbols,decorations.pathmorphing,decorations.pathreplacing,decorations.markings,patterns,shadows,mindmap,trees,fit,backgrounds,matrix,intersections,angles,quotes}
\usepackage{pgfplots}
\usepackage[version=4]{mhchem} % équations nucléaires \ce{^{235}_{92}U}
\usepackage[european,siunitx]{circuitikz} % schémas électriques
\pgfplotsset{compat=1.18}
\usepgfplotslibrary{fillbetween,groupplots} % aires entre courbes (calcul intégral)
\usepackage{enumitem}
\usepackage{fancyhdr}
% [most] charge toutes les bibliothèques tcolorbox (listings, minted, raster,
% documentation...) et gonfle inutilement la mémoire TeX sur un document long
% (~300 boîtes) : on ne charge que ce qui est réellement utilisé.
\usepackage[skins,breakable]{tcolorbox}
\usepackage{graphicx}
\usepackage{colortbl}
\usepackage{booktabs}
\usepackage{tabularx}
\usepackage{diagbox} % case d'en-tête diagonale des tableaux à double entrée
% Colonnes extensibles centrée / gauche / droite (C, L, R), souvent utilisées par les modèles
\newcolumntype{C}{>{\centering\arraybackslash}X}
\newcolumntype{L}{>{\raggedright\arraybackslash}X}
\newcolumntype{R}{>{\raggedleft\arraybackslash}X}
\usepackage{array}
\usepackage{multirow}
\usepackage{multicol}
\usepackage{siunitx}
% parse-numbers=false : accepte aussi \SI{2,5 \times 10^{-3}}{...}
\sisetup{locale=FR,output-decimal-marker={,},group-minimum-digits=4,parse-numbers=false}
\DeclareSIUnit\ohm{\ensuremath{\symup{\Omega}}} % U+2126 absent de la police
\DeclareSIUnit\division{div} % réglages d'oscilloscope (ms/div, V/div)
\DeclareSIUnit\tr{tr} % tours (vitesse de rotation, ex. \SI{1500}{\tr\per\minute})
\DeclareSIUnit\bit{bit}
\DeclareSIUnit\byte{o} % octet
\DeclareSIUnit\inch{po} % pouce
\DeclareSIUnit\ppp{ppp} % points par pouce (résolution d'image)
\usepackage{qrcode}
% Blocs de code source (PHP, C, HTML, JavaScript, SQL) — spécifique à la
% matière Informatique : listings + habillage aux couleurs OnBuch+ (voir le
% style \lstset ci-dessous, à côté des autres boîtes pédagogiques).
\usepackage{listings}
% \degree : commande fréquemment utilisée par les modèles pour un angle ou une
% température en dehors de \SI{}{} (non fournie par les paquets chargés).
\providecommand{\degree}{^{\circ}}
% \qed (fin de démonstration), utilisé par les modèles sans amsthm
\providecommand{\qed}{\hfill$\square$}
% \barre : double barre des valeurs interdites dans un tableau de variations (tkz-tab)
\providecommand{\barre}{\ensuremath{\|}}
% environnement proof (amsthm non chargé) utilisé par les modèles
\ifdefined\proof\else\newenvironment{proof}[1][Démonstration]{\par\noindent\textit{#1.}\ }{\qed\par}\fi
\usepackage{lastpage}
\usepackage{hyperref}
\hypersetup{hidelinks}

% ============================================================
% Palette « Pop » OnBuch+ — mêmes hex que la classe P (plus_ui.dart)
% ============================================================
\definecolor{poporange}{HTML}{F59321}
\definecolor{popdark}{HTML}{E07A0C}
\definecolor{popcream}{HTML}{FFF6E8}
\definecolor{popink}{HTML}{1C1714}
\definecolor{poppurple}{HTML}{7A5AE0}
\definecolor{popgreen}{HTML}{2E9E6A}
\definecolor{poppink}{HTML}{E0457B}
\definecolor{popblue}{HTML}{2F7DE1}
\definecolor{popgold}{HTML}{C98A12}
\definecolor{popmuted}{HTML}{8A7F76}
\definecolor{popline}{HTML}{EADFD2}
\definecolor{popgray}{HTML}{8A7F76}
\colorlet{popgrey}{popgray}
% Alias tolérants (le modèle nomme parfois les couleurs différemment)
\colorlet{popwhite}{white}
\colorlet{popblack}{popink}
\colorlet{popred}{poppink}
\colorlet{popyellow}{popgold}
% Teintes claires pour fonds de tableaux / zones
\colorlet{popblueL}{popblue!10!white}
\colorlet{poporangeL}{poporange!14!white}
\colorlet{popgreenL}{popgreen!12!white}
\colorlet{poppurpleL}{poppurple!12!white}
\colorlet{poppinkL}{poppink!12!white}
\colorlet{popgoldL}{popgold!14!white}

\pagecolor{popcream}

% ============================================================
% Typographie
% ============================================================
\defaultfontfeatures{Ligatures=TeX}
\setmainfont{PlusJakartaSans-Medium.ttf}[Path=../../../fonts/,BoldFont=PlusJakartaSans-Bold.ttf]
% Maths en sans-serif (Fira Math) avec chiffres et lettres droites de Plus
% Jakarta, pour que les formules se fondent dans le texte.
\usepackage[math-style=ISO,bold-style=ISO]{unicode-math}
\setmathfont{FiraMath-Regular.otf}[Path=../../../fonts/]
\setmathfont{PlusJakartaSans-Medium.ttf}[Path=../../../fonts/,range={up/num,up/{Latin,latin},\mathup}]
\usepackage{icomma} % virgule décimale sans espace en mode math
% Caractères mathématiques Unicode absents des polices de texte (Plus Jakarta,
% Archivo Black) : rendus par leur équivalent mathématique, en texte comme en
% formule — sinon ils s'affichent en carré vide (ex. « Arithmétique dans ℤ »).
\usepackage{newunicodechar}
\newunicodechar{ℝ}{\ensuremath{\mathbb{R}}}
\newunicodechar{ℤ}{\ensuremath{\mathbb{Z}}}
\newunicodechar{ℂ}{\ensuremath{\mathbb{C}}}
\newunicodechar{ℕ}{\ensuremath{\mathbb{N}}}
\newunicodechar{ℚ}{\ensuremath{\mathbb{Q}}}
\newunicodechar{𝔻}{\ensuremath{\mathbb{D}}}
\newunicodechar{α}{\ensuremath{\alpha}}
\newunicodechar{β}{\ensuremath{\beta}}
\newunicodechar{γ}{\ensuremath{\gamma}}
\newunicodechar{δ}{\ensuremath{\delta}}
\newunicodechar{ε}{\ensuremath{\varepsilon}}
\newunicodechar{θ}{\ensuremath{\theta}}
\newunicodechar{λ}{\ensuremath{\lambda}}
\newunicodechar{μ}{\ensuremath{\mu}}
\newunicodechar{σ}{\ensuremath{\sigma}}
\newunicodechar{τ}{\ensuremath{\tau}}
\newunicodechar{φ}{\ensuremath{\varphi}}
\newunicodechar{ψ}{\ensuremath{\psi}}
\newunicodechar{ω}{\ensuremath{\omega}}
\newunicodechar{Σ}{\ensuremath{\Sigma}}
% majuscules produites par \MakeUppercase dans les titres (α-aminés -> Α)
\newunicodechar{Α}{\ensuremath{\alpha}}
\newunicodechar{Β}{\ensuremath{\beta}}
\newunicodechar{Γ}{\ensuremath{\Gamma}}
\newunicodechar{Δ}{\ensuremath{\Delta}}
\newunicodechar{Ε}{\ensuremath{\varepsilon}}
\newunicodechar{Θ}{\ensuremath{\Theta}}
\newunicodechar{Λ}{\ensuremath{\Lambda}}
\newunicodechar{Μ}{\ensuremath{\mu}}
\newunicodechar{Τ}{\ensuremath{\tau}}
\newunicodechar{Φ}{\ensuremath{\Phi}}
\newunicodechar{Ψ}{\ensuremath{\Psi}}
\newunicodechar{Ω}{\ensuremath{\Omega}}
\newunicodechar{↦}{\ensuremath{\mapsto}}
\newunicodechar{∈}{\ensuremath{\in}}
\newunicodechar{∉}{\ensuremath{\notin}}
\newunicodechar{∧}{\ensuremath{\wedge}}
\newunicodechar{⇔}{\ensuremath{\Leftrightarrow}}
\newunicodechar{⟺}{\ensuremath{\Longleftrightarrow}}
\newunicodechar{⇒}{\ensuremath{\Rightarrow}}
\newunicodechar{⟹}{\ensuremath{\Longrightarrow}}
\newunicodechar{∘}{\ensuremath{\circ}}
\newunicodechar{∪}{\ensuremath{\cup}}
\newunicodechar{∩}{\ensuremath{\cap}}
\newunicodechar{≡}{\ensuremath{\equiv}}
\newunicodechar{⊂}{\ensuremath{\subset}}
\newunicodechar{⊥}{\ensuremath{\perp}}
\newunicodechar{∃}{\ensuremath{\exists}}
\newunicodechar{∀}{\ensuremath{\forall}}
\newunicodechar{∛}{\ensuremath{\sqrt[3]{\,}}}
\newunicodechar{∜}{\ensuremath{\sqrt[4]{\,}}}
\newunicodechar{⃗}{\ensuremath{{}^{\to}}}
\newunicodechar{ⁿ}{\ifmmode^{n}\else\textsuperscript{\ensuremath{n}}\fi}
\newunicodechar{ˣ}{\ifmmode^{x}\else\textsuperscript{\ensuremath{x}}\fi}
\newunicodechar{ᵇ}{\ifmmode^{b}\else\textsuperscript{\ensuremath{b}}\fi}
\newunicodechar{ʳ}{\ifmmode^{r}\else\textsuperscript{\ensuremath{r}}\fi}
\newunicodechar{ᵘ}{\ifmmode^{u}\else\textsuperscript{\ensuremath{u}}\fi}
\newunicodechar{ʸ}{\ifmmode^{y}\else\textsuperscript{\ensuremath{y}}\fi}
\newunicodechar{ᵃ}{\ifmmode^{a}\else\textsuperscript{\ensuremath{a}}\fi}
\newunicodechar{ᵉ}{\ifmmode^{e}\else\textsuperscript{\ensuremath{e}}\fi}
\newunicodechar{ᵏ}{\ifmmode^{k}\else\textsuperscript{\ensuremath{k}}\fi}
\newunicodechar{ᵖ}{\ifmmode^{p}\else\textsuperscript{\ensuremath{p}}\fi}
\newunicodechar{ᴺ}{\ifmmode^{N}\else\textsuperscript{\ensuremath{N}}\fi}
\newunicodechar{ᶜ}{\ifmmode^{c}\else\textsuperscript{\ensuremath{c}}\fi}
\newunicodechar{ʰ}{\ifmmode^{h}\else\textsuperscript{\ensuremath{h}}\fi}
\newunicodechar{ᵠ}{\ifmmode^{q}\else\textsuperscript{\ensuremath{q}}\fi}
\newunicodechar{ₙ}{\ifmmode_{n}\else\textsubscript{\ensuremath{n}}\fi}
\newunicodechar{ₐ}{\ifmmode_{a}\else\textsubscript{\ensuremath{a}}\fi}
\newunicodechar{ᵢ}{\ifmmode_{i}\else\textsubscript{\ensuremath{i}}\fi}
\newunicodechar{ᵣ}{\ifmmode_{r}\else\textsubscript{\ensuremath{r}}\fi}
\newunicodechar{ₖ}{\ifmmode_{k}\else\textsubscript{\ensuremath{k}}\fi}
\newunicodechar{ᵦ}{\ifmmode_{\beta}\else\textsubscript{\ensuremath{\beta}}\fi}
\newunicodechar{ₓ}{\ifmmode_{x}\else\textsubscript{\ensuremath{x}}\fi}
\newfontfamily\popdisplay{ArchivoBlack-Regular.ttf}[Path=../../../fonts/]
\newfontfamily\poplogo{SpaceGrotesk-Bold.ttf}[Path=../../../fonts/]
\newfontfamily\popsemi{PlusJakartaSans-SemiBold.ttf}[Path=../../../fonts/]
\newfontfamily\popxbold{PlusJakartaSans-ExtraBold.ttf}[Path=../../../fonts/]
\newfontfamily\popxboldls{PlusJakartaSans-ExtraBold.ttf}[Path=../../../fonts/,LetterSpace=3.0]

\setlist[enumerate]{leftmargin=1.7em,itemsep=5pt,topsep=4pt}
\setlist[itemize]{leftmargin=1.7em,itemsep=4pt,topsep=4pt,label={\color{poporange}\textbullet}}
\setlength{\parindent}{0pt}
\setlength{\parskip}{5pt}
\renewcommand{\arraystretch}{1.25}

% pgfplots : style Pop par défaut
\pgfplotsset{
  every axis/.append style={axis line style={popink,line width=1pt},tick style={popink},
    grid style={popline,line width=0.5pt},label style={font=\small},tick label style={font=\footnotesize}},
  popcurve/.style={poporange,line width=1.8pt,no markers,smooth},
}
% Même style disponible dans un \draw TikZ simple (hors pgfplots)
\tikzset{popcurve/.style={poporange,line width=1.8pt}}
\tikzset{popgrid/.style={popline,very thin}}
\colorlet{popcurve}{poporange} % le nom du style est parfois utilisé comme couleur

% ============================================================
% Pill / squiggle
% ============================================================
\newcommand{\poppill}[3]{%
  \tikz[baseline=-0.55ex]\node[draw=popink,line width=1.2pt,rounded corners=2.6pt,%
    fill=#2,inner xsep=6.5pt,inner ysep=3pt]{%
    \popxboldls\fontsize{7}{7}\selectfont\color{#3}\MakeUppercase{#1}};%
}
\newcommand{\popsquiggle}[1]{%
  \tikz[baseline=-0.4ex]{
    \draw[#1,line width=2.6pt,line cap=round]
      plot [smooth] coordinates {(0,0.09) (0.34,0.01) (0.68,0.21) (1.02,0.01) (1.36,0.10)};
  }%
}
% Mot-clé surligné (marqueur orange sous le texte)
\newcommand{\cle}[1]{{\popxbold\color{popdark}#1}}

% ============================================================
% En-tête / pied
% ============================================================
\pagestyle{fancy}
\fancyhf{}
\renewcommand{\headrulewidth}{0pt}
\renewcommand{\footrulewidth}{0pt}
\lhead{\raisebox{-0.15ex}{\poplogo\fontsize{13}{13}\selectfont\color{popink}OnBuch\color{poporange}+}}
\rhead{\poppill{\DOCMATIERE\ \textbullet\ \DOCNIVEAU\ \textbullet\ Leçon \DOCLECON}{white}{popink}}
\lfoot{\popxbold\fontsize{7.6}{7.6}\selectfont\color{popmuted}\MakeUppercase{Cours signé OnBuch+}\ \textbullet\ {\popsemi\DOCTITRE}}
\rfoot{\tikz[baseline=-0.6ex]\node[rounded corners=8.5pt,fill=popink,minimum height=17pt,minimum width=17pt,inner xsep=5pt,inner sep=0]{\popxbold\fontsize{8}{8}\selectfont\color{popcream}\thepage\,/\,\pageref*{LastPage}};}

% ============================================================
% Titres
% ============================================================
\newcommand{\chaptitre}[2]{%
  \begin{tcolorbox}[enhanced,colback=poporange,colframe=popink,boxrule=2.6pt,arc=11pt,
      drop shadow=popink,left=15pt,right=15pt,top=12pt,bottom=11pt,before skip=3pt,after skip=18pt]
    \poppill{#2}{popink}{popcream}\par\vspace{9pt}
    {\raggedright\hyphenpenalty=10000\exhyphenpenalty=10000\popdisplay\fontsize{22}{24}\selectfont\color{popcream}\MakeUppercase{#1}\par}
  \end{tcolorbox}
}
% Section numérotée : pastille orange avec le numéro + titre Archivo
\newcounter{popsec}
\newcommand{\coursec}[1]{%
  \stepcounter{popsec}\setcounter{popsub}{0}%
  \par\vspace{16pt}\noindent
  \begin{minipage}{\linewidth}
  \tikz[baseline=-0.7ex]\node[draw=popink,line width=1.6pt,fill=poporange,rounded corners=4pt,minimum size=22pt,inner sep=0]{\popdisplay\fontsize{12}{12}\selectfont\color{popcream}\thepopsec};%
  \hspace{8pt}{\popdisplay\fontsize{15}{17}\selectfont\color{popink}\MakeUppercase{#1}}\par
  \vspace{4pt}\noindent{\color{poporange}\rule{36pt}{3.6pt}}
  \end{minipage}\par\nopagebreak\vspace{8pt}
}
\newcounter{popsub}
\newcommand{\courssub}[1]{%
  \stepcounter{popsub}%
  \par\vspace{9pt}\noindent{\popxbold\fontsize{11.5}{13}\selectfont\color{popink}{\color{poporange}\thepopsec.\thepopsub}\hspace{6pt}#1}\par\nopagebreak\vspace{4pt}
}

% ============================================================
% Boîtes pédagogiques — même recette : fond blanc, bordure épaisse colorée,
% ombre dure encre, badge plein. Titre optionnel [#1].
% ============================================================
\tcbset{popbase/.style={breakable,enhanced,colback=white,boxrule=2.3pt,arc=9pt,
  shadow={3pt}{-3pt}{0pt}{popink},left=14pt,right=14pt,top=11pt,bottom=11pt,before skip=11pt,after skip=15pt,
  fontupper=\popsemi\fontsize{10.6}{14.5}\selectfont\color{popink},
  fontlower=\popsemi\fontsize{10.6}{14.5}\selectfont\color{popink}}}
% Coche dessinée (le glyphe ✓ n'existe pas dans les polices du document)
\newcommand{\popcheck}[1]{\tikz[baseline=-0.5ex]\draw[#1,line width=1.6pt,line cap=round,line join=round](-0.13,0.0)--(-0.04,-0.09)--(0.13,0.1);}
\providecommand{\coche}{\popcheck{popgreen}}
\providecommand{\resultat}[1]{\par\smallskip\noindent\fcolorbox{popgreen}{popgreenL}{\parbox{\dimexpr\linewidth-2\fboxsep-2\fboxrule\relax}{\textbf{Résultat.} #1}}\par\smallskip}
\newcommand{\popboxhead}[3]{\poppill{#1}{#2}{white}\ifx\relax#3\relax\else\hspace{6pt}{\popxbold\fontsize{10.6}{12}\selectfont\color{popink}#3}\fi\par\vspace{7pt}}

\newtcolorbox{aretenir}[1][]{popbase,colframe=popgreen,before upper={\popboxhead{À retenir}{popgreen}{#1}}}
\newtcolorbox{exemplebox}[1][]{popbase,colframe=poporange,before upper={\popboxhead{Exemple}{poporange}{#1}}}
\newtcolorbox{definition}[1][]{popbase,colframe=popblue,before upper={\popboxhead{Définition}{popblue}{#1}}}
\newtcolorbox{propriete}[1][]{popbase,colframe=poppurple,before upper={\popboxhead{Propriété}{poppurple}{#1}}}
\newtcolorbox{methode}[1][]{popbase,colframe=popgold,before upper={\popboxhead{Méthode}{popgold}{#1}}}
\newtcolorbox{attention}[1][]{popbase,colframe=poppink,before upper={\popboxhead{Attention}{poppink}{#1}}}
\newtcolorbox{experience}[1][]{popbase,colframe=popblue,colback=popblueL,before upper={\popboxhead{Expérience}{popblue}{#1}}}
% « Le savais-tu ? » : carte violette pleine (contexte, culture, Cameroun)
\newtcolorbox{savaistu}[1][]{popbase,colback=poppurple,colframe=popink,
  fontupper=\popsemi\fontsize{10.4}{14.2}\selectfont\color{white},
  before upper={\poppill{Le savais-tu\,?}{white}{poppurple}\ifx\relax#1\relax\else\hspace{6pt}{\popxbold\color{white}#1}\fi\par\vspace{7pt}}}
% Exercice résolu : énoncé en haut, \tcblower puis la solution sur fond crème
\newcounter{popexo}\newcounter{popexr}
\newtcolorbox{exoresolu}[1][]{popbase,colframe=poporange,colbacklower=popcream,
  segmentation style={popink,line width=1.2pt,dashed},
  before upper={\stepcounter{popexr}\popboxhead{Exercice résolu \thepopexr}{poporange}{#1}},
  before lower={\poppill{Solution}{popink}{popcream}\par\vspace{6pt}}}
% Exercice (non résolu en place) — la correction est dans la partie Corrigés
\newtcolorbox{exercice}[1][]{popbase,colframe=popink,
  before upper={\stepcounter{popexo}\popboxhead{Exercice \thepopexo}{popink}{#1}}}
% Corrigé
\newtcolorbox{corrige}[1][]{popbase,colframe=popgreen,colback=popgreenL,
  before upper={\popboxhead{Corrigé}{popgreen}{#1}}}
% Objectifs / prérequis / bilan
\newtcolorbox{objectifs}{popbase,colframe=popink,colback=poporangeL,
  before upper={\popboxhead{Objectifs de la leçon}{popink}{}}}
\newtcolorbox{prerequis}{popbase,colframe=popink,colback=popblueL,
  before upper={\popboxhead{Prérequis}{popblue}{}}}
\newtcolorbox{bilan}{popbase,colframe=popink,colback=white,boxrule=2.8pt,
  before upper={\popboxhead{Fiche bilan}{poporange}{L'essentiel à connaître pour l'examen}}}
% Figure encadrée avec légende
\newtcolorbox{popfigure}[1][]{enhanced,breakable=false,colback=white,colframe=popline,boxrule=1.4pt,arc=8pt,
  left=10pt,right=10pt,top=10pt,bottom=8pt,before skip=10pt,after skip=12pt,halign=center,
  lower separated=false,halign lower=center,
  fontlower=\popsemi\fontsize{9}{11.5}\selectfont\color{popmuted},
  #1}
\newcommand{\legende}[1]{\par\vspace{6pt}{\popsemi\fontsize{9}{11.5}\selectfont\color{popmuted}#1}}

% ============================================================
% Bloc de code source — \begin{lstlisting}[language=Php]...\end{lstlisting}
% (ou language=C, HTML, JavaScript, SQL ; option omise = pseudo-code brut).
% Même recette visuelle que les figures (\popfigure) : cadre encre épais,
% fond clair (popcream), légende possible juste après via \legende{...}
% comme pour une figure (listings ne sait pas arrondir ses coins : on garde
% un cadre rectangulaire, cohérent avec les autres tableaux du document).
% CONTRAT : les modèles ne doivent utiliser QUE \begin{lstlisting}[...] tel
% quel (pas de \begin{tcblisting}, pas de \verb pour un bloc multi-lignes).
% ============================================================
\lstdefinestyle{poplst}{
  backgroundcolor=\color{popcream},
  basicstyle=\popsemi\fontsize{8.6}{11}\selectfont\color{popink},
  keywordstyle=\bfseries\color{popblue},
  commentstyle=\itshape\color{popmuted},
  stringstyle=\color{popgreen},
  numberstyle=\tiny\color{popmuted},
  identifierstyle=\color{popink},
  frame=l,
  rulecolor=\color{popink},
  framerule=1.6pt,
  framesep=7pt,
  framexleftmargin=7pt,
  framexrightmargin=7pt,
  xleftmargin=2pt,
  xrightmargin=2pt,
  breaklines=true,
  breakatwhitespace=false,
  showstringspaces=false,
  showspaces=false,
  showtabs=false,
  tabsize=2,
  columns=flexible,
  keepspaces=true,
  numbers=none,
  aboveskip=10pt,
  belowskip=10pt,
  captionpos=b,
}
\lstset{style=poplst, language=PHP}
% La distribution TeX embarquée ne fournit pas le fichier de langue
% JavaScript de listings (lstlang*.dat) : on la redéfinit nous-mêmes
% pour que \begin{lstlisting}[language=JavaScript] fonctionne.
\lstdefinelanguage{JavaScript}{
  keywords={break,case,catch,class,const,continue,debugger,default,delete,do,
    else,export,extends,finally,for,function,if,import,in,instanceof,let,new,
    return,static,super,switch,this,throw,try,typeof,var,void,while,with,yield,
    async,await,of,get,set},
  keywordstyle=\bfseries\color{popblue},
  ndkeywords={true,false,null,undefined,NaN,Infinity},
  ndkeywordstyle=\bfseries\color{poporange},
  identifierstyle=\color{popink},
  sensitive=true,
  comment=[l]{//},
  morecomment=[s]{/*}{*/},
  string=[b]",
  morestring=[b]',
  morestring=[b]`,
}
% Même limitation pour CSS et les fichiers de configuration (apacheconf,
% nginx, ini...) : ils ne sont pas fournis. CSS et un style générique de
% type "config" (commentaires # ou ;, pas de mots-clés) couvrent les besoins.
\lstdefinelanguage{CSS}{
  keywords={color,background,background-color,margin,padding,border,
    border-radius,font,font-size,font-weight,font-family,width,height,
    display,position,top,left,right,bottom,float,clear,text-align,
    line-height,list-style,cursor,overflow,box-shadow,transition,
    flex,grid,align-items,justify-content,z-index},
  keywordstyle=\bfseries\color{popblue},
  identifierstyle=\color{popink},
  sensitive=false,
  comment=[s]{/*}{*/},
  string=[b]",
  morestring=[b]',
}
\lstdefinelanguage{apacheconf}{
  sensitive=false,
  morecomment=[l]{\#},
}
\lstdefinelanguage{Apache}{
  sensitive=false,
  morecomment=[l]{\#},
}
\lstdefinelanguage{text}{sensitive=false}
\lstdefinelanguage{powershell}{
  keywords={New-Object,New-SmbShare,Get-Acl,Set-Acl,Get-ChildItem,Set-Content,
    Get-Content,Write-Host,ForEach-Object,Where-Object,Select-Object,
    Test-Path,Remove-Item,Copy-Item,Move-Item,Start-Process,Stop-Process,
    If,Else,ElseIf,ForEach,While,Function,Return,Try,Catch},
  keywordstyle=\bfseries\color{popblue},
  identifierstyle=\color{popink},
  sensitive=false,
  comment=[l]{\#},
  string=[b]",
  morestring=[b]',
}
\lstdefinelanguage{ini}{
  sensitive=false,
  morecomment=[l]{;},
  morecomment=[l]{\#},
}

% ============================================================
% Signature OnBuch+
% ============================================================
\newcommand{\poplogoinline}[1]{{\poplogo\fontsize{#1}{#1}\selectfont\color{popink}OnBuch\color{poporange}+}}
\newcommand{\signatureonbuch}{%
  \begin{tcolorbox}[enhanced,colback=popink,colframe=popink,boxrule=2.2pt,arc=10pt,
      drop shadow=poporange,left=14pt,right=14pt,top=10pt,bottom=10pt,before skip=0pt,after skip=0pt]
    \tikz[baseline=-0.7ex]\node[circle,fill=popgreen,draw=popcream,line width=1.3pt,minimum size=22pt,inner sep=0]{\popcheck{white}};%
    \hspace{9pt}\begin{minipage}[c]{0.62\linewidth}
      {\popxbold\fontsize{12}{13}\selectfont\color{popcream}Signé\ \textbullet\ {\poplogo OnBuch\color{poporange}+}}\par\vspace{2pt}
      {\raggedright\popsemi\fontsize{8}{10.5}\selectfont\color{popline}Ludovic A.\\\MakeUppercase{Conforme au programme officiel MINESEC}\par}
    \end{minipage}\hfill
    {\poplogo\fontsize{22}{22}\selectfont\color{popcream}OnBuch\color{poporange}+}
  \end{tcolorbox}%
}

% ============================================================
% Page de garde
% #1 titre · #2 matière · #3 niveau · #4 module · #5 n° leçon · #6 durée ·
% #7 description / objectifs (texte ou itemize)
% ============================================================
\newcommand{\pagegarde}[7]{%
  \thispagestyle{empty}
  \newgeometry{a4paper,margin=1.7cm,top=1.6cm,bottom=1.4cm}%
  \begin{tikzpicture}[remember picture,overlay]
    \fill[poporange!18] ([xshift=-3.2cm,yshift=-3.6cm]current page.north east) circle (4.6cm);
    \fill[poppurple!14] ([xshift=2.4cm,yshift=7.6cm]current page.south west) circle (3.8cm);
  \end{tikzpicture}%
  {\poplogo\fontsize{30}{30}\selectfont\color{popink}OnBuch\color{poporange}+}\hfill
  \raisebox{0.6ex}{\poppill{Cours signé \textbullet\ Premium}{popink}{popcream}}\par
  \vspace{1pt}\popsquiggle{poppurple}\par
  \vspace{8pt}
  \begin{tcolorbox}[enhanced,colback=poporange,colframe=popink,boxrule=2.8pt,arc=13pt,
      drop shadow=popink,left=18pt,right=18pt,top=12pt,bottom=13pt,after skip=10pt]
    \poppill{#2 \textbullet\ #3}{popink}{popcream}\hspace{5pt}\poppill{#4}{white}{popink}\par\vspace{8pt}
    {\popdisplay\fontsize{14}{14}\selectfont\color{popink}LEÇON #5}\par\vspace{4pt}
    {\raggedright\hyphenpenalty=10000\exhyphenpenalty=10000\popdisplay\fontsize{25}{27}\selectfont\color{popcream}\MakeUppercase{#1}\par}
    \vspace{8pt}
    {\popxbold\fontsize{9.5}{9.5}\selectfont\color{popink}\MakeUppercase{Durée conseillée : #6}}
  \end{tcolorbox}
  \begin{tcolorbox}[enhanced,colback=white,colframe=popink,boxrule=2.2pt,arc=10pt,
      drop shadow=popink,left=16pt,right=16pt,top=10pt,bottom=10pt,after skip=8pt]
    \poppill{Dans cette leçon}{poppurple}{white}\par\vspace{5pt}
    \popsemi\fontsize{9.3}{12.6}\selectfont\color{popink}#7
  \end{tcolorbox}
  \noindent\begin{minipage}[t]{0.487\linewidth}
    \begin{tcolorbox}[enhanced,colback=popgreen,colframe=popink,boxrule=2.2pt,arc=10pt,
        drop shadow=popink,left=11pt,right=11pt,top=10pt,bottom=10pt,height=3.1cm,valign=center]
      \tcbox[colback=white,colframe=popink,boxrule=1.3pt,arc=5pt,boxsep=0pt,nobeforeafter,
        left=3pt,right=3pt,top=3pt,bottom=3pt,valign=center]{\qrcode[height=1.9cm]{https://play.google.com/store/apps/details?id=com.luvvix.onbuch&hl=fr}}%
      \hspace{8pt}\begin{minipage}[c]{0.5\linewidth}
        {\popdisplay\fontsize{11}{12.5}\selectfont\color{white}\MakeUppercase{Télécharge OnBuch}}\par\vspace{4pt}
        {\raggedright\popsemi\fontsize{8.4}{11}\selectfont\color{white}Gratuit \textbullet\ Google Play\\Tuteur IA, annales, concours, résultats\par}
      \end{minipage}
    \end{tcolorbox}
  \end{minipage}\hfill
  \begin{minipage}[t]{0.487\linewidth}
    \begin{tcolorbox}[enhanced,colback=poppurple,colframe=popink,boxrule=2.2pt,arc=10pt,
        drop shadow=popink,left=11pt,right=11pt,top=10pt,bottom=10pt,height=3.1cm,valign=center]
      \tikz[baseline=-0.7ex]\node[circle,fill=white,draw=popink,line width=1.3pt,minimum size=1.6cm,inner sep=0]{\popdisplay\fontsize{24}{24}\selectfont\color{poppurple}+};%
      \hspace{8pt}\begin{minipage}[c]{0.6\linewidth}
        {\popdisplay\fontsize{11}{12.5}\selectfont\color{white}\MakeUppercase{Passe à OnBuch+}}\par\vspace{4pt}
        {\raggedright\popsemi\fontsize{8.4}{11}\selectfont\color{white}Tous les cours signés, Tuteur IA illimité, fiches PDF — onglet OnBuch+ de l'app\par}
      \end{minipage}
    \end{tcolorbox}
  \end{minipage}
  \vspace{8pt}
  \popcontenu
  \vfill
  \signatureonbuch
  \restoregeometry
  \newpage
}

% Rangée « Ce cours contient » (page de garde)
\newcommand{\poptile}[3]{% #1 couleur #2 gros texte #3 légende
  \begin{tcolorbox}[enhanced,colback=#1,colframe=popink,boxrule=2pt,arc=9pt,shadow={2.5pt}{-2.5pt}{0pt}{popink},
      left=6pt,right=6pt,top=9pt,bottom=9pt,halign=center,valign=center,height=1.75cm,width=\linewidth,nobeforeafter]
    {\popdisplay\fontsize{12.5}{13}\selectfont\color{white}\MakeUppercase{#2}}\par\vspace{3pt}
    {\centering\popsemi\fontsize{7.8}{9.5}\selectfont\color{white}#3\par}
  \end{tcolorbox}}
\newcommand{\popcontenu}{%
  \par\noindent{\popxboldls\fontsize{7.5}{7.5}\selectfont\color{popmuted}CE COURS SIGNÉ CONTIENT}\par\vspace{7pt}
  \noindent\begin{minipage}[t]{0.235\linewidth}\poptile{popblue}{Cours}{complet, illustré, pas à pas}\end{minipage}\hfill
  \begin{minipage}[t]{0.235\linewidth}\poptile{poporange}{Méthodes}{et exercices résolus}\end{minipage}\hfill
  \begin{minipage}[t]{0.235\linewidth}\poptile{poppink}{Exercices}{type examen + corrigés}\end{minipage}\hfill
  \begin{minipage}[t]{0.235\linewidth}\poptile{popgreen}{Fiche bilan}{l'essentiel à retenir}\end{minipage}\par}

% Sommaire
\newtcolorbox{sommaire}{enhanced,colback=white,colframe=popink,boxrule=2.2pt,arc=10pt,
  drop shadow=popink,left=18pt,right=18pt,top=13pt,bottom=13pt,before skip=6pt,after skip=20pt,
  before upper={\poppill{Sommaire}{poppurple}{white}\par\vspace{9pt}\popsemi\fontsize{10.6}{14.5}\selectfont\color{popink}}}

% Signature de fin de cours — poussée tout en bas de la dernière page
\newcommand{\findecours}{%
  \par\vfill\nopagebreak
  \begin{center}\popsquiggle{poporange}\end{center}\vspace{4pt}
  \signatureonbuch
}
\AtBeginDocument{\def\llbracket{\mathopen{[\mkern-3.5mu[}}\def\rrbracket{\mathclose{]\mkern-3.5mu]}}} % intervalles d'entiers (glyphe ⟦ absent de la police)
% Glyphes absents de Fira Math : redéfinis à partir de glyphes disponibles
\AtBeginDocument{%
  \def\setminus{\mathbin{\backslash}}\let\smallsetminus\setminus
  \def\vdots{\mathord{\vbox{\baselineskip3.2pt\lineskiplimit0pt\kern4pt\hbox{.}\hbox{.}\hbox{.}}}}%
  \def\ddots{\mathinner{\mkern1mu\raise6.5pt\hbox{.}\mkern2mu\raise3.6pt\hbox{.}\mkern2mu\raise0.7pt\hbox{.}\mkern1mu}}%
  \def\triangle{\mathord{\scalebox{0.9}{$\symup{\Delta}$}}}%
  \def\nabla{\mathord{\raisebox{\depth}{\scalebox{1}[-1]{$\symup{\Delta}$}}}}%
  \def\checkmark{\popcheck{popdark}}%
  \def\star{\mathbin{\ast}}\def\bigstar{\mathord{\ast}}%
  \def\diamond{\mathbin{\scalebox{0.6}{$\Diamond$}}}%
  \def\bigcup{\mathop{\mathchoice{\raisebox{-0.3em}{\scalebox{1.7}{$\cup$}}}{\scalebox{1.25}{$\cup$}}{\cup}{\cup}}\displaylimits}%
  \def\bigcap{\mathop{\mathchoice{\raisebox{-0.3em}{\scalebox{1.7}{$\cap$}}}{\scalebox{1.25}{$\cap$}}{\cap}{\cap}}\displaylimits}%
  \def\lesseqqgtr{\mathrel{\lessgtr}}\def\bigoplus{\mathop{\oplus}\displaylimits}%
}
\providecommand{\ssi}{si et seulement si} % abréviation fréquente
\AtBeginDocument{\providecommand{\wideparen}{\overparen}} % arc de cercle

%% FIGFIX-BEGIN v5 — correctifs globaux des figures (docs/onbuch-generation/tools/figfix)
\usepackage{adjustbox}\usepackage{etoolbox}\tcbuselibrary{raster}
% toute figure TikZ/pgfplots plus large que la ligne est réduite à la largeur disponible
\BeforeBeginEnvironment{tikzpicture}{\begin{adjustbox}{max width=\linewidth}}
\AfterEndEnvironment{tikzpicture}{\end{adjustbox}}
% légendes pgfplots lisibles par-dessus les courbes
\pgfplotsset{every axis/.append style={legend style={fill=white,fill opacity=0.92,text opacity=1,draw=black!15,inner sep=3pt}}}
% étiquettes d'arêtes (syntaxe edge["..."]) avec fond blanc
\tikzset{every edge quotes/.append style={fill=white,inner sep=1.2pt}}
%% FIGFIX-END
\begin{document}

\pagegarde{Application cumulée de l’éthique et de l’empathie}{Informatique}{3ème}{Informatique – classe de 3ème}{N11}{1 séance (fiche de progression harmonisée)}{Leçon de synthèse et d'application de l'unité « Éthique et empathie » : l'élève apprend à analyser des situations concrètes de la vie numérique et sociale, à adopter une attitude citoyenne vis-à-vis d'autrui et à rédiger une démarche justifiée. Elle prépare directement aux épreuves de rédaction argumentée du BEPC.
\vspace{4pt}
\begin{multicols}{2}\small
\begin{itemize}
\item À la fin de cette leçon, je sais définir l'éthique numérique et l'empathie et les distinguer clairement.
\item À la fin de cette leçon, je sais identifier dans une situation concrète les comportements contraires à l'éthique (cyberharcèlement, plagiat, atteinte à la vie privée).
\item À la fin de cette leçon, je sais adopter une attitude citoyenne vis-à-vis d'autrui en ligne et hors ligne.
\item À la fin de cette leçon, je sais analyser une situation-problème en suivant une démarche en étapes (observer, identifier, décider, agir, justifier).
\item À la fin de cette leçon, je sais rédiger une réponse argumentée et justifiée à une situation d'éthique numérique.
\item À la fin de cette leçon, je sais proposer des solutions concrètes et respectueuses d'autrui face à un conflit numérique.
\end{itemize}\end{multicols}}

\begin{sommaire}
\begin{multicols}{2}
\begin{enumerate}
\item Rappels structurants : éthique, empathie, citoyenneté numérique
\item Identifier les conduites non éthiques dans la vie courante
\item Adopter une attitude citoyenne vis-à-vis d'autrui
\item La démarche d'analyse d'une situation d'éthique (méthode ODARJ)
\item Rédiger et justifier une réponse argumentée
\item Synthèse et auto-évaluation citoyenne
\item Activité d'intégration
\item Méthodes et exercices résolus
\item Exercices
\item Corrigés des exercices
\item Fiche bilan
\end{enumerate}
\end{multicols}
\end{sommaire}

\begin{objectifs}
\begin{itemize}
\item À la fin de cette leçon, je sais définir l'éthique numérique et l'empathie et les distinguer clairement.
\item À la fin de cette leçon, je sais identifier dans une situation concrète les comportements contraires à l'éthique (cyberharcèlement, plagiat, atteinte à la vie privée).
\item À la fin de cette leçon, je sais adopter une attitude citoyenne vis-à-vis d'autrui en ligne et hors ligne.
\item À la fin de cette leçon, je sais analyser une situation-problème en suivant une démarche en étapes (observer, identifier, décider, agir, justifier).
\item À la fin de cette leçon, je sais rédiger une réponse argumentée et justifiée à une situation d'éthique numérique.
\item À la fin de cette leçon, je sais proposer des solutions concrètes et respectueuses d'autrui face à un conflit numérique.
\item À la fin de cette leçon, je sais citer les textes camerounais protégeant les personnes dans le numérique (loi de 2010 sur la cybersécurité).
\item À la fin de cette leçon, je sais évaluer ma propre conduite numérique à l'aide d'une grille de critères citoyens.
\end{itemize}
\end{objectifs}

\begin{prerequis}
\begin{itemize}
\item Notions d'éthique et de morale vues dans les leçons précédentes de l'unité
\item Notion d'empathie et de respect d'autrui (leçons N9 et N10)
\item Utilisation courante des réseaux sociaux et de la messagerie (WhatsApp, Facebook)
\item Notion de citoyenneté vue en Éducation à la citoyenneté (4ème)
\item Règles de base de rédaction d'un paragraphe argumenté
\end{itemize}
\end{prerequis}

\newpage

\coursec{Rappels structurants : éthique, empathie, citoyenneté numérique}

\courssub{Situation d'accroche et problématique}

Imagine la scène : nous sommes au quartier \textbf{Nkolbisson} à Yaoundé, un mardi soir après les cours. \textbf{Amina}, élève de 3ème, tient son smartphone. Une notification WhatsApp vibre. C'est le groupe de la classe « \texttt{3ème A - Les Géniaux} ». Un camarade vient de poster une photo détournée d'un autre élève, \textbf{Paul}, qui marche avec une canne depuis son accident de moto l'an dernier. La photo est accompagnée d'un montage grossier et d'une légende moqueuse : « \emph{Le Flash version handicapé} ». Trois émojis « mort de rire » suivent. Deux minutes plus tard, cinq autres élèves ont transféré le message. Un sixième écrit : « \texttt{Transfère Amina, c'est juste pour rigoler entre nous, y'a pas de mal !} »

Amina hésite. Son pouce flotte au-dessus de l'écran. Elle connaît Paul. Elle sait la douleur que cache son sourire en classe. Elle sent une boule au ventre : \emph{ce n'est pas drôle}. Mais comment le dire sans passer pour la « rabat-joie » du groupe ? Que lui disent exactement l'éthique numérique et l'empathie ? Et surtout, \textbf{comment rédiger une réponse argumentée, ferme et respectueuse, qui fera réfléchir le groupe au lieu de l'enflammer} ?

C'est toute la finalité de cette leçon : mobiliser les notions d'éthique, d'empathie et de citoyenneté numérique pour \textbf{analyser une situation concrète}, \textbf{prendre position} et \textbf{justifier sa démarche} par écrit. Nous commençons par poser les fondations conceptuelles solides.

\courssub{Qu'est-ce que l'éthique ? Du bien et du mal dans la société}

\begin{definition}[Éthique (définition générale)]
L'\cle{éthique} est l'ensemble des \textbf{règles de conduite} (normes, valeurs, principes) qui, dans une société donnée, permettent de \textbf{distinguer le bien du mal}, le juste de l'injuste, l'acceptable de l'inacceptable. Elle guide le jugement moral et l'action humaine.
\end{definition}

\noindent Contrairement à la loi qui s'impose de l'extérieur (sanction étatique), l'éthique naît d'une \textbf{réflexion intérieure} et d'un \textbf{consensus social}. Elle répond à la question : « \emph{Est-ce que je \textbf{dois} faire cela, même si personne ne me voit et que la loi ne me l'interdit pas explicitement ?} »

\begin{exemplebox}[Exemple courant : la place dans le bus]
Au carrefour \textbf{Carrefour Warda} à Douala, un bus SOTRA s'arrête. Une personne âgée monte. Aucun siège n'est libre. La loi ne t'oblige pas à céder ta place (sauf places réservées). Pourtant, l'éthique sociale camerounaise — le respect des aînés — te dicte de te lever. C'est un \textbf{devoir éthique}, non légal.
\end{exemplebox}

\courssub{L'éthique numérique : l'éthique appliquée au monde connecté}

\begin{definition}[Éthique numérique]
L'\cle{éthique numérique} est l'application des règles éthiques à l'\textbf{usage des technologies de l'information et de la communication} (Internet, réseaux sociaux, messagerie instantanée, téléphonie mobile, intelligence artificielle, etc.). Elle définit ce qui est \textbf{correct, respectueux et responsable} dans nos interactions en ligne.
\end{definition}

\noindent Le numérique amplifie les conséquences de nos actes : \textbf{vitesse} (un clic propage à des milliers), \textbf{permanence} (une capture d'écran survit à la suppression), \textbf{portée} (le village planétaire voit tout). L'éthique numérique n'est pas une « éthique à part » : c'est la \textbf{même éthique}, transposée dans un environnement où l'absence de visage physique (écran interposé) facilite la désinhibition et la cruauté.

\begin{propriete}[Champ d'application de l'éthique numérique]
L'éthique numérique couvre notamment :
\begin{itemize}
    \item La \textbf{protection de la vie privée} (ne pas publier photos/infos d'autrui sans consentement).
    \item Le \textbf{respect de la dignité} (pas de moqueries, harcèlement, discours de haine).
    \item L'\textbf{honnêteté intellectuelle} (citer ses sources, ne pas plagier, ne pas propager de fausses informations).
    \item La \textbf{sécurité} (ne pas partager de mots de passe, ne pas pirater, signaler les failles).
    \item La \textbf{responsabilité} (assumer ses publications, corriger ses erreurs).
\end{itemize}
\end{propriete}

\courssub{L'empathie : le moteur invisible de l'éthique}

\begin{definition}[Empathie]
L'\cle{empathie} est la \textbf{capacité à se mettre à la place d'autrui}, à \textbf{comprendre} ses émotions, ses pensées, sa perspective, et à \textbf{les partager} (ressentir \emph{avec} l'autre), sans les confondre avec les siennes. Elle se distingue de la sympathie (ressentir \emph{pour} l'autre, pitié) et de la contagion émotionnelle (subir l'émotion sans la comprendre).
\end{definition}

\noindent L'empathie a deux composantes indissociables :
\begin{enumerate}
    \item \textbf{Cognitive} : comprendre intellectuellement la situation de l'autre (« \emph{Paul a mal à la jambe, il a subi des opérations, il fait des efforts pour marcher} »).
    \item \textbf{Affective} : ressentir une résonance émotionnelle (« \emph{Je sens sa honte, sa colère, sa tristesse quand il verra cette photo} »).
\end{enumerate}

\begin{aretenir}[Lien empathie $\rightarrow$ éthique $\rightarrow$ citoyenneté]
L'empathie est le \textbf{moteur interne} : elle nous \textbf{fait sentir} la blessure d'autrui.
L'éthique est la \textbf{boussole} : elle nous \textbf{dit} que blesser autrui est mal.
La \textbf{citoyenneté numérique} est l'\textbf{acte extérieur} : elle \textbf{traduit} ce ressenti et ce jugement en \textbf{conduites concrètes} (ne pas transférer, signaler, défendre la victime, éduquer le groupe).
\end{aretenir}

\noindent Sans empathie, l'éthique reste une règle abstraite qu'on contourne. Sans éthique, l'empathie reste une émotion stérile. Sans citoyenneté, les deux restent privées et ne changent pas le monde numérique.

\courssub{Les trois piliers de la conduite citoyenne en ligne}

\begin{propriete}[Les trois piliers de la citoyenneté numérique (programme officiel)]
Toute conduite citoyenne en ligne repose sur trois piliers indissociables :
\begin{center}
\begin{tabularx}{\linewidth}{|>{\bfseries}l|X|X|}
\hline
\textbf{Pilier} & \textbf{Sens concret en ligne} & \textbf{Exemple d'action (Amina)} \\
\hline
\cle{Respect} & Reconnaître la dignité, l'intimité, l'opinion, la différence d'autrui. Ne pas outrager, moquer, harceler. & Ne pas transférer la photo. Répondre : « \emph{Cette photo blesse Paul. On ne se moque pas d'un handicap.} » \\
\hline
\cle{Responsabilité} & Assumer ses actes numériques (publications, partages, commentaires). Mesurer les conséquences. Corriger ses erreurs. & Si elle a déjà ri (réaction émoji), elle retire son « like », s'excuse auprès de Paul en privé, et explique pourquoi au groupe. \\
\hline
\cle{Solidarité} & Soutenir les plus vulnérables. Signaler les abus. Protéger la vie privée collective. Ne pas être spectateur passif. & Signaler le message aux admins du groupe / à la modération WhatsApp. Proposer au groupe une charte « \emph{Zéro moquerie sur les différences} ». \\
\hline
\end{tabularx}
\end{center}
\end{propriete}

\courssub{Distinction cruciale : Légal vs Éthique}

C'est l'erreur n°1 des élèves (et des adultes) : \textbf{confondre « c'est permis par la loi » et « c'est bien »}.

\begin{definition}[Légal vs Éthique]
\begin{itemize}
    \item \textbf{Légal} : conforme à la \textbf{loi} en vigueur (code pénal, loi sur la cybersécurité, droit d'auteur, RGPD local). Sanction : amende, prison, dommages-intérêts (procès).
    \item \textbf{Éthique} : conforme au \textbf{bien moral}, aux valeurs humaines universelles et culturelles (respect, dignité, vérité, justice). Sanction : perte de confiance, réputation ternie, exclusion sociale, remords, \textbf{conscience coupable}.
\end{itemize}
\end{definition}

\begin{attention}[Piège fréquent : « Ce n'est pas illégal, donc je le fais »]
\begin{itemize}
    \item \textbf{Publier une rumeur non vérifiée} sur un camarade (ex : « \emph{Untel a volé le téléphone du prof} ») : peut ne pas être \emph{immédiatement} poursuivi pénalement (si pas de plainte, si diffamation non constituée), mais est \textbf{éthiquement condamnable} (atteinte à la réputation, injustice, manque de vérité).
    \item \textbf{Partager une photo intime} reçue en confidence : même si la loi ne l'interdit pas explicitement dans tous les cas (selon les juridictions), c'est une \textbf{trahison éthique majeure} (violation de la confiance, de l'intimité, de la dignité).
    \item \textbf{Utiliser une faille de sécurité} d'un site scolaire pour voir les notes des autres : techniquement possible, parfois zone grise légale, mais \textbf{totalement anti-éthique} (honnêteté, respect du travail d'autrui).
\end{itemize}
\emph{Règle d'or :} \textbf{La loi pose le plancher minimum ; l'éthique vise le plafond de l'humain.}
\end{attention}

\begin{savaistu}[Le cadre légal camerounais : Loi n°2010/012]
Le Cameroun dispose d'un arsenal juridique fort. La \textbf{loi n°2010/012 du 21 décembre 2010 relative à la cybersécurité et à la cybercriminalité} punit notamment :
\begin{itemize}
    \item L'\textbf{atteinte à la vie privée} (Art. 7-1) : interception, conservation, divulgation de correspondances privées sans consentement.
    \item La \textbf{diffamation et l'injure} par voie électronique (Art. 8) : allégation ou imputation d'un fait portant atteinte à l'honneur.
    \item Le \textbf{harcèlement sexuel ou moral} en ligne (Art. 9).
    \item L'\textbf{usurpation d'identité} numérique (Art. 10).
\end{itemize}
Les peines vont jusqu'à \textbf{5 ans de prison et 10 millions de FCFA d'amende}. La loi rattrape souvent ce que l'éthique a laissé faire. Mieux vaut \textbf{agir éthiquement \emph{avant}} que la loi n'intervienne.
\end{savaistu}

\courssub{Illustration : Empathie, Éthique, Citoyenneté — une architecture en trois cercles}

\begin{popfigure}
\begin{tikzpicture}[
    every node/.style={font=\sffamily\small},
    circle1/.style={draw=popblue, thick, fill=popblueL, minimum width=#1},
    circle2/.style={draw=poporange, thick, fill=poporangeL, minimum width=#1},
    circle3/.style={draw=popgreen, thick, fill=popgreenL, minimum width=#1},
    arrow/.style={->, >=Stealth, thick, popdark},
    labelstyle/.style={align=center, text width=2.8cm}
]
% Cercles concentriques (du plus grand au plus petit pour le dessin)
% Citoyenneté numérique (extérieur)
\node[circle3=10cm] (cit) at (0,0) {};
\node[labelstyle, text=popgreen!80!black, align=center] at (0, 3.8){\textbf{CITOYENNETÉ NUMÉRIQUE}\\\emph{Actes concrets : Respect, Responsabilité, Solidarité}};

% Éthique (milieu)
\node[circle2=7cm] (eth) at (0,0) {};
\node[labelstyle, text=poporange!80!black, align=center] at (0, 1.8){\textbf{ÉTHIQUE NUMÉRIQUE}\\\emph{Règles de conduite : Bien vs Mal, Dignité, Vie privée}};

% Empathie (centre)
\node[circle1=4cm] (emp) at (0,0) {};
\node[labelstyle, text=popblue!80!black, align=center] at (0, 0){\textbf{EMPATHIE}\\\emph{Moteur interne : Comprendre \& Ressentir l'autre}};

% Flèches d'influence (centrifuges)
\draw[arrow] (emp.north) -- ++(0,0.5) node[above, labelstyle, text=popdark] {inspire \& guide};
\draw[arrow] (eth.north) -- ++(0,0.5) node[above, labelstyle, text=popdark] {se traduit en};
\end{tikzpicture}
\legende{Architecture de la conduite citoyenne : l'empathie (cœur) motive l'éthique (boussole) qui s'incarne en citoyenneté numérique (actions).}
\end{popfigure}

\courssub{Tableau comparatif : Légal vs Éthique — Deux exemples camerounais}

\begin{popfigure}
\begin{center}
\begin{tabularx}{\linewidth}{|>{\bfseries}l|X|X|}
\hline
\rowcolor{popblueL}
\textbf{Situation concrète (Cameroun)} & \textbf{Analyse LÉGALE (Loi n°2010/012, Code pénal)} & \textbf{Analyse ÉTHIQUE (Valeurs humaines)} \\
\hline
\textbf{1. Amina transfère la photo moquant le handicap de Paul dans le groupe WhatsApp de la classe (50 élèves).}
& \textbf{Zone grise / Risque légal :} Pourrait relever de l'atteinte à la vie privée (Art. 7-1) ou harcèlement moral (Art. 9) si répétition. Poursuite possible si Paul/parents portent plainte. Souvent non poursuivi « entre élèves ». & \textbf{Clairément NON ÉTHIQUE :} Violation de la \textbf{dignité} (handicap), de l'\textbf{intimité} (image), de la \textbf{bienveillance}. Trahison de la \textbf{confiance} (groupe classe). Banalisation du \textbf{handicapisme}. Amina devient \textbf{complice active} de la souffrance de Paul. \\
\hline
\textbf{2. Un élève crée un faux compte Facebook au nom du proviseur et poste : « Demain, cours annulés, fête au lycée ! »}
& \textbf{ILLÉGAL :} \textbf{Usurpation d'identité numérique} (Art. 10 Loi 2010/012) + \textbf{Trouble à l'ordre public} (fausse info). Peine : prison + amende. L'administration \textbf{portera plainte}. & \textbf{Clairément NON ÉTHIQUE :} \textbf{Mensonge}, \textbf{manipulation}, \textbf{irresponsabilité} (élèves qui viennent pour rien, danger routier), \textbf{manque de respect} à l'autorité et à la communauté éducative. \\
\hline
\end{tabularx}
\end{center}
\legende{Comparatif Légal vs Éthique : un acte peut être « pas encore illégal » ou « zone grise légale » mais profondément anti-éthique. L'éthique exige plus que la loi.}
\end{popfigure}

\courssub{Exemple résolu : Analyser la situation d'Amina avec la grille éthique}

\begin{exoresolu}[Analyse éthique de la situation d'accroche]
\textbf{Énoncé :} Reprenons la situation d'Amina. En utilisant les trois piliers (Respect, Responsabilité, Solidarité) et la distinction Légal/Éthique, analyse la conduite de : (a) l'auteur de la photo, (b) les 5 transféreurs, (c) Amina si elle transfère, (d) Amina si elle refuse et s'exprime.

\tcblower
\textbf{Solution détaillée :}

\begin{center}
\begin{tabularx}{\linewidth}{|l|X|X|X|}
\hline
\rowcolor{popblueL}
\textbf{Acteur} & \textbf{Respect} & \textbf{Responsabilité} & \textbf{Solidarité} \\
\hline
\textbf{Auteur de la photo} & \textcolor{popred}{Violé} : Atteinte grave à la dignité de Paul (handicap), utilisation non consentie de son image. & \textcolor{popred}{Absente} : Ne assume pas (anonyme ou pseudo), ne mesure pas l'impact viral. & \textcolor{popred}{Absente} : Agresse le plus vulnérable au lieu de le protéger. \\
\hline
\textbf{5 Transféreurs} & \textcolor{popred}{Violé} : Participent à l'humiliation publique. « Juste pour rire » nie la souffrance. & \textcolor{poporange}{Faible} : Se cachent derrière l'effet de groupe (« tout le monde le fait »). & \textcolor{popred}{Absente} : Deviennent \textbf{harceleurs par relais}. Amplifient la portée. \\
\hline
\textbf{Amina (si transfère)} & \textcolor{popred}{Violé} : Trahit sa propre conscience (elle « sent » que c'est mal). & \textcolor{popred}{Absente} : Cède à la pression sociale, perd son autonomie morale. & \textcolor{popred}{Absente} : Abandonne Paul. Devient \textbf{bourreau par omission/action}. \\
\hline
\textbf{Amina (si refuse \& s'exprime)} & \textcolor{popgreen}{Honoré} : Reconnaît la dignité de Paul. Dit « non » à l'indigne. & \textcolor{popgreen}{Assumée} : Prend le risque social (moqueries « rabat-joie ») pour rester cohérente. & \textcolor{popgreen}{Vécue} : Protège Paul (signalement, soutien privé), éduque le groupe (argumentation), fait avancer la norme du groupe. \\
\hline
\end{tabularx}
\end{center}

\textbf{Conclusion de l'analyse :} Seule la posture (d) est \textbf{pleinement citoyenne}. Elle demande du \textbf{courage moral} (vertu cardinale), fruit de l'empathie (je sens la douleur de Paul) et de l'éthique (je sais que c'est mal). La réponse d'Amina doit \textbf{nommer l'émotion} (« \emph{Paul a mal} »), \textbf{nommer le principe} (« \emph{On ne se moque pas du handicap} »), \textbf{poser l'acte} (« \emph{Je ne transfère pas, je signale} »).
\end{exoresolu}

\courssub{Mini-TP guidé : « Je m'entraîne à voir l'invisible »}

\begin{experience}[Exercice d'empathie cognitive et affective — 10 min]
\textbf{Objectif :} Développer la composante cognitive (comprendre) et affective (ressentir) de l'empathie face à une situation de cyberharcèlement banalisé.

\textbf{Matériel :} Une feuille, un stylo, ton honnêteté.

\textbf{Consignes :}
\begin{enumerate}
    \item \textbf{Lis} ce message fictuel reçu sur WhatsApp (groupe « Amis du Lycée ») : \\
    \texttt{« Regardez cette vidéo de Kevin (3ème B) qui tombe en EPS. Il court comme une vache boiteuse lol. Partagez max ! »} \\
    (La vidéo montre Kevin, essoufflé, trébuchant, le short déchiré, des rires en fond sonore.)
    \item \textbf{Étape 1 — Cognitive (2 min) :} Écris 3 phrases commençant par « \emph{Kevin, en ce moment, il doit...} » (ex : « \emph{...se sentir humilié devant tout le lycée} », « \emph{...avoir peur de venir en cours demain} », « \emph{...se demander qui a filmé au lieu d'aider} » »).
    \item \textbf{Étape 2 — Affective (2 min) :} Ferme les yeux 30 secondes. Visualise Kevin. Note \textbf{un mot} l'émotion qui monte en \textbf{toi} (pas ce que Kevin ressent, ce que \textbf{toi} tu ressens : honte ? colère ? tristesse ? malaise ?).
    \item \textbf{Étape 3 — Décision éthique (1 min) :} Écris \textbf{la réponse exacte} que tu envoies dans le groupe. Contrainte : \textbf{pas d'insulte, pas de sermon, 2 phrases max}.
    \item \textbf{Étape 4 — Partage (5 min) :} En binôme, lis ta réponse à ton voisin. Discutez : laquelle est la plus \textbf{efficace pour arrêter le partage} sans braquer ?
\end{enumerate}

\textbf{Observations attendues :}
\begin{itemize}
    \item Les réponses « \texttt{Arrêtez, c'est nul} » sont éthiques mais peu efficaces (braquent).
    \item Les réponses empathiques + fermes marchent mieux : \texttt{« Kevin est mon pote. Cette vidéo le détruit. Je la signale. Faites pareil. »} ou \texttt{« Imaginez votre petite sœur à sa place. Supprimez-la. Maintenant. »}
\end{itemize}

\textbf{Interprétation :} L'empathie \textbf{guide le vocabulaire} (mots justes, pas de violence verbale) et \textbf{renforce la légitimité} (tu parles au nom de la victime, pas au tien). C'est la base de la \textbf{rhétorique citoyenne} qu'on approfondira en Section 5.
\end{experience}

\courssub{Synthèse de la section : ce qu'il faut retenir avant d'aller plus loin}

\begin{aretenir}[Checklist des acquis — Section 1]
\begin{itemize}
    \item \textbf{Éthique} = règles collectives bien/mal. \textbf{Éthique numérique} = même chose en ligne (amplifié par vitesse/portée/permanence).
    \item \textbf{Empathie} = comprendre (cognitif) + ressentir (affectif) la situation d'autrui. C'est le \textbf{moteur} de l'éthique.
    \item \textbf{Citoyenneté numérique} = \textbf{traduction en actes} (Respect, Responsabilité, Solidarité).
    \item \textbf{Légal $\neq$ Éthique}. La loi = plancher (sanction étatique). L'éthique = plafond (exigence humaine). \textbf{Ne jamais dire « c'est pas interdit » pour justifier un acte blessant.}
    \item \textbf{Amina} : son dilemme est le nôtre. La leçon entière vise à lui donner (et à te donner) les \textbf{outils d'analyse} (grille 3 piliers, distinction légal/éthique) et les \textbf{outils d'expression} (rédaction justifiée, méthode ODARJ en Section 4) pour transformer l'hésitation en \textbf{action citoyenne efficace}.
\end{itemize}
\end{aretenir}

\noindent Nous avons posé les fondations. La section suivante nous apprendra à \textbf{repérer concrètement} les conduites non éthiques dans la vie courante (cyberharcèlement, fake news, vol de données, etc.) pour ne plus jamais dire « \emph{je ne savais pas} ».

\coursec{Identifier les conduites non éthiques dans la vie courante}

Dans la section précédente, nous avons posé les fondations : l'\cle{éthique} comme ensemble de règles de conduite, l'\cle{empathie} comme capacité à se mettre à la place d'autrui, et la \cle{citoyenneté numérique} comme comportement responsable dans le monde connecté. Mais reconnaître les valeurs ne suffit pas : encore faut-il savoir repérer, dans la vie de tous les jours, les situations où ces valeurs sont bafouées. C'est précisément l'objectif de cette section : t'apprendre à \textbf{identifier} les conduites non éthiques, à les \textbf{nommer} correctement et à les \textbf{qualifier} à l'aide de critères précis. Car on ne peut combattre que ce que l'on sait reconnaître.

\courssub{Pourquoi apprendre à identifier les conduites non éthiques ?}

Imagine la scène suivante, très fréquente dans nos collèges : pendant la récréation, un élève montre à ses amis une photo modifiée d'un camarade, tout le monde rit, et la photo est transférée dans le groupe WhatsApp de la classe. Le soir même, l'élève photographié ne veut plus venir au collège. Qui a mal agi ? Celui qui a créé la photo ? Ceux qui ont ri ? Ceux qui l'ont transférée ? Ceux qui sont restés silencieux ?

Sans une méthode claire, il est difficile de répondre. Or, au BEPC comme dans la vie, tu devras \textbf{analyser des situations concrètes} et \textbf{justifier ta réponse}. Identifier une conduite non éthique, c'est donc :

\begin{itemize}
\item \textbf{protéger les autres} : repérer tôt une situation de souffrance permet d'intervenir avant qu'elle ne s'aggrave ;
\item \textbf{se protéger soi-même} : beaucoup d'élèves participent à des torts sans s'en rendre compte (un simple transfert peut te rendre complice) ;
\item \textbf{devenir un citoyen numérique lucide} : capable de distinguer une plaisanterie acceptable d'une atteinte à la dignité.
\end{itemize}

\courssub{Les six grandes conduites non éthiques à connaître}

Passons en revue, une par une, les conduites que tu dois savoir reconnaître. Pour chacune, nous partons d'une situation de la vie courante, puis nous donnons la définition rigoureuse.

\textbf{1. Le cyberharcèlement.} Marlyse, élève en classe de 3ème, reçoit chaque soir des SMS insultants sur son apparence physique. Les mêmes moqueries réapparaissent sous ses photos sur les réseaux sociaux. Ce n'est pas une dispute ponctuelle : c'est \emph{répété}, \emph{voulu} et \emph{blessant}.

\begin{definition}[Cyberharcèlement]
Le \cle{cyberharcèlement} est l'ensemble des moqueries, insultes, menaces ou humiliations \textbf{répétées} à l'encontre d'une personne, par voie numérique : réseaux sociaux, SMS, messageries instantanées, groupes de discussion. Il se distingue de la simple dispute par trois éléments : l'\textbf{intention de nuire}, la \textbf{répétition} dans le temps et le \textbf{déséquilibre} entre l'agresseur (souvent entouré ou anonyme) et la victime (souvent isolée).
\end{definition}

\textbf{2. La violation de la vie privée.} Junior photographie son voisin en train de dormir en classe et publie la photo avec un commentaire moqueur, sans lui demander son avis. Même si la photo est « vraie », sa publication est une faute.

\begin{definition}[Violation de la vie privée]
La \cle{violation de la vie privée} consiste à publier, diffuser ou exploiter la photo, la voix, les informations personnelles (adresse, notes, situation familiale, état de santé) d'une personne \textbf{sans son consentement}. Chaque individu a un droit sur son image et sur ses données : c'est le \textbf{droit à l'image} et le droit à la protection de la vie privée.
\end{definition}

\textbf{3. Le plagiat et la triche numérique.} Pour son exposé d'histoire, Sandrine copie intégralement un texte trouvé sur Internet et le présente comme son propre travail. Elle obtient une bonne note, mais cette note est volée.

\begin{definition}[Plagiat]
Le \cle{plagiat} consiste à reproduire totalement ou partiellement le travail d'autrui (texte, image, devoir, idée) \textbf{en le présentant comme le sien}, sans citer la source. Sur Internet, la facilité du « copier-coller » rend le plagiat tentant, mais il reste une \textbf{triche} : il trompe l'enseignant et vole le travail de l'auteur. Utiliser une information trouvée en ligne est permis à condition de \textbf{citer sa source} et de la reformuler.
\end{definition}

\textbf{4. La désinformation.} Une rumeur circule : « les examens sont annulés cette année ». Sans vérifier auprès de l'école ou du MINESEC, des centaines d'élèves la partagent. Le désordre s'installe.

\begin{definition}[Fake news (désinformation)]
Une \cle{fake news} (fausse information) est une information \textbf{fausse ou trompeuse}, créée ou relayée volontairement ou par négligence, et présentée comme vraie. La \cle{désinformation} est la conduite qui consiste à créer ou à propager de telles informations \textbf{sans les avoir vérifiées} auprès de sources fiables (site officiel, administration, enseignant, média reconnu).
\end{definition}

\textbf{5. L'exclusion numérique.} Le groupe WhatsApp de la classe de 3ème B sert à partager les devoirs et les leçons. Mais les administrateurs refusent d'y ajouter Étienne, ou le retirent volontairement, parce qu'il est « différent ». Étienne rate les informations importantes.

\begin{definition}[Exclusion numérique]
L'\cle{exclusion numérique} est le fait d'\textbf{écarter volontairement} une personne d'un espace numérique collectif (groupe de classe, forum, liste de diffusion) auquel elle a légitimement sa place. Elle blesse doublement : elle prive la victime d'informations utiles et lui signifie qu'elle n'est pas des nôtres. C'est une atteinte directe à la dignité et à la solidarité.
\end{definition}

\textbf{6. L'usurpation d'identité.} Rodrigue connaît le mot de passe de son camarade. Pour « rigoler », il se connecte à son compte et envoie des messages embarrassants en son nom. Les destinataires croient que c'est le camarade qui parle.

\begin{definition}[Usurpation d'identité]
L'\cle{usurpation d'identité} consiste à \textbf{se faire passer pour autrui} en utilisant son compte, son nom, sa photo ou ses identifiants, afin d'agir en son nom (envoyer des messages, publier, nuire à sa réputation). C'est une faute grave : elle trompe les tiers, salit la victime et détruit la confiance, qui est le fondement de toute communication numérique.
\end{definition}

\courssub{Les quatre critères de reconnaissance d'une conduite non éthique}

Face à une situation ambiguë (« est-ce une blague ou du harcèlement ? »), il faut un outil d'analyse. Retiens ces \textbf{quatre critères} :

\begin{methode}[Qualifier une conduite : les 4 critères]
Pour déterminer si une conduite numérique est non éthique, vérifie successivement :
\begin{enumerate}
\item \textbf{L'intention} : l'auteur cherche-t-il à nuire, à humilier, à tromper ou à s'enrichir indûment ? Une intention de nuire est l'indice le plus fort.
\item \textbf{La répétition} : le fait se reproduit-il dans le temps ? Un acte unique peut être une maladresse ; des actes répétés constituent un harcèlement.
\item \textbf{Le consentement} : la personne concernée a-t-elle donné son accord (pour sa photo, ses informations, l'usage de son compte) ? L'absence de consentement rend la conduite fautive, même si le contenu est « vrai » ou « drôle ».
\item \textbf{La dignité} : la conduite rabaisse-t-elle, humilie-t-elle ou exclut-elle la personne ? Toute atteinte à la dignité d'autrui est non éthique, même sans intention avouée.
\end{enumerate}
\textbf{Règle de décision} : si l'un de ces critères est clairement violé, la conduite est non éthique. Plusieurs critères violés ensemble signalent une faute grave.
\end{methode}

\begin{attention}[« Je n'ai fait que transférer ! »]
Erreur très fréquente : croire que \textbf{relayer} un message moqueur, une photo volée ou une fausse rumeur « n'engage pas », parce qu'on n'en est pas l'auteur. C'est faux. \textbf{Le relais est aussi une participation au tort} : chaque transfert amplifie l'humiliation, élargit l'audience et prolonge la souffrance de la victime. Devant un contenu blessant, la conduite citoyenne est de \textbf{ne pas transférer, ne pas commenter, ne pas « liker »}, et si possible d'alerter un adulte ou de soutenir la victime. Le silence complice et le partage amusé font partie du problème.
\end{attention}

\courssub{L'arbre de décision : « Ce message est-il éthique ? »}

Pour t'aider à décider \emph{avant} de publier ou de transférer, voici un raisonnement simple en trois questions. Il fonctionne comme un petit algorithme : chaque réponse « non » t'arrête immédiatement.

\begin{popfigure}
\begin{tikzpicture}[
  font=\small,
  decision/.style={diamond, draw=popblue, fill=popblueL, text width=2.6cm, align=center, inner sep=1pt, aspect=2},
  action/.style={rectangle, rounded corners, draw=popgreen, fill=popgreenL, text width=2.8cm, align=center, inner sep=4pt},
  stopbox/.style={rectangle, rounded corners, draw=poporange, fill=poporangeL, text width=3.2cm, align=center, inner sep=4pt},
  startbox/.style={rectangle, rounded corners, draw=popdark, fill=popcream, text width=3.4cm, align=center, inner sep=4pt},
  fleche/.style={-{Stealth[length=2.5mm]}, thick, popdark}
]
\node[startbox] (start) {Je veux publier ou transférer un message};
\node[decision, below=0.9cm of start] (q1) {Ai-je le consentement des personnes concernées ?};
\node[decision, below=1.1cm of q1] (q2) {Le message respecte-t-il la dignité d'autrui ?};
\node[decision, below=1.1cm of q2] (q3) {L'information est-elle vraie et vérifiée ?};
\node[action, below=1.1cm of q3] (ok) {Je peux publier en toute éthique};
\node[stopbox, right=2.2cm of q2] (stop) {Je ne publie pas, je ne transfère pas};

\draw[fleche] (start) -- (q1);
\draw[fleche] (q1) -- node[left]{Oui} (q2);
\draw[fleche] (q2) -- node[left]{Oui} (q3);
\draw[fleche] (q3) -- node[left]{Oui} (ok);
\draw[fleche] (q1.east) -- node[above, pos=0.25]{Non} (stop.north west);
\draw[fleche] (q2.east) -- node[above]{Non} (stop.west);
\draw[fleche] (q3.east) -- node[above, pos=0.3]{Non} (stop.south west);
\end{tikzpicture}
\legende{Arbre de décision « Ce message est-il éthique ? » : trois questions (consentement, respect, vérité). Un seul « Non » suffit à bloquer la publication.}
\end{popfigure}

Ce schéma résume la logique de la méthode : le \textbf{consentement} (critère 3), le \textbf{respect de la dignité} (critère 4) et la \textbf{véracité} (qui protège contre la désinformation) sont trois barrières successives. Un message qui franchit les trois peut être publié ; un seul échec doit te faire renoncer.

\courssub{Une situation analysée pas à pas : la photo moquante}

Voyons comment ces critères s'appliquent à une scène concrète, racontée en trois temps.

\begin{popfigure}
\begin{tikzpicture}[
  font=\small,
  case/.style={rectangle, draw=popdark, fill=white, text width=3.4cm, minimum height=3.2cm, align=center, inner sep=5pt},
  titre/.style={rectangle, rounded corners, draw=popblue, fill=popblue, text=white, font=\small\bfseries, inner sep=3pt},
  fleche/.style={-{Stealth[length=3mm]}, very thick, poporange}
]
\node[titre] at (0,2.1) {1. Réception};
\node[case] (c1) at (0,0) {Amina reçoit une photo moquante d'un camarade, envoyée dans le groupe de la classe.};
\node[titre] at (5,2.1) {2. Le choix};
\node[case] (c2) at (5,0) {Elle hésite : rire et transférer, ou refuser de relayer et soutenir le camarade ?};
\node[titre] at (10,2.1) {3. Les conséquences};
\node[case] (c3) at (10,0) {Transférer : l'humiliation grandit, elle devient complice. Refuser et alerter : la chaîne s'arrête.};
\draw[fleche] (c1.east) -- (c2.west);
\draw[fleche] (c2.east) -- (c3.west);
\end{tikzpicture}
\legende{Bande dessinée schématisée en trois cases : réception d'un contenu blessant, moment du choix, conséquences des deux attitudes possibles.}
\end{popfigure}

Analysons avec nos quatre critères. La photo moquante viole la \textbf{dignité} (elle ridiculise), elle est publiée sans \textbf{consentement}, et si les moqueries se répètent, le critère de \textbf{répétition} bascule la situation vers le cyberharcèlement. Le choix d'Amina à la case 2 est donc décisif : transférer la rend participante du tort, refuser fait d'elle un maillon de la solution.

\begin{exemplebox}[Au collège de Bafoussam : trois conduites en une seule situation]
Au collège de Bafoussam, le groupe WhatsApp de la classe de 3ème C sert à échanger les devoirs. Les administrateurs \textbf{retirent volontairement} un élève, Nkeng, du groupe. Puis ils \textbf{publient ses notes} obtenues au dernier devoir, accompagnées de \textbf{commentaires moqueurs} repris chaque jour pendant une semaine.
\par
Trois conduites non éthiques sont présentes :
\begin{enumerate}
\item \textbf{L'exclusion numérique} : retirer Nkeng du groupe de classe le prive des informations scolaires et le rejette de la communauté ;
\item \textbf{La violation de la vie privée} : diffuser ses notes sans son consentement expose ses informations personnelles ;
\item \textbf{Le cyberharcèlement} : les commentaires moqueurs, répétés pendant une semaine avec intention d'humilier, correspondent exactement à la définition.
\end{enumerate}
\end{exemplebox}

\courssub{Exercice résolu : qualifier des conduites}

\begin{exoresolu}[À toi de juger]
Pour chacune des situations suivantes, identifie la conduite non éthique mise en jeu et justifie ta réponse à l'aide d'au moins un des quatre critères (intention, répétition, consentement, dignité).
\begin{enumerate}
\item Bertrand recopie la rédaction de français trouvée sur un site Internet et la rend en signant de son nom.
\item Chaque soir depuis deux semaines, des SMS insultants sont envoyés à Solange par un numéro inconnu.
\item Vanessa partage dans son statut une « information » annonçant la fermeture de tous les collèges, sans l'avoir vérifiée.
\item Alain utilise le téléphone de son ami, laissé sans surveillance, pour envoyer un message d'amour à une camarade en se faisant passer pour lui.
\end{enumerate}
\tcblower
\textbf{Solution détaillée.}
\begin{enumerate}
\item \textbf{Plagiat (triche numérique).} Bertrand présente le travail d'autrui comme le sien, sans citer la source. Critère violé : l'\textbf{intention} de tromper l'enseignant pour obtenir une note imméritée. La conduite correcte aurait été de reformuler et de citer le site utilisé.
\item \textbf{Cyberharcèlement.} Trois critères sont réunis : l'\textbf{intention} de nuire (insultes), la \textbf{répétition} (chaque soir, depuis deux semaines) et l'atteinte à la \textbf{dignité}. L'anonymat du numéro aggrave la lâcheté de l'agresseur. Solange doit conserver les messages comme preuves et en parler à un adulte (parent, enseignant, censeur).
\item \textbf{Désinformation (relais d'une fake news).} Vanessa propage une fausse information sans vérification auprès d'une source fiable (le collège, le MINESEC). Même sans mauvaise intention, le relais négligent est fautif : avant de partager, il faut \textbf{vérifier} la source.
\item \textbf{Usurpation d'identité.} Alain se fait passer pour son ami en utilisant son téléphone et son compte, \textbf{sans son consentement}. Même présenté comme une plaisanterie, cet acte trompe la destinataire et engage la réputation de la victime. C'est une faute, pas une blague.
\end{enumerate}
\end{exoresolu}

\courssub{Ce qu'il faut retenir}

\begin{aretenir}[Identifier une conduite non éthique]
\begin{itemize}
\item Les six conduites non éthiques à reconnaître : \cle{cyberharcèlement}, \cle{violation de la vie privée}, \cle{plagiat}, \cle{désinformation}, \cle{exclusion numérique}, \cle{usurpation d'identité}.
\item Quatre critères de qualification : \textbf{intention} de nuire, \textbf{répétition}, absence de \textbf{consentement}, atteinte à la \textbf{dignité}. Un seul critère violé suffit à qualifier la conduite de non éthique.
\item Avant de publier ou transférer : trois questions — consentement ? respect ? vérité ? Un seul « non » impose de s'abstenir.
\item Le relais (transfert, commentaire, « like ») est une \textbf{participation au tort}, pas un acte neutre.
\end{itemize}
\end{aretenir}

Maintenant que tu sais \emph{identifier} et \emph{nommer} ces conduites, la prochaine étape consiste à apprendre à \emph{réagir} : c'est l'objet de la section suivante, consacrée à l'adoption d'une attitude citoyenne vis-à-vis d'autrui.

\coursec{Adopter une attitude citoyenne vis-à-vis d'autrui}

Après avoir appris à \emph{identifier} les conduites non éthiques (harcèlement, diffusion de rumeurs, usurpation d'identité, plagiat), nous devons maintenant passer de l'observation à l'\textbf{action}. Savoir reconnaître une injustice ne suffit pas : le citoyen numérique responsable est celui qui \textbf{choisit d'agir} pour protéger la dignité d'autrui, réparer ses propres erreurs et contribuer à un climat de confiance sur le réseau. Cette section te donne les repères concrets et la méthode pour transformer l'empathie en gestes quotidiens.

\courssub{Respecter autrui : les fondations du vivre-ensemble numérique}

Le respect en ligne ne diffère pas du respect dans la cour de récréation ou au marché : il s'agit de \textbf{reconnaître l'autre comme une personne qui a des droits}, et non comme un objet ou un simple pseudonyme anonyme. Trois piliers structurent ce respect : le consentement, le langage correct et le respect des opinions différentes.

\begin{definition}[Consentement à la publication]
Le \cle{consentement} est l'accord libre, éclairé et explicite d'une personne avant que son image, sa voix, ses propos ou ses données personnelles ne soient diffusés sur un espace public ou semi-public (groupe WhatsApp, story Instagram, page Facebook, TikTok). L'absence de refus ne vaut pas consentement.
\end{definition}

\begin{exemplebox}[Photo de classe partagée sans accord]
Lors de la fête de fin d'année, ton ami Kevin prend une photo de groupe où l'on voit Amina en train de manger, bouche ouverte, l'air fatigué. Kevin poste la photo sur le groupe WhatsApp de la classe « 3\up{e} A » avec la légende « La team qui a trop mangé ». Amina n'a pas été consultée. Elle se sent humiliée.
\begin{itemize}
    \item \textbf{Non éthique :} publication sans consentement, moquerie publique.
    \item \textbf{Éthique :} Kevin demande à Amina (et aux autres) : « Ça te va si je poste cette photo ? » Si Amina refuse, Kevin ne poste pas. Point final.
\end{itemize}
\end{exemplebox}

\begin{propriete}[Langage correct et opinions différentes]
Sur un espace numérique, le \cle{langage correct} exclut les insultes, le langage SMS excessif qui nuit à la compréhension, les majuscules criardes (QUI ÉQUIVAUDRAIENT À CRIER) et les propos discriminatoires (origine, genre, handicap, religion, niveau scolaire). \textbf{Respecter les opinions différentes} signifie accepter qu'un camarade pense autrement que toi sur un sujet (choix d'orientation, avis sur un jeu, opinion personnelle) sans l'attaquer personnellement.
\end{propriete}

\begin{attention}[Piège de l'« humour » et du « second degré »]
Écrire « C'est une blague » ou « Je rigole » après une remarque blessante \textbf{ne l'efface pas}. L'écrit ne transmet pas le ton de la voix. Ce qui te semble drôle peut être perçu comme une agression par celui qui le reçoit, surtout s'il est fragile ou déjà harcelé par ailleurs. \textbf{Règle d'or :} si tu dois expliquer que c'est une blague, c'est que ce n'en était pas une pour la victime.
\end{attention}

\courssub{Protéger autrui : ne pas rester spectateur face au harcèlement}

Lorsque tu es témoin d'une situation d'injustice (harcèlement, diffusion de photo intime, propos haineux), ton silence n'est pas neutre : il \textbf{renforce l'agresseur} en lui donnant l'impression que son comportement est toléré par le groupe.

\begin{propriete}[Effet du silence du témoin]
Le silence du témoin renforce l'agresseur ; l'intervention du témoin protège la victime. Ce principe, validé par la psychologie sociale (effet du spectateur, en anglais \emph{bystander effect}), s'applique pleinement aux espaces numériques où l'audience est large et la responsabilité diluée.
\end{propriete}

\begin{definition}[Témoin actif]
Un \cle{témoin actif} (en anglais \emph{bystander}) est une personne qui assiste à une situation d'injustice ou de souffrance en ligne et \textbf{choisit d'intervenir positivement} : en soutenant la victime, en signalant le contenu, en alertant un adulte de confiance ou un modérateur, sans pour autant s'exposer inutilement au danger.
\end{definition}

\begin{popfigure}
\begin{tikzpicture}[
    box/.style={rectangle, rounded corners, draw=popdark, thick, fill=popcream, text width=4.6cm, align=center, minimum height=1.6cm},
    arrow/.style={-Stealth, thick},
    scene/.style={rectangle, rounded corners, draw=poporange, thick, fill=poporangeL, text width=4.8cm, align=center, minimum height=1.8cm},
    passive/.style={rectangle, rounded corners, draw=poppink, thick, fill=poppinkL, text width=4.2cm, align=center, minimum height=1.8cm},
    active/.style={rectangle, rounded corners, draw=popgreen, thick, fill=popgreenL, text width=4.2cm, align=center, minimum height=1.8cm},
]
% Scène centrale
\node[scene] (scene) {\textbf{Scène :} un élève poste une moquerie méchante sur le groupe de classe : « Regardez les chaussures de Pierre, quelle honte ! »};

% Chemin passif (gauche)
\node[passive, left=1.2cm of scene, yshift=2.2cm, align=center] (passif1){\textbf{Spectateur passif}\\\emph{« Je ne veux pas de problèmes »}};
\node[passive, left=1.2cm of scene, yshift=-1.2cm, align=center] (passif2){Regarde sans réagir\\Ne signale pas\\Ne parle pas à Pierre};
\node[passive, below=0.6cm of passif2] (conseqP) {\textbf{Conséquence :} Pierre se sent seul, l'agresseur recommence, le groupe banalise la méchanceté};

\draw[arrow, poppink] (scene.west) -- ++(-0.4,0) |- (passif1.east);
\draw[arrow, poppink] (passif1.south) |- (passif2.west);
\draw[arrow, poppink] (passif2) -- (conseqP);

% Chemin actif (droite)
\node[active, right=1.2cm of scene, yshift=2.2cm, align=center] (actif1){\textbf{Témoin actif}\\\emph{« Je peux aider »}};
\node[active, right=1.2cm of scene, yshift=-1.2cm, align=center] (actif2){1. Soutient Pierre en privé\\2. Signale le message\\3. Alerte le professeur principal\\4. Ne like pas, ne partage pas};
\node[active, below=0.6cm of actif2] (conseqA) {\textbf{Conséquence :} Pierre se sent soutenu, le message est supprimé, l'agresseur comprend que c'est inacceptable};

\draw[arrow, popgreen] (scene.east) -- ++(0.4,0) |- (actif1.west);
\draw[arrow, popgreen] (actif1.south) |- (actif2.east);
\draw[arrow, popgreen] (actif2) -- (conseqA);

% Légendes des chemins
\node[above=0.15cm of passif1, text=poppink, font=\bfseries\small] {Inaction};
\node[above=0.15cm of actif1, text=popgreen, font=\bfseries\small] {Action citoyenne};
\end{tikzpicture}
\legende{Spectateur passif ou témoin actif : deux choix face à la même scène de moquerie. Les conséquences sur le climat de classe sont radicalement différentes.}
\end{popfigure}

\begin{methode}[Comment intervenir concrètement sans se mettre en danger]
\begin{enumerate}
    \item \textbf{Ne participe pas :} ne like pas, ne commente pas, ne transfère pas le message harcelant.
    \item \textbf{Soutiens la victime en privé :} envoie un message direct : « J'ai vu ce qui s'est passé, c'est injuste, tu n'as rien fait de mal. » Cela brise l'isolement immédiat.
    \item \textbf{Signale le contenu :} utilise les outils de la plateforme (bouton « Signaler »). C'est anonyme et efficace.
    \item \textbf{Alerte un adulte de confiance :} professeur principal, censeur, parent, éducateur. Tu n'es pas un « rapporteur », tu es un \textbf{protecteur}.
    \item \textbf{Ne confronte pas l'agresseur seul publiquement :} cela peut dégénérer en « guerre de commentaires ». Laisse l'autorité (modérateur, adulte) gérer la sanction.
\end{enumerate}
\end{methode}

\courssub{Réparer un tort : l'art de l'excuse sincère et de la correction}

Nous commettons tous des erreurs : un commentaire trop vif, une photo partagée trop vite, une rumeur relayée sans vérifier. Le citoyen numérique se distingue par sa capacité à \textbf{reconnaître sa faute et à la corriger}.

\begin{exemplebox}[Relai d'une fausse information]
Tu as partagé sur ton statut WhatsApp une capture d'écran annonçant : « \textbf{Annulation du BEPC suite à une fuite massive des sujets !} » avec le logo du MINESEC. Deux heures plus tard, le communiqué officiel du Ministère dément formellement cette information.
\begin{itemize}
    \item \textbf{Mauvaise réaction :} ne rien faire, espérer que cela s'oublie, ou dire « Je l'ai juste partagé, ce n'est pas moi qui l'ai écrit ».
    \item \textbf{Bonne réaction (réparation) :}
    \begin{enumerate}
        \item \textbf{Supprimer} le statut erroné immédiatement.
        \item \textbf{Publier un démenti clair} au même endroit : « RECTIFICATIF : l'information que j'ai partagée sur l'annulation du BEPC est FAUSSE. Le MINESEC a démenti. Désolé pour l'inquiétude causée. Vérifiez toujours sur les canaux officiels du MINESEC. »
        \item \textbf{S'excuser} auprès de ceux que l'on a inquiétés (camarades, parents) si nécessaire.
    \end{enumerate}
\end{itemize}
\end{exemplebox}

\begin{propriete}[Les 3 R de la réparation numérique]
\begin{enumerate}
    \item \textbf{Reconnaître :} « J'ai eu tort, j'ai publié cela sans vérifier / j'ai blessé X par ce commentaire. »
    \item \textbf{Rétablir :} supprimer le contenu offensant ou faux ; publier la correction (démenti, excuse) \textbf{au même endroit} et avec une visibilité au moins équivalente.
    \item \textbf{Réparer la relation :} s'excuser sincèrement auprès de la personne concernée (en privé si c'était personnel, en public si c'était public), sans condition (« Je m'excuse \textbf{si} tu t'es vexé » n'est pas une vraie excuse).
\end{enumerate}
\end{propriete}

\courssub{L'empathie appliquée : la pause avant de publier}

L'empathie n'est pas un sentiment vague, c'est une \textbf{compétence qui s'entraîne} : la capacité à se mettre à la place de l'autre. Avant toute publication (message, commentaire, story, transfert), impose-toi une pause de quelques secondes.

\begin{methode}[La règle des 3 questions avant de publier]
Avant d'appuyer sur « Envoyer », « Publier » ou « Transférer », pose-toi ces trois questions dans l'ordre :
\begin{enumerate}
    \item \textbf{Est-ce VRAI ?} Ai-je vérifié la source ? Est-ce une rumeur, une satire, une fausse information ? (Si non $\rightarrow$ \textbf{ne pas publier}.)
    \item \textbf{Est-ce RESPECTUEUX ?} Est-ce que cela blesse, humilie, exclut ou discrimine quelqu'un ? Est-ce que j'accepterais qu'on me le dise ? (Si non $\rightarrow$ \textbf{ne pas publier}.)
    \item \textbf{Est-ce UTILE ?} Cela apporte-t-il une information, un soutien, une joie, une réflexion constructive ? Ou est-ce du bruit, de la provocation ? (Si non $\rightarrow$ \textbf{s'abstenir}.)
\end{enumerate}
\textbf{Seulement si les 3 réponses sont OUI}, tu publies en toute conscience.
\end{methode}

\begin{aretenir}[L'empathie sous forme d'algorithme]
Cette méthode transforme l'empathie en \textbf{algorithme de décision} simple, que tu peux appliquer mentalement :
\begin{lstlisting}
Algorithme Decision_Publication
Debut
    Si (information NON verifiee) Alors
        Ecrire("Stop : information non verifiee")
    Sinon
        Si (contenu NON respectueux) Alors
            Ecrire("Stop : risque de blesser quelqu'un")
        Sinon
            Si (contenu NON utile) Alors
                Ecrire("Stop : publication inutile")
            Sinon
                Ecrire("Publication autorisee")
            FinSi
        FinSi
    FinSi
Fin
\end{lstlisting}
Applique cet algorithme mentalement avant chaque publication : il devient un réflexe au bout de quelques semaines.
\end{aretenir}

\courssub{Les 5 gestes citoyens quotidiens : une routine numérique saine}

Au-delà des situations de crise, la citoyenneté numérique se construit par de \textbf{petits gestes répétés}. Voici une frise des 5 habitudes à ancrer dès la classe de 3\up{e}.

\begin{popfigure}
\begin{tikzpicture}[
    stepnode/.style={rectangle, rounded corners, draw=popblue, thick, fill=popblueL, text width=2.6cm, align=center, minimum height=2.4cm, font=\small},
    arrownode/.style={-Stealth, thick, draw=popblue, shorten >=2pt, shorten <=2pt},
    num/.style={circle, fill=popblue, text=white, font=\bfseries\small, minimum size=0.6cm, inner sep=0},
]
% Nœuds de la frise
\node[stepnode, align=center] (g1){\textbf{VÉRIFIER}\\avant de partager\\Source officielle ?\\Date ? Auteur ?};
\node[num, left=0.05cm of g1.north west, anchor=south east] {1};

\node[stepnode, right=0.4cm of g1, align=center] (g2){\textbf{CITER}\\ses sources\\Auteur + lien\\+ date};
\node[num, left=0.05cm of g2.north west, anchor=south east] {2};

\node[stepnode, right=0.4cm of g2, align=center] (g3){\textbf{SIGNALER}\\les contenus haineux\\Bouton « Signaler »\\anonyme et rapide};
\node[num, left=0.05cm of g3.north west, anchor=south east] {3};

\node[stepnode, right=0.4cm of g3, align=center] (g4){\textbf{SOUTENIR}\\les victimes\\Message privé\\Alerter un adulte};
\node[num, left=0.05cm of g4.north west, anchor=south east] {4};

\node[stepnode, right=0.4cm of g4, align=center] (g5){\textbf{RÉFLÉCHIR}\\3 questions\\Vrai ? Respectueux ?\\Utile ?};
\node[num, left=0.05cm of g5.north west, anchor=south east] {5};

% Flèches
\draw[arrownode] (g1) -- (g2);
\draw[arrownode] (g2) -- (g3);
\draw[arrownode] (g3) -- (g4);
\draw[arrownode] (g4) -- (g5);

% Boucle de retour (habitude)
\draw[arrownode, dashed] (g5.south) to[bend right=25] node[below, font=\footnotesize, text=popmuted] {Routine quotidienne} (g1.south);
\end{tikzpicture}
\legende{Frise des 5 gestes citoyens quotidiens. Ces micro-habitudes, répétées, construisent un espace numérique plus sûr pour tous.}
\end{popfigure}

\begin{savaistu}[Le signalement au Cameroun : accessible à tous, y compris aux mineurs]
Au Cameroun, les principales plateformes utilisées (Facebook, WhatsApp, Instagram, TikTok) disposent de \textbf{boutons « Signaler » accessibles à tous les utilisateurs}, sans condition d'âge.
\begin{itemize}
    \item \textbf{WhatsApp :} appui long sur le message $\rightarrow$ « Signaler » (ou « Signaler et bloquer »). Le signalement est transmis de façon anonyme.
    \item \textbf{Facebook / Instagram :} les trois points (⋮) sur la publication, le commentaire ou le profil $\rightarrow$ « Signaler » $\rightarrow$ choisir la raison (harcèlement, discours de haine, fausse information, etc.).
    \item \textbf{TikTok :} bouton Partager (flèche) $\rightarrow$ « Signaler » $\rightarrow$ sélectionner la catégorie.
    \item \textbf{Cadre légal :} la \textbf{loi n\up{o} 2010/012 du 21 décembre 2010} relative à la cybersécurité et à la cybercriminalité au Cameroun punit notamment le harcèlement par voie électronique, la diffusion de fausses informations et l'atteinte à la vie privée. Signaler, c'est contribuer à faire respecter la loi.
\end{itemize}
\textbf{Astuce :} fais une capture d'écran (preuve) \textbf{avant} que le contenu ne soit supprimé par son auteur ou par le modérateur, surtout si un dépôt de plainte est envisagé.
\end{savaistu}

\courssub{Exemple résolu : analyse complète d'une situation vécue}

\begin{exoresolu}[Gestion d'une « story » qui dérape]
\textbf{Énoncé :} Chloé (3\up{e} C) poste une story Instagram filmant Maxime (3\up{e} C) qui trébuche dans l'escalier du collège. Elle ajoute une musique comique et écrit « Le clown de l'année ». La story est vue par 120 élèves. Maxime la voit, est très gêné et ne vient plus en cours pendant deux jours. Toi, tu as vu la story et tu connais Maxime. Que fais-tu \textbf{étape par étape}, en justifiant chaque action par les notions du cours ?
\tcblower
\textbf{Solution détaillée :}
\begin{enumerate}
    \item \textbf{Analyse éthique :} c'est du \cle{cyberharcèlement} : moquerie publique, large audience (120 élèves contre 1), atteinte à la dignité, absence de consentement de Maxime pour le filmage et la diffusion.
    \item \textbf{Action immédiate (témoin actif) :}
        \begin{itemize}
            \item \textbf{Ne pas réagir sur la story} (pas de like, pas de commentaire qui valide la moquerie).
            \item \textbf{Signaler la story} : ⋮ $\rightarrow$ Signaler $\rightarrow$ Harcèlement / intimidation. \emph{Justification :} le silence renforce l'agresseur ; le signalement est anonyme et relève du devoir de protection.
        \end{itemize}
    \item \textbf{Soutien à la victime (empathie concrète) :}
        \begin{itemize}
            \item \textbf{Message privé à Maxime :} « Salut, j'ai vu la story de Chloé. C'est injuste et méchant. Tu n'es pas un clown. Tu viens en cours demain ? Si tu as besoin de parler, je suis là. » \emph{Justification :} briser l'isolement, reconnaître l'émotion de Maxime, offrir une présence.
        \end{itemize}
    \item \textbf{Alerte à un adulte (responsabilité) :}
        \begin{itemize}
            \item \textbf{Parler au professeur principal ou au censeur} : « Une story qui harcèle Maxime circule ; il ne vient plus en cours. » \emph{Justification :} l'établissement a le devoir de protéger les élèves ; l'adulte a l'autorité pour sanctionner et médiatiser.
        \end{itemize}
    \item \textbf{Demande de réparation à l'auteure :}
        \begin{itemize}
            \item Si tu es proche de Chloé, \textbf{message privé calme} : « Chloé, ta story sur Maxime est du harcèlement. Il ne vient plus en cours. Tu dois la supprimer et t'excuser auprès de lui, sinon le professeur principal devra intervenir. » \emph{Justification :} donner une chance de réparer (les 3 R) avant une sanction plus lourde.
        \end{itemize}
    \item \textbf{Auto-évaluation :} j'applique la règle des 3 questions à cette story : Vrai ? Oui, il est tombé. Respectueux ? \textbf{Non}. Utile ? \textbf{Non}. $\rightarrow$ Je ne l'aurais jamais publiée ni partagée.
\end{enumerate}
\textbf{Conclusion :} l'enchaînement \textbf{Signaler $\rightarrow$ Soutenir $\rightarrow$ Alerter $\rightarrow$ Exiger la réparation} est la marque du citoyen numérique responsable.
\end{exoresolu}

\courssub{Synthèse : ton « permis de citoyen numérique »}

\begin{aretenir}[Check-list de l'attitude citoyenne en 3\up{e}]
\begin{center}
\begin{tabularx}{\linewidth}{|l|X|}
\hline
\rowcolor{popblueL}\textbf{Compétence} & \textbf{Indicateur de réussite (auto-évaluation)} \\
\hline
Respecter & Je demande \textbf{toujours} l'accord avant de poster une photo ou une information concernant autrui. Je n'insulte jamais, même « pour rire ». \\
\hline
Protéger & Je \textbf{signale} systématiquement le harcèlement dont je suis témoin. J'envoie un message de soutien à la victime. J'alerte un adulte. \\
\hline
Réparer & Si je me trompe (fausse information, propos blessant) : je \textbf{supprime}, je \textbf{rectifie} au même endroit, je \textbf{m'excuse} sincèrement. \\
\hline
Empathie active & J'applique la \textbf{règle des 3 questions} (Vrai ? Respectueux ? Utile ?) avant \textbf{chaque} publication. \\
\hline
Gestes quotidiens & Je vérifie les sources, je cite mes sources, je signale la haine, je soutiens les victimes, je réfléchis avant de publier. \\
\hline
\end{tabularx}
\end{center}
\end{aretenir}

\begin{exercice}[Mise en situation : le groupe WhatsApp « Délégués de 3\up{e} »]
Tu es délégué(e) de ta classe. Sur le groupe WhatsApp « Délégués de 3\up{e} » (élèves + professeur principal), un élève poste un montage photo grossier du professeur de mathématiques avec la légende « Le prof qui ne comprend rien lui-même ». Plusieurs élèves réagissent par des émoticônes de rire.
\begin{enumerate}
    \item Identifie les manquements éthiques (respect, consentement, rôle du délégué).
    \item En tant que délégué(e) \textbf{et} témoin actif, rédige le message \textbf{exact} que tu postes dans le groupe pour gérer la situation (ton ferme et respectueux, rappel du cadre, action demandée).
    \item Décris l'action que tu mènes \textbf{en parallèle} (hors du groupe) auprès de la victime (le professeur) et de l'auteur du montage.
    \item Applique la « règle des 3 questions » à la publication initiale de l'élève : quel est le résultat ?
\end{enumerate}
\end{exercice}

\begin{corrige}[Exercice : mise en situation « Délégués de 3\up{e} »]
\begin{enumerate}
    \item \textbf{Manquements :} atteinte à la dignité du professeur (harcèlement en ligne), absence de consentement pour le montage photo, violation du devoir d'exemplarité et de modération des délégués, banalisation de l'insulte par les réactions de rire.
    \item \textbf{Message au groupe :} \\
    « Bonjour à tous. En tant que délégués, nous avons pour rôle de représenter la classe dans le respect. Le montage photo du professeur de mathématiques est une insulte publique et constitue du harcèlement : ce n'est pas drôle, et c'est puni par la loi (loi n\up{o} 2010/012 sur la cybercriminalité). Je demande à l'auteur de le SUPPRIMER IMMÉDIATEMENT et de présenter ses excuses au professeur. Les réactions de rire cautionnent cet acte : je vous invite à les retirer. Si le contenu n'est pas supprimé d'ici une heure, je le signale à l'administration du groupe et au professeur principal. Merci de votre compréhension. »
    \item \textbf{Actions parallèles :}
        \begin{itemize}
            \item \textbf{Vers le professeur (victime) :} aller le voir discrètement : « Monsieur, nous avons vu le montage. C'est inacceptable. Nous vous soutenons et nous avons demandé sa suppression. »
            \item \textbf{Vers l'auteur :} message privé ferme mais éducatif : « Supprime le montage et excuse-toi. Si tu ne le fais pas, l'administration sanctionnera (conseil de discipline). Tu risques gros pour une soi-disant blague. »
        \end{itemize}
    \item \textbf{Règle des 3 questions appliquée à la publication initiale :}
        \begin{itemize}
            \item Vrai ? \textbf{Non} (montage, caricature mensongère).
            \item Respectueux ? \textbf{Non} (insulte, humiliation publique).
            \item Utile ? \textbf{Non} (détruit l'autorité du professeur et dégrade le climat de classe).
        \end{itemize}
        $\rightarrow$ \textbf{Publication interdite.}
\end{enumerate}
\end{corrige}

\begin{attention}[Erreur fréquente : « Je ne suis pas concerné, je n'ai fait que regarder »]
Le \textbf{simple fait de visionner une story harcelante sans réagir, de rester dans un groupe où circulent des photos intimes volées, ou de ne pas signaler un propos haineux} fait de toi un \textbf{complice passif}. La morale comme la loi camerounaise (loi n\up{o} 2010/012) exigent le signalement. \textbf{Ne pas agir, c'est choisir le camp de l'agresseur.} Sois le maillon qui brise la chaîne.
\end{attention}

\coursec{La démarche d'analyse d'une situation d'éthique (méthode ODARJ)}

Après avoir identifié les conduites non éthiques et compris ce qu'implique une attitude citoyenne, nous disposons maintenant d'une « boîte à outils » conceptuelle (éthique, empathie, loi, respect de la vie privée, propriété intellectuelle). Mais face à une situation réelle, complexe et souvent émotionnelle, comment passer de l'intuition (« je sens que ce n'est pas bien ») à une décision argumentée, solide et défendable ? C'est tout l'objet de cette section : t'offrir une \textbf{méthode structurée, reproductible et évaluée au BEPC} pour analyser n'importe quel dilemme numérique. Cette méthode, baptisée \textbf{ODARJ}, transforme ton jugement moral en une démarche intellectuelle rigoureuse.

\courssub{Pourquoi une méthode d'analyse ? Le piège de l'improvisation}

Imagine : tu reçois sur WhatsApp une photo embarrassante d'un camarade, prise à son insu dans les vestiaires du collège. Ton premier réflexe ? Peut-être la transférer à ton meilleur ami « pour rigoler », ou l'ignorer, ou la signaler. Sans méthode, tu réagis sous le coup de l'émotion (peur, amusement, pression du groupe) ou de l'habitude. Or, l'éthique numérique exige de la \textbf{réflexivité} : la capacité à mettre de la distance entre le stimulus (le message reçu) et ta réponse (ton action). La méthode ODARJ est ton « pare-feu » cognitif : elle t'oblige à ralentir, à structurer ta pensée et à produire une justification que tu pourras assumer devant ton principal, tes parents, ou toi-même dans dix ans.

\begin{attention}[Erreur fréquente — Sauter directement à la solution]
Beaucoup d'élèves (et d'adultes !) lisent une situation-problème et écrivent immédiatement : « Je signale au professeur » ou « Je ne fais rien ». \textbf{C'est insuffisant.} Au BEPC comme dans la vie, ce qui est noté et respecté, ce n'est pas seulement \emph{quoi} tu décides, mais \emph{pourquoi} et \emph{comment} tu y es parvenu. Une bonne décision mal justifiée vaut moins qu'une décision moyenne parfaitement argumentée. La méthode ODARJ te force à expliciter le cheminement invisible qui mène au choix final.
\end{attention}

\courssub{La méthode ODARJ : cinq étapes pour une décision citoyenne}

L'acronyme \cle{ODARJ} résume les cinq temps obligatoires de l'analyse. Chaque lettre correspond à une question clé et à un livrable écrit.

\begin{methode}[ODARJ — Observer, Diagnostiquer, Alternatives, Réagir, Justifier]
\begin{enumerate}
    \item \textbf{O — Observer (Décrire objectivement)} : \textit{« Que se passe-t-il vraiment ? »} Note les faits bruts, sans jugement, sans interprétation, sans émotion. Qui ? Quoi ? Où ? Quand ? Comment ? C'est le « procès-verbal » de la situation.
    \item \textbf{D — Diagnostiquer (Identifier les enjeux éthiques)} : \textit{« Qu'est-ce qui cloche ? »} Repère les conduites en jeu (partage non consenti, harcèlement, plagiat, intrusion, etc.) et les critères violés (respect de la vie privée, dignité, loi, propriété intellectuelle, empathie). Fais le lien avec le vocabulaire du cours.
    \item \textbf{A — Alternatives (Envisager les options)} : \textit{« Que puis-je faire ? »} Propose \textbf{au moins deux} réactions distinctes (ex. : « Je transfère » / « Je signale » / « J'ignore » / « J'en parle à la victime »). Pour chaque alternative, liste les conséquences probables \textbf{pour chaque acteur} (victime, auteur, témoins, toi, institution).
    \item \textbf{R — Réagir (Choisir et agir)} : \textit{« Que fais-je concrètement ? »} Sélectionne l'alternative la plus citoyenne. Décris l'action précise : à qui tu t'adresses, par quel canal, avec quels mots, dans quel délai. C'est le passage à l'acte responsable.
    \item \textbf{J — Justifier (Argumenter le choix)} : \textit{« Pourquoi ce choix est-il le meilleur ? »} Rédige un paragraphe structuré qui appuie ta décision sur : (1) une notion précise du cours, (2) un raisonnement logique (si... alors...), (3) l'analyse des conséquences (étape A). C'est la « preuve » de ta maturité numérique.
\end{enumerate}
\end{methode}

\begin{popfigure}
\begin{tikzpicture}[
    node distance=0.9cm,
    box/.style={rectangle, rounded corners, draw=popblue, thick, fill=popblueL, text width=3.6cm, align=center, minimum height=1.7cm, font=\small\bfseries},
    arrow/.style={-Stealth, thick, popdark},
    labelbox/.style={rectangle, draw=poporange, thick, fill=poporangeL, text width=5.2cm, align=left, minimum height=1.6cm, font=\small}
]
% Nœuds principaux (colonne de gauche)
\node[box, align=center] (O){\textbf{O — OBSERVER}\\\textit{Les faits bruts}\\(Qui ? Quoi ? Où ? Quand ?)};
\node[box, below=of O, align=center] (D){\textbf{D — DIAGNOSTIQUER}\\\textit{Enjeux \& Critères}\\(Vie privée, Loi, Empathie...)};
\node[box, below=of D, align=center] (A){\textbf{A — ALTERNATIVES}\\\textit{Options \& Conséquences}\\(Min. 2 scénarios comparés)};
\node[box, below=of A, align=center] (R){\textbf{R — RÉAGIR}\\\textit{Choix \& Action}\\(Concrète, datée, ciblée)};
\node[box, below=of R, align=center] (J){\textbf{J — JUSTIFIER}\\\textit{Argumentation}\\(Notion + Logique + Conséquences)};

% Flèches verticales
\draw[arrow] (O.south) -- (D.north);
\draw[arrow] (D.south) -- (A.north);
\draw[arrow] (A.south) -- (R.north);
\draw[arrow] (R.south) -- (J.north);

% Annotations latérales (colonne de droite, alignées sur les nœuds)
\node[labelbox, right=1.2cm of O, align=center] (lO){\textit{Question clé :} « Qu'ai-je vu/lu exactement ? »\\ \textit{Sortie :} description neutre, chronologique.};
\node[labelbox, right=1.2cm of D, align=center] (lD){\textit{Question clé :} « Quelle règle est brisée ? »\\ \textit{Sortie :} liste des violations (vie privée, secret des correspondances, charte...).};
\node[labelbox, right=1.2cm of A, align=center] (lA){\textit{Question clé :} « Quels sont mes choix ? »\\ \textit{Sortie :} tableau comparatif (Option / Acteurs / Conséquences).};
\node[labelbox, right=1.2cm of R, align=center] (lR){\textit{Question clé :} « Que fais-je maintenant ? »\\ \textit{Sortie :} plan d'action : « Je supprime le fichier », « Je préviens un adulte ».};
\node[labelbox, right=1.2cm of J, align=center] (lJ){\textit{Question clé :} « Pourquoi ce choix ? »\\ \textit{Sortie :} paragraphe argumenté type BEPC.};

% Flèches vers annotations
\draw[popmuted, dashed, -Stealth] (O.east) -- (lO.west);
\draw[popmuted, dashed, -Stealth] (D.east) -- (lD.west);
\draw[popmuted, dashed, -Stealth] (A.east) -- (lA.west);
\draw[popmuted, dashed, -Stealth] (R.east) -- (lR.west);
\draw[popmuted, dashed, -Stealth] (J.east) -- (lJ.west);
\end{tikzpicture}
\legende{Organigramme de la méthode ODARJ : une chaîne logique où chaque étape alimente la suivante. Sauter une étape fragilise la justification finale.}
\end{popfigure}

\courssub{Étape 1 — Observer : la « photographie » factuelle}

C'est l'étape la plus simple en apparence, mais la plus souvent bâclée. Observer, ce n'est pas donner son avis (« C'est méchant »), c'est établir un \textbf{constat} précis et vérifiable.

\begin{propriete}[Règles de l'observation objective]
\begin{itemize}
    \item \textbf{Exclusivité des faits :} seuls les éléments vérifiables (dates, heures, noms, propos exacts, captures d'écran) ont leur place.
    \item \textbf{Interdiction des adjectifs qualificatifs :} pas de « horrible », « drôle », « bête », « gentil ». Remplace par « photo prise sans consentement », « message insultant reçu à 14h23 ».
    \item \textbf{Exhaustivité contextuelle :} précise l'environnement (groupe WhatsApp « 3ème A », profil public TikTok, session de TP en salle informatique, cybercafé du quartier).
\end{itemize}
\end{propriete}

\begin{exemplebox}[Observation ratée vs Observation réussie]
\textbf{Situation :} Amine a publié sur TikTok une vidéo de son professeur qui s'énerve en classe.
\begin{itemize}
    \item \textcolor{poporange}{\textbf{Mauvaise observation :}} « Amine a fait une vidéo humiliante de notre prof qui pète un câble. C'est pas respectueux. »
    \item \textcolor{popgreen}{\textbf{Bonne observation :}} « Le 12 octobre 2025, vers 10h15, lors du cours de SVT en salle 12, Amine K. (élève de 3ème B) a enregistré sans autorisation une séquence de 43 secondes où M. Fouda (professeur) élève la voix face à une perturbation. La vidéo a été publiée à 10h22 sur le compte TikTok public @amine\_k\_237 avec la légende "Mon prof pète un plomb lol \#collège \#fail". 147 vues, 23 commentaires à 11h00. »
\end{itemize}
\end{exemplebox}

\courssub{Étape 2 — Diagnostiquer : qualifier les faits au regard de l'éthique et de la loi}

Ici, tu changes de casquette : tu deviens analyste. Tu prends les faits bruts (Étape O) et tu les \textbf{qualifies} à l'aide du vocabulaire précis de l'unité. C'est le moment de mobiliser les notions du cours : \cle{atteinte à la vie privée}, \cle{droit à l'image}, \cle{cyberharcèlement}, \cle{propriété intellectuelle}, \cle{usurpation d'identité}, \cle{non-respect des conditions d'utilisation}, \cle{manque d'empathie}.

\begin{definition}[Diagnostiquer]
Établir la correspondance entre les faits observés et les normes de référence (lois camerounaises, règlement intérieur de l'établissement, principes éthiques, règles de la nétiquette) pour identifier \textbf{quoi} est éthiquement inacceptable et \textbf{pourquoi}.
\end{definition}

\begin{methode}[Comment diagnostiquer efficacement]
\begin{enumerate}
    \item Relis ton observation (Étape O).
    \item Pour chaque fait marquant, demande-toi : « Quel droit, quelle valeur, quelle règle est violé(e) ? »
    \item Cite la référence : « le droit au respect de la vie privée », « la loi n°2010/012 du 21 décembre 2010 relative à la cybersécurité et à la cybercriminalité au Cameroun (qui sanctionne par exemple la diffusion d'images ou de données d'autrui sans consentement) », « le règlement intérieur du collège (interdiction de filmer en classe) », « le principe d'empathie : se mettre à la place du professeur filmé à son insu ».
    \item Identifie les \textbf{acteurs} : auteur (responsable principal), victime(s), témoins/complices (ceux qui likent, partagent, commentent), institution (collège, plateforme).
\end{enumerate}
\end{methode}

\courssub{Étape 3 — Alternatives : l'exercice de la pensée complexe}

C'est le cœur de la compétence « Rédiger et justifier une démarche ». Il est \textbf{interdit} de n'avoir qu'une seule option. La vie réelle offre toujours au moins deux chemins : l'action facile et immédiate (souvent non éthique) et l'action responsable (souvent plus coûteuse en effort). Tu dois les mettre en regard.

\begin{propriete}[Structure obligatoire du tableau d'alternatives]
Pour chaque alternative envisagée, analyse les conséquences sur \textbf{quatre dimensions} :
\begin{enumerate}
    \item \textbf{Pour la/les victime(s)} (protection, réparation, aggravation).
    \item \textbf{Pour l'auteur/autrice de la conduite non éthique} (sanction, prise de conscience, impunité).
    \item \textbf{Pour toi (l'observateur/acteur)} (charge mentale, risque disciplinaire, cohérence avec tes valeurs).
    \item \textbf{Pour le collectif / l'institution} (climat de confiance, précédent, image de l'établissement).
\end{enumerate}
\end{propriete}

\begin{attention}[Piège classique : l'alternative « fantôme »]
Ne propose pas une alternative irréaliste pour faire joli. Exemple : « Alternative 1 : Je signale. Alternative 2 : J'appelle le Président de la République. » $\rightarrow$ \textbf{Zéro point.} Les alternatives doivent être \textbf{réellement accessibles} à un élève de 3ème ici et maintenant (signaler à la vie scolaire, en parler à ses parents ou à un adulte de confiance, supprimer le contenu, ne rien faire, confronter l'auteur en privé, saisir les services compétents comme la police ou l'ANTIC en cas d'infraction grave, etc.).
\end{attention}

\courssub{Étape 4 — Réagir : l'engagement concret}

L'éthique sans action est stérile. Cette étape exige une \textbf{description opérationnelle} : pas « je vais signaler », mais « j'envoie un message au censeur ce soir avant 19h avec la capture d'écran en pièce jointe et j'en informe mon professeur principal demain matin à 7h30 ». La précision prouve que tu as vraiment l'intention d'agir.

\courssub{Étape 5 — Justifier : l'argumentation de niveau BEPC}

C'est l'étape qui vaut le plus de points en examen. Une justification n'est pas une opinion (« Je pense que c'est bien »). C'est une \textbf{démonstration}.

\begin{aretenir}[Les 3 piliers d'une justification recevable (critères BEPC)]
\begin{enumerate}
    \item \textbf{Référence explicite à une notion du cours :} « Ce choix s'appuie sur le principe de \textbf{l'empathie numérique} : se mettre à la place de la victime pour anticiper sa souffrance. » / « Il respecte le \textbf{droit à l'image} et la \textbf{vie privée}, protégés par la loi camerounaise. »
    \item \textbf{Argument logique (chaîne causale) :} « Si je ne signale pas (hypothèse), la vidéo reste en ligne (fait), la victime continue de subir le regard des autres (conséquence), l'auteur se sent légitimé (conséquence), le climat de classe se dégrade (conséquence). Donc, signaler rompt cette chaîne. »
    \item \textbf{Lien avec l'analyse des conséquences (Étape A) :} « L'alternative "ne rien faire" entraîne l'aggravation pour la victime (colonne 1 du tableau) et l'impunité pour l'auteur (colonne 2), ce qui contredit mon devoir de citoyen numérique. L'alternative choisie minimise les dommages et active la protection institutionnelle. »
\end{enumerate}
\end{aretenir}

\begin{exemplebox}[Justification type « 4/4 points »]
\textit{« Je choisis de signaler immédiatement à la vie scolaire et d'en informer mes parents (Réagir). Ce choix est justifié par trois raisons. (1) Il respecte le \textbf{droit à l'image} et la \textbf{vie privée} : M. Fouda n'a pas consenti à être filmé ni diffusé, et la loi n°2010/012 sur la cybersécurité et la cybercriminalité sanctionne la publication de données ou d'images d'autrui sans autorisation. (2) L'analyse des alternatives (Étape A) montre que "ne rien faire" laisse la vidéo se propager (147 vues en 38 minutes), amplifiant l'atteinte à la dignité du professeur et banalisant le non-respect de l'autorité pédagogique. (3) En signalant, j'active la \textbf{protection institutionnelle} (vie scolaire, administration, signalement à la plateforme) qui seule peut faire retirer le contenu rapidement et engager une médiation éducative, conformément à mon devoir de \textbf{citoyenneté numérique}. »}
\end{exemplebox}

\courssub{Application guidée : Le cas d'Amina (accroche de l'unité)}

Reprenons la situation d'accroche de l'unité pour dérouler la méthode ODARJ de bout en bout.

\begin{center}
\begin{minipage}{0.92\linewidth}
\textbf{Contexte (rappel) :} Amina, élève de 3ème, découvre sur le téléphone déverrouillé de sa meilleure amie Fatou, posé sur la table du CDI, des captures d'écran de conversations privées entre Fatou et un garçon. Amina est tentée de les lire, voire de les photographier pour les montrer à d'autres copines « pour rigoler ».
\end{minipage}
\end{center}

\vspace{0.5em}

\begin{experience}[Application complète ODARJ — Cas d'Amina]
\textbf{Objectif :} produire une analyse écrite complète, prête à être rendue en devoir ou au BEPC.

\textbf{Matériel :} feuille de brouillon, stylo, fiche de cours (notions : vie privée, secret des correspondances, empathie, loyauté).

\textbf{Manipulation (rédaction guidée) :}

\noindent\textbf{1. O — OBSERVER (faits bruts)}
\begin{itemize}
    \item \textbf{Quand :} mardi 18 novembre 2025, 14h30 (heure de permanence au CDI).
    \item \textbf{Où :} CDI du collège, table n°4.
    \item \textbf{Qui :} Amina (observatrice), Fatou (propriétaire du téléphone, amie proche), tiers non identifiés (destinataires potentiels).
    \item \textbf{Quoi :} le smartphone de Fatou (écran allumé, déverrouillé) est posé face visible. L'application « Galerie » est ouverte sur un album intitulé « Captures écran » contenant environ 12 images. Les vignettes montrent des bulles de messagerie (WhatsApp) avec du texte lisible (prénoms, confidences).
    \item \textbf{Action d'Amina (tentative) :} Amina s'approche, penche la tête, commence à lire le premier message. Elle hésite à prendre son propre téléphone pour photographier l'écran de Fatou.
\end{itemize}

\noindent\textbf{2. D — DIAGNOSTIQUER (enjeux éthiques et juridiques)}
\begin{center}
\begin{tabularx}{\linewidth}{|l|X|}\hline
\textbf{Conduite identifiée} & \textbf{Critère / Notion / Loi violée} \\\hline
Lecture de messages privés sans consentement & Atteinte au \textbf{secret des correspondances} et à la \textbf{vie privée} : les échanges privés d'autrui sont protégés par la loi camerounaise ; y accéder sans autorisation est une violation. \\\hline
Intention de photographier pour diffuser & \textbf{Atteinte au droit à l'image} et aux \textbf{données personnelles} (la loi n°2010/012 relative à la cybersécurité et à la cybercriminalité réprime la collecte et la diffusion illicites de données d'autrui). Rupture de la \textbf{loyauté} et de la \textbf{confiance} (amitié). \\\hline
Diffusion potentielle au groupe « copines » & Début de \textbf{cyberharcèlement} (diffusion de contenu intime). Manque total d'\textbf{empathie} : ignorance de la honte, de la trahison et de l'isolement que subirait Fatou. \\\hline
Non-respect de la charte du CDI & Règlement intérieur : « Interdiction d'utiliser le matériel d'autrui sans accord », « Respect de la vie privée des usagers ». \\\hline
\end{tabularx}
\end{center}

\noindent\textbf{3. A — ALTERNATIVES (tableau comparatif des conséquences)}
\begin{center}
\begin{tabularx}{\linewidth}{|l|X|X|X|X|}\hline
\textbf{Alternative} & \textbf{Fatou (victime)} & \textbf{Amina (acteur)} & \textbf{Groupe / Collectif} & \textbf{Institution (CDI/Collège)} \\\hline
\textbf{A1 : Lire, photographier, partager} & Trahison, humiliation publique, perte de confiance, risque de déprime, plainte possible. & Complice active, risque de sanction lourde (conseil de discipline, plainte), culpabilité durable, mauvaise réputation. & Climat toxique, banalisation du vol de données. & Échec de la mission éducative, image du collège dégradée, procédure disciplinaire lourde. \\\hline
\textbf{A2 : Ne rien dire, ne pas regarder, partir} & Vie privée préservée \emph{pour l'instant}, mais téléphone toujours exposé (risque qu'un tiers en profite). & Conscience tranquille (non-participation), mais \textbf{non-assistance} morale : Fatou reste vulnérable. & Aucun impact positif. Problème non résolu. & Aucun signalement, faille non corrigée. \\\hline
\textbf{A3 : Alerter Fatou discrètement, refermer le téléphone, n'en parler à personne} & \textbf{Protection immédiate}, prise de conscience de sa négligence (téléphone déverrouillé), confiance renforcée. & \textbf{Posture citoyenne} : loyauté, empathie, responsabilité. Modèle pour les pairs. Aucun risque disciplinaire. & Message fort : « On veille les uns sur les autres ». Climat de confiance. & Situation gérée à la source, prévention d'un incident futur. \\\hline
\textbf{A4 : Signaler au documentaliste / à la vie scolaire sans toucher au téléphone} & Protection institutionnelle, mais Fatou peut se sentir « dénoncée » par son amie (rupture du lien). & Décharge la responsabilité sur un adulte, mais perd l'occasion d'un geste d'amitié et d'autonomie. & Dépendance à l'autorité, moins d'autorégulation. & Procédure formelle engagée, trace administrative. \\\hline
\end{tabularx}
\end{center}

\noindent\textbf{4. R — RÉAGIR (choix et plan d'action)}

\textbf{Choix : Alternative A3 (alerter Fatou discrètement + sécuriser le téléphone).}

\textbf{Plan d'action concret :}
\begin{enumerate}
    \item \textit{Maintenant (14h32) :} je retourne le téléphone face contre table sans le déverrouiller davantage, sans lire la suite.
    \item \textit{Immédiatement :} je m'approche de Fatou (à 2 mètres, en train de chercher un livre) et je chuchote : « Fatou, ton téléphone est ouvert sur tes captures WhatsApp à la table 4. Je l'ai retourné. Verrouille-le tout de suite. »
    \item \textit{Dans l'heure :} je lui rappelle (en privé) l'importance du verrouillage automatique (code, schéma ou empreinte) et de ne pas laisser son téléphone sans surveillance.
    \item \textit{Si elle refuse d'écouter ou si le contenu a déjà fuité :} j'envisage l'Alternative A4 (signaler à un adulte de confiance) pour la protéger malgré elle.
\end{enumerate}

\noindent\textbf{5. J — JUSTIFIER (argumentation finale)}

\textit{« Je choisis l'alternative A3 (alerter Fatou discrètement et sécuriser son téléphone) pour trois raisons fondées sur les notions de l'unité. \\
\textbf{1. Respect de la vie privée et du secret des correspondances :} les messages privés de Fatou sont protégés par la loi. En refermant le téléphone sans lire ni photographier, je cesse immédiatement l'atteinte en cours et j'évite de commettre moi-même une infraction. Je traite les données de Fatou comme je voudrais qu'on traite les miennes (règle d'or de l'\textbf{empathie}). \\
\textbf{2. Analyse des conséquences (Étape A) :} le tableau montre que l'A3 est la \textbf{seule} option qui protège Fatou \textbf{et} préserve notre lien d'amitié (loyauté), tout en me positionnant en acteur responsable (citoyenneté). L'A1 est dévastatrice (sur le plan moral et même judiciaire). L'A2 est passive (Fatou reste en danger). L'A4, bien que légitime, brise la confiance entre amies et infantilise Fatou. \\
\textbf{3. Proportionnalité :} la réponse est proportionnée à la gravité (incident mineur encore contenu) : je gère au niveau le plus simple, entre pairs, avant d'escalader vers l'institution si nécessaire. C'est l'application concrète de l'\textbf{éthique de la responsabilité} : je réponds de mes actes et de leurs conséquences sur autrui. »}
\end{experience}

\begin{popfigure}
\begin{center}
\begin{tabularx}{\linewidth}{|>{\bfseries}l|X|}\hline
\textbf{Étape ODARJ} & \textbf{Contenu produit pour le cas d'Amina} \\\hline
\textbf{O — Observer} & Description factuelle : date, lieu, acteurs, état du téléphone (déverrouillé, galerie ouverte sur des captures WhatsApp), action tentée (lecture, photo). \textbf{Zéro adjectif subjectif}. \\\hline
\textbf{D — Diagnostiquer} & 4 violations identifiées : (1) secret des correspondances, (2) vie privée / données personnelles (loi n°2010/012), (3) loyauté / empathie (valeurs), (4) charte du CDI (règlement). Acteurs : Amina (auteur potentiel), Fatou (victime), tiers (risque). \\\hline
\textbf{A — Alternatives} & 4 scénarios comparés (A1 Partage, A2 Fuite, A3 Alerte discrète, A4 Signalement à un adulte) sur 4 dimensions (Victime, Acteur, Collectif, Institution). A3 sort optimal. \\\hline
\textbf{R — Réagir} & Choix : A3. Plan d'action précis : 1. Retourner l'écran, 2. Alerter Fatou à l'oral (chuchotement), 3. Conseil de sécurité (verrouillage automatique), 4. Clause de repli (A4 si échec). \\\hline
\textbf{J — Justifier} & Paragraphe structuré en 3 piliers : (1) référence aux notions (vie privée, empathie), (2) logique causale via le tableau A (A3 seul protège + loyauté), (3) principes éthiques (proportionnalité, responsabilité). \\\hline
\end{tabularx}
\end{center}
\legende{Synthèse de l'application ODARJ au cas d'Amina. Ce tableau constitue une « fiche modèle » à reproduire en examen.}
\end{popfigure}

\courssub{Entraînement BEPC : le barème de la méthode}

\begin{exemplebox}[Sur 20 points au BEPC, une réponse à situation-problème structurée en 5 étapes ODARJ peut valoriser jusqu'à 4 points de méthode]
\begin{center}
\begin{tabular}{|c|c|p{8cm}|}\hline
\textbf{Étape} & \textbf{Points} & \textbf{Attendus pour le point maximal} \\\hline
O & 0,5 & Faits uniquement, 5W (Qui, Quoi, Où, Quand, Comment), pas d'opinion. \\\hline
D & 1,0 & Au moins 3 violations distinctes nommées avec le vocabulaire du cours + référence à une loi ou à une règle. \\\hline
A & 1,0 & Minimum 3 alternatives réalistes + tableau de conséquences sur 4 dimensions (Victime, Auteur, Soi, Collectif). \\\hline
R & 0,5 & Action précise, datée, adressée, avec clause de repli si pertinent. \\\hline
J & 1,0 & Paragraphe unique intégrant : (1) notion du cours citée, (2) lien logique Si/Alors, (3) appui sur le tableau A. \\\hline
\textbf{Total Méthode} & \textbf{4,0} & \textit{(Le fond / la justesse du choix vaut les points restants de l'exercice, souvent 6 à 10 points selon l'énoncé.)} \\\hline
\end{tabular}
\end{center}
\end{exemplebox}

\begin{savaistu}[L'éthique par la méthode : une longue tradition]
La méthode ODARJ n'est pas une invention scolaire. Elle s'inspire de la \textbf{casuistique} (étude des cas de conscience) pratiquée depuis des siècles en philosophie morale, et reprise aujourd'hui en éthique médicale (méthode des « 4 principes » : autonomie, bienfaisance, non-malfaisance, justice) et dans les métiers de l'ingénieur, où l'on apprend à \textbf{formaliser} un dilemme avant de trancher. Au Cameroun, les grandes écoles de formation technique intègrent des modules de déontologie où l'on enseigne exactement cela : \textbf{ne pas décider sous le coup de l'émotion, mais décider par la raison éclairée par la règle}. ODARJ est ta version « collège » de cette exigence professionnelle.
\end{savaistu}

\courssub{Synthèse de la section : ton « algorithme » moral}

\begin{aretenir}[Ce qu'il faut retenir de la méthode ODARJ]
\begin{itemize}
    \item \textbf{O — Observer} : faire le « procès-verbal » des faits (Qui, Quoi, Où, Quand, Comment). \textit{Sortie : description neutre.}
    \item \textbf{D — Diagnostiquer} : qualifier les faits avec le vocabulaire du cours (vie privée, secret des correspondances, cyberharcèlement, propriété intellectuelle, empathie, loi). \textit{Sortie : liste des violations + acteurs.}
    \item \textbf{A — Alternatives} : construire un tableau comparatif (min. 3 options $\times$ 4 dimensions de conséquences). \textit{Sortie : tableau décisionnel.}
    \item \textbf{R — Réagir} : choisir l'option optimale et écrire le plan d'action concret (Quoi, Qui, Quand, Comment). \textit{Sortie : plan opérationnel.}
    \item \textbf{J — Justifier} : rédiger l'argumentation en 3 piliers : \textbf{notion du cours} + \textbf{logique Si/Alors} + \textbf{preuve par le tableau A}. \textit{Sortie : paragraphe argumenté noté 4/4.}
\end{itemize}
\textbf{Ordre strict :} O $\rightarrow$ D $\rightarrow$ A $\rightarrow$ R $\rightarrow$ J. On ne justifie (J) que ce qu'on a choisi (R), on ne choisit (R) qu'après avoir comparé (A), on ne compare (A) qu'après avoir diagnostiqué (D), on ne diagnostique (D) qu'après avoir observé (O).
\end{aretenir}

\begin{exercice}[Entraînement ODARJ — Cas « Le devoir de maths »]
\textbf{Situation :} Kévin, en 3ème, n'a pas fait son devoir de maths (noté sur 20, coefficient 3). Son camarade Yannick, premier de la classe, lui envoie la photo de son devoir parfait par WhatsApp à 6h du matin, juste avant le cours : « Copie vite, efface après. » Kévin recopie en 10 minutes et rend sa copie. Le professeur, en corrigeant, remarque une erreur \textbf{identique, bizarre et spécifique} (même faute de signe à la ligne 4, même abréviation « dc » au lieu de « donc ») sur les deux copies.

\textbf{Consigne :} applique la méthode ODARJ complète pour analyser la situation du point de vue de \textbf{Kévin} (au moment où il reçoit le message à 6h00) puis du point de vue du \textbf{professeur} (au moment de la correction). Rédige l'analyse complète sur une feuille (O, D, A, R, J pour chaque acteur).
\end{exercice}

\begin{corrige}[Exercice n°1 — Le devoir de maths (Éléments de correction)]
\textbf{Point de vue de Kévin (6h00) :}
\begin{itemize}
    \item \textbf{O :} réception d'un message WhatsApp à 6h00, photo du devoir de Yannick, instruction « copie/efface ». Devoir à rendre à 8h00. Kévin n'a pas travaillé.
    \item \textbf{D :} \textbf{fraude / tricherie} (règlement intérieur), \textbf{non-respect du travail d'autrui} (le devoir de Yannick est le fruit de son effort), \textbf{manque d'intégrité} (honnêteté intellectuelle), \textbf{risque disciplinaire} (zéro, conseil de discipline).
    \item \textbf{A :} A1 : Copier (gain de note, mais risque de zéro si détecté, trahison de soi-même, injustice envers les autres). A2 : Refuser, dire « Non merci, je rends page blanche » (note de zéro assumée, honneur sauvé, respect du professeur). A3 : Répondre « Je n'ai pas fait, j'assume » + demander de l'aide pour comprendre après (responsabilité, apprentissage).
    \item \textbf{R :} choix A3. Action : répondre immédiatement « Non, je ne triche pas. Je rends blanc. On révise ensemble ce week-end ? ». Rendre une copie blanche à 8h00.
    \item \textbf{J :} choix A3 justifié par l'\textbf{intégrité} (valeur fondamentale de la citoyenneté numérique) : « Mieux vaut un zéro honnête qu'un 20 volé ». Logique : si je triche (A1), je valide mon incompétence, je risque une sanction lourde et je commets une injustice collective. A3 assume l'échec, préserve l'estime de soi et ouvre sur une remédiation. Le tableau A montre que A3 est la seule option sans conséquence négative durable.
\end{itemize}
\textbf{Point de vue du professeur (correction) :}
\begin{itemize}
    \item \textbf{O :} deux copies (Kévin, Yannick) avec une erreur signature identique (même faute de signe, même abréviation « dc »). Devoir coefficient 3.
    \item \textbf{D :} \textbf{fraude avérée} (indices convergents). Violation du règlement sur les devoirs. Atteinte à l'\textbf{équité} (justice entre les élèves).
    \item \textbf{A :} A1 : Mettre zéro aux deux et signaler au conseil de discipline (sanction lourde, mais Yannick est peut-être victime d'un envoi détourné). A2 : Mettre zéro à Kévin seul et avertir Yannick « ne partage plus » (juste pour Kévin, mais sans enquête). A3 : Entendre les deux élèves séparément puis ensemble, puis décider (procédure contradictoire et éducative).
    \item \textbf{R :} choix A3 (pédagogie + justice). Convocation de Kévin et Yannick. Sanction : zéro pour Kévin, avertissement pour Yannick (négligence dans le partage), information des parents des deux.
    \item \textbf{J :} A3 respecte le \textbf{droit de s'expliquer} (contradictoire) et la \textbf{finalité éducative} de la sanction (comprise, pas seulement subie). Le tableau A montre que A1 est injuste pour Yannick (double peine possible) et A2 trop expéditif (pas d'enquête). A3 seul permet de distinguer l'auteur (Kévin) du facilitateur (Yannick) et d'éduquer les deux.
\end{itemize}
\end{corrige}

\begin{attention}[Dernier rappel avant la section suivante]
La méthode ODARJ n'est pas une « usine à gaz » réservée à l'examen. C'est ton \textbf{processus mental par défaut} face à tout dilemme numérique : un message suspect qui te demande tes identifiants (hameçonnage) ? $\rightarrow$ ODARJ. Une rumeur sur les réseaux sociaux ? $\rightarrow$ ODARJ. Une demande de mot de passe d'un « ami » ? $\rightarrow$ ODARJ. Plus tu t'exerces à l'écrire (sur le cas d'Amina, sur l'exercice de Kévin, sur tes propres situations vécues), plus elle devient automatique et rapide, et plus elle te protège, toi et les autres. La section suivante te proposera une synthèse globale et une auto-évaluation pour valider tes acquis de l'unité.
\end{attention}

\coursec{Rédiger et justifier une réponse argumentée}

Après avoir analysé une situation avec la méthode \textbf{ODARJ} (Observer, Définir le problème, Analyser les enjeux, Rechercher des solutions, Juger et agir), tu dois maintenant \textbf{communiquer ton jugement} de façon claire, structurée et convaincante. Au BEPC comme dans la vie citoyenne, une bonne décision mal expliquée perd de sa valeur. Cette section t'apprend à transformer ton analyse en un texte argumenté rigoureux, que l'on puisse lire, évaluer et transmettre.

\courssub{La structure TAEC : le squelette de ton argumentation}

Tout paragraphe argumenté solide repose sur quatre piliers, dans cet ordre immuable : \textbf{Thèse} $\rightarrow$ \textbf{Argument} $\rightarrow$ \textbf{Exemple} $\rightarrow$ \textbf{Conclusion}. C'est le plan \textbf{TAEC}. Respecter cet ordre, c'est guider ton lecteur (le correcteur, ton professeur, ton interlocuteur) sans le perdre.

\begin{popfigure}
\begin{tikzpicture}[
    node distance=1.2cm and 0.8cm,
    funnel/.style={trapezium, trapezium left angle=70, trapezium right angle=110, minimum width=3cm, minimum height=1.2cm, draw=popdark, thick, fill=popblueL, text=popdark, font=\bfseries\small, align=center},
    arrow/.style={-Stealth, thick, popdark},
    labelstyle/.style={font=\small, text=popmuted, anchor=north, yshift=-0.2cm}
]
\node[funnel, align=center] (thesis){THÈSE\\Affirmation claire};
\node[labelstyle] at (thesis.south) {« Je refuse de partager cette photo »};

\node[funnel, below=of thesis, align=center] (argument){ARGUMENT\\Justification théorique};
\node[labelstyle] at (argument.south) {« Atteinte à la dignité, violation éthique »};

\node[funnel, below=of argument, align=center] (example){EXEMPLE\\Ancrage concret / Empathie};
\node[labelstyle] at (example.south) {« À sa place, je serais humilié »};

\node[funnel, below=of example, align=center] (conclusion){CONCLUSION\\Engagement citoyen};
\node[labelstyle] at (conclusion.south) {« Je supprime et je soutiens la victime »};

\draw[arrow] (thesis.south) -- (argument.north);
\draw[arrow] (argument.south) -- (example.north);
\draw[arrow] (example.south) -- (conclusion.north);

\node[left=0.3cm of thesis.west, font=\bfseries\tiny, text=poporange] {1.};
\node[left=0.3cm of argument.west, font=\bfseries\tiny, text=poporange] {2.};
\node[left=0.3cm of example.west, font=\bfseries\tiny, text=poporange] {3.};
\node[left=0.3cm of conclusion.west, font=\bfseries\tiny, text=poporange] {4.};
\end{tikzpicture}
\legende{Schéma en entonnoir TAEC : la thèse large s'affine par l'argument, se concrétise par l'exemple, aboutit à l'engagement.}
\end{popfigure}

\begin{methode}[Le plan TAEC — Thèse, Argument, Exemple, Conclusion]
\noindent\textbf{Étape 1 — Thèse (Affirmation).} Énonce ta position en une phrase courte, directe, sans ambiguïté. C'est la réponse à la question « Que décides-tu ? ».
\smallskip
\noindent\textbf{Étape 2 — Argument (Justification).} Explique \emph{pourquoi} cette thèse est la bonne. Mobilise les notions du cours : \cle{éthique}, \cle{empathie}, \cle{respect d'autrui}, \cle{responsabilité}, \cle{cyberharcèlement}, \cle{droit à l'image}. Si la situation relève du pénal, cite la \textbf{loi n°2010/012 du 21 décembre 2010 relative à la cybersécurité et la cybercriminalité au Cameroun}.
\smallskip
\noindent\textbf{Étape 3 — Exemple / Empathie (Ancrage).} Rends l'argument vivant. Mets-toi à la place de la victime (\cle{empathie}) ou donne un cas concret observé. C'est la preuve que tu as \emph{ressenti} la situation, pas seulement compris la règle.
\smallskip
\noindent\textbf{Étape 4 — Conclusion citoyenne (Engagement).} Propose une action concrète, réaliste, immédiate : « Je supprime le fichier », « Je signale au modérateur ou au professeur », « J'accompagne la victime chez le conseiller d'orientation », « Je m'engage à ne plus liker de contenus humiliants ». Pas de vœu pieux (« Il faut que tout le monde soit gentil »), mais \emph{ton} acte.
\end{methode}

\courssub{Les connecteurs logiques : le ciment du raisonnement}

Un texte sans connecteurs est une suite de phrases disjointes. Les connecteurs montrent la \textbf{logique} qui relie tes idées. Utilise-les \emph{systématiquement}.

\begin{center}
\begin{tabularx}{\linewidth}{|l|X|X|}
\hline
\rowcolor{popblueL}
\textbf{Fonction} & \textbf{Connecteurs recommandés} & \textbf{Exemple d'usage} \\
\hline
\textbf{Introduire la thèse} & \textbf{D'abord,} \textbf{En premier lieu,} & D'abord, je refuse de transférer ce message. \\
\hline
\textbf{Enchaîner l'argument} & \textbf{Ensuite,} \textbf{Par ailleurs,} \textbf{De plus,} & Ensuite, ce partage violerait le droit à l'image. \\
\hline
\textbf{Expliquer (cause)} & \textbf{Car,} \textbf{En effet,} \textbf{Parce que,} & Car la loi camerounaise protège l'intimité. \\
\hline
\textbf{Illustrer} & \textbf{Par exemple,} \textbf{Ainsi,} \textbf{C'est le cas quand} & Par exemple, si c'était ma photo, je me sentirais trahi. \\
\hline
\textbf{Conclure / Conséquence} & \textbf{Par conséquent,} \textbf{C'est pourquoi,} \textbf{Aussi,} & C'est pourquoi je le supprime immédiatement. \\
\hline
\textbf{Nuancer / Opposition} & \textbf{Toutefois,} \textbf{Cependant,} \textbf{Néanmoins,} & Toutefois, je ne dois pas insulter l'expéditeur. \\
\hline
\end{tabularx}
\end{center}

\begin{attention}[Erreur fréquente — La « liste de courses » sans fil conducteur]
\noindent Beaucoup d'élèves écrivent : « Je ne partage pas. C'est mal. La loi l'interdit. Moi je n'aimerais pas. Donc je ne le fais pas. » \textbf{Ce n'est pas une argumentation, c'est une énumération.} Il manque les \textbf{connecteurs} (\texttt{car}, \texttt{c'est pourquoi}) et surtout le \textbf{lien explicite} entre la notion du cours (ex. \cle{empathie}) et l'exemple personnel. Le correcteur ne doit pas \emph{deviner} ton raisonnement : tu dois l'\textbf{écrire}.
\end{attention}

\courssub{Justifier par les notions du cours et par la loi}

Ton argumentation doit prouver que tu as \textbf{maîtrisé le programme}. Les mots-clés suivants doivent apparaître \emph{naturellement} dans ton texte, pas collés artificiellement :

\begin{itemize}
    \item \textbf{\cle{Éthique}} : ensemble de principes moraux guidant l'usage du numérique (respect, honnêteté, responsabilité).
    \item \textbf{\cle{Empathie}} : capacité à se mettre à la place d'autrui pour ressentir sa souffrance ou sa joie.
    \item \textbf{\cle{Respect d'autrui}} : reconnaissance de la dignité, de la vie privée, du droit à l'image de l'autre.
    \item \textbf{\cle{Responsabilité}} : assumer les conséquences de ses actes en ligne (partager, liker, commenter engage ta responsabilité).
    \item \textbf{\cle{Cyberharcèlement}} / \textbf{\cle{Atteinte à l'intimité}} : qualifications précises des conduites interdites.
\end{itemize}

\begin{propriete}[Référence légale — Loi camerounaise n°2010/012]
\noindent Quand la situation implique une infraction (photo intime volée, moqueries répétées, usurpation d'identité, menace), tu peux citer la \textbf{loi n°2010/012 du 21 décembre 2010 relative à la cybersécurité et la cybercriminalité au Cameroun}, qui punit notamment :
\begin{itemize}
    \item les \textbf{atteintes à la vie privée} commises par un moyen de communication électronique (diffusion de photos, vidéos ou correspondances privées sans consentement) ;
    \item le \textbf{harcèlement} exercé à travers les technologies de l'information et de la communication (messages répétés, insultes, menaces) ;
    \item l'\textbf{usurpation d'identité} et la diffusion de contenus illicites.
\end{itemize}
\noindent\emph{Formulation type :} « Cette action constitue une atteinte à la vie privée, punie par la loi n°2010/012 du 21 décembre 2010 relative à la cybersécurité et la cybercriminalité au Cameroun. »
\end{propriete}

\courssub{Rédiger une conclusion citoyenne : de l'intention à l'acte}

La conclusion n'est pas une formule de politesse. C'est ton \textbf{engagement personnel}. Elle doit être :
\begin{enumerate}
    \item \textbf{Concrète} : une action faisable \emph{maintenant} (supprimer, signaler, parler à un adulte de confiance, bloquer l'auteur).
    \item \textbf{Protectrice} : elle vise à stopper le dommage et à soutenir la victime.
    \item \textbf{Non-violente} : pas de menace de représailles, pas d'insulte en retour. La citoyenneté numérique, c'est aussi le \cle{self-control}.
\end{enumerate}

\begin{exemplebox}[Modèle complet TAEC — Situation : on t'envoie une photo humiliante d'un camarade]
\noindent\textbf{Thèse :} \textbf{Je refuse catégoriquement de transférer cette photo à qui que ce soit.}
\smallskip
\noindent\textbf{Argument :} \textbf{En effet,} diffuser une image dégradante sans consentement porte atteinte à la \cle{dignité} et au \cle{droit à l'image} de mon camarade. C'est une violation de l'\cle{éthique} numérique et de l'\cle{empathie} : nous avons la \cle{responsabilité} de protéger les plus vulnérables. \textbf{Par ailleurs,} sur le plan juridique, cet acte constitue une atteinte à la vie privée punie par la \textbf{loi n°2010/012} relative à la cybersécurité et la cybercriminalité au Cameroun.
\smallskip
\noindent\textbf{Exemple :} \textbf{Par exemple,} si c'était ma propre photo qui circulait ainsi devant tout le collège, je me sentirais humilié, trahi, et j'aurais peur d'aller en cours. L'\cle{empathie} m'oblige à ressentir cette souffrance à sa place.
\smallskip
\noindent\textbf{Conclusion citoyenne :} \textbf{C'est pourquoi} je supprime immédiatement le fichier de mon téléphone, je bloque l'expéditeur, je signale l'incident à mon professeur principal et j'apporte mon soutien moral à la victime en l'accompagnant voir le conseiller d'orientation.
\end{exemplebox}

\courssub{Correction guidée d'une copie-type : repérer les forces et les oublis}

Voici une production d'élève face au sujet : \emph{« Un élève de ta classe a créé un faux compte au nom d'un professeur pour poster des insultes. Tu découvres l'auteur. Que fais-tu ? Justifie ta démarche. »}

\begin{popfigure}
\begin{tikzpicture}[
    node distance=0.4cm,
    box/.style={rectangle, draw, rounded corners=3pt, thick, inner sep=4pt, text width=\linewidth-16pt, font=\small},
    greenbox/.style={box, draw=popgreen, fill=popgreenL, text=popdark},
    orangebox/.style={box, draw=poporange, fill=poporangeL, text=popdark},
    redbox/.style={box, draw=popink, fill=poppinkL, text=popdark},
    labelstyle/.style={font=\tiny\bfseries, anchor=north west, xshift=4pt, yshift=-2pt}
]
\node[greenbox] (thesis) {\textbf{Thèse (correcte) :} Je dénonce l'auteur à l'administration et je refuse de liker ou partager les publications.};
\node[labelstyle, text=popgreen] at (thesis.north west) {FORCE : Position claire, action immédiate.};

\node[orangebox, below=of thesis] (arg) {\textbf{Argument (moyen) :} C'est pas bien de faire ça. Le prof a le droit au respect. C'est du cyberharcèlement.};
\node[labelstyle, text=poporange] at (arg.north west) {MANQUE : Vocabulaire familier (« C'est pas bien »). Notions du cours absentes (éthique, empathie, responsabilité). Loi non citée.};

\node[redbox, below=of arg] (ex) {\textbf{Exemple (faux) :} Moi aussi j'ai déjà eu un faux compte.};
\node[labelstyle, text=popink] at (ex.north west) {ERREUR : Exemple hors sujet (expérience de l'élève comme auteur, pas comme victime ou témoin). Pas d'empathie pour le professeur.};

\node[orangebox, below=of ex] (conc) {\textbf{Conclusion (moyenne) :} Du coup j'arrête.};
\node[labelstyle, text=poporange] at (conc.north west) {FAIBLE : « Du coup » est oral. Action floue (« j'arrête » quoi ?). Pas d'engagement de signalement concret ni de soutien à la victime.};

\draw[popgreen, thick, decorate, decoration={brace, amplitude=6pt}] (thesis.north east) -- (thesis.south east) node[midway, right=8pt, text=popgreen, font=\tiny\bfseries] {VALIDÉ};
\draw[poporange, thick, decorate, decoration={brace, amplitude=6pt}] (arg.north east) -- (conc.south east) node[midway, right=8pt, text=poporange, font=\tiny\bfseries] {À REVOIR};
\end{tikzpicture}
\legende{Copie annotée : zones vertes (acquis), orange (incomplet ou imprécis), rose (erreur ou hors-sujet).}
\end{popfigure}

\begin{exoresolu}[Rédaction corrigée et commentée]
\noindent\textbf{Sujet :} Un élève a créé un faux compte au nom d'un professeur pour poster des insultes. Tu découvres l'auteur. Que fais-tu ? Justifie ta démarche.
\tcblower
\noindent\textbf{Thèse :} \textbf{Je signale immédiatement l'usurpation d'identité et le harcèlement à la direction de l'établissement.}
\smallskip
\noindent\textbf{Argument :} \textbf{En effet,} créer un faux compte pour insulter au nom d'un tiers constitue une \textbf{usurpation d'identité numérique} et un \textbf{cyberharcèlement} envers le professeur. Cela viole les principes d'\cle{éthique} (honnêteté, respect) et d'\cle{empathie} : la victime subit une atteinte à sa \cle{réputation} et à son \cle{autorité} pédagogique. \textbf{De plus,} sur le plan légal, la \textbf{loi n°2010/012 du 21 décembre 2010} punit l'usurpation d'identité, les atteintes à la vie privée et le harcèlement par voie électronique. Ma \cle{responsabilité} de témoin est d'agir pour faire cesser le trouble.
\smallskip
\noindent\textbf{Exemple :} \textbf{Par exemple,} imaginons que ce professeur soit mon père ou mon grand frère : le voir humilié publiquement, incapable de se défendre car l'auteur se cache, provoquerait en moi de la colère et de la honte. L'\cle{empathie} me commande de ne pas laisser cet acte impuni.
\smallskip
\noindent\textbf{Conclusion citoyenne :} \textbf{C'est pourquoi} je remets une capture d'écran des preuves au proviseur (sans la diffuser moi-même), j'encourage la victime à signaler les faits aux autorités compétentes et je m'engage à sensibiliser ma classe sur les dangers de l'usurpation d'identité lors de la prochaine heure de vie de classe.
\end{exoresolu}

\begin{aretenir}[Check-list avant de rendre ta copie]
\noindent Avant de considérer ta réponse comme finie, vérifie ces 5 points :
\begin{enumerate}
    \item \textbf{Thèse} : une phrase, affirmative, qui répond directement à la question.
    \item \textbf{Argument} : au moins \textbf{deux notions du cours} citées explicitement (\texttt{éthique}, \texttt{empathie}, \texttt{responsabilité}...) + \textbf{la loi} si infraction (loi n°2010/012).
    \item \textbf{Exemple} : une mise en situation \textbf{personnelle} (« À ma place... », « Si c'était mon frère... ») qui prouve l'\cle{empathie}, pas un fait divers vague.
    \item \textbf{Connecteurs} : \texttt{En effet}, \texttt{Par ailleurs}, \texttt{Par exemple}, \texttt{C'est pourquoi} présents et bien placés.
    \item \textbf{Conclusion} : une action \textbf{concrète, immédiate, non-violente, protectrice}.
\end{enumerate}
\end{aretenir}

\begin{savaistu}[Le savais-tu ? — Le barème invisible du BEPC]
\noindent Au BEPC, dans une question portant sur l'éthique numérique, la \textbf{justification} (le « pourquoi ») compte souvent \textbf{plus que la décision} (le « quoi »). Une réponse « Je signale » sans argumentation vaut peu de points. La même réponse structurée en TAEC, avec notions du cours, référence à la loi, empathie et connecteurs, vaut la quasi-totalité des points. \textbf{Le correcteur cherche la trace de ton raisonnement citoyen, pas seulement ton bon sens.}
\end{savaistu}

\begin{exercice}[Entraînement — À toi de jouer]
\noindent\textbf{Situation :} Dans le groupe WhatsApp de la classe, un camarade poste un montage photo grossier d'une fille de la classe avec la légende « La reine du bahut ». Tout le monde rit et like. Tu sais que cette fille est timide et qu'elle a déjà pleuré pour moins que ça.
\smallskip
\noindent\textbf{Consigne :} Rédige ta réponse complète en respectant \textbf{strictement le plan TAEC}. Utilise au moins \textbf{3 connecteurs logiques}, cite \textbf{2 notions du cours} et \textbf{la loi camerounaise} si applicable. Ton engagement final doit être \textbf{réalisable demain matin au collège}.
\end{exercice}

\begin{corrige}[Exercice — Proposition de correction]
\noindent\textbf{Thèse :} Je ne like pas, je ne commente pas, et je signale le message aux administrateurs du groupe ainsi qu'au professeur principal.
\smallskip
\noindent\textbf{Argument :} \textbf{En effet,} ce montage constitue une \cle{atteinte à la dignité} et au \cle{droit à l'image} de ma camarade. Le fait que le groupe rie aggrave le \cle{cyberharcèlement} par effet de meute. L'\cle{éthique} numérique m'interdit de cautionner par mon silence ou mon like une humiliation publique. \textbf{Par ailleurs,} la diffusion d'un montage dégradant sans consentement constitue une atteinte à la vie privée et un harcèlement punis par la \textbf{loi n°2010/012 du 21 décembre 2010} relative à la cybersécurité et la cybercriminalité au Cameroun.
\smallskip
\noindent\textbf{Exemple :} \textbf{Par exemple,} connaissant sa timidité, je sais qu'elle ne s'en remettra pas facilement ; à sa place, je me sentirais détruit de voir mon visage déformé validé par des dizaines de « j'aime ». L'\cle{empathie} m'impose de briser cette chaîne de moqueries.
\smallskip
\noindent\textbf{Conclusion citoyenne :} \textbf{C'est pourquoi} je demande la suppression immédiate du message, j'envoie un message privé de soutien à la victime pour lui dire qu'elle n'est pas seule, et j'alerte le professeur principal dès la rentrée des classes demain matin avec une capture d'écran horodatée.
\end{corrige}

\coursec{Synthèse et auto-évaluation citoyenne}

Dans la section précédente, tu as appris à rédiger et justifier une réponse argumentée grâce à la méthode TAEC (Thèse, Argument, Exemple, Conclusion). Tu sais désormais analyser une situation d'éthique numérique avec la démarche ODARJ et défendre ton point de vue par écrit. Mais une leçon d'éthique n'a de valeur que si elle transforme ta conduite réelle, chaque jour, sur ton téléphone ou ton ordinateur. Cette dernière section est donc un moment de bilan et d'engagement : nous allons rassembler toutes les notions de l'unité, évaluer honnêtement ta propre conduite numérique, construire ensemble une charte pour le groupe WhatsApp de la classe, et formuler un engagement personnel concret. C'est le couronnement de l'unité : passer du \emph{savoir} au \emph{savoir-être}.

\courssub{Bilan des notions clés de l'unité}

Avant de t'évaluer, reprenons calmement les cinq piliers de cette unité. Chacun répond à une question précise de la vie quotidienne.

\begin{aretenir}[Les cinq piliers de l'unité]
\begin{itemize}
\item \cle{L'éthique} : ensemble de règles de conduite qui nous aident à distinguer le bien du mal, y compris dans le monde numérique. Question clé : \emph{« Est-ce bien ce que je m'apprête à faire ? »}
\item \cle{L'empathie} : capacité à se mettre à la place d'autrui, à ressentir ce que l'autre peut ressentir. Question clé : \emph{« Comment réagirais-je si on me faisait cela ? »}
\item \cle{La citoyenneté numérique} : ensemble des comportements responsables, respectueux et solidaires d'un internaute envers les autres et envers la communauté. Question clé : \emph{« Suis-je un bon citoyen sur Internet ? »}
\item \cle{La démarche ODARJ} : méthode d'analyse d'une situation en cinq étapes — Observer, Décrire, Analyser, Réagir, Justifier.
\item \cle{La rédaction TAEC} : méthode pour construire une réponse argumentée — Thèse, Argument, Exemple, Conclusion.
\end{itemize}
\end{aretenir}

Pour mémoriser l'ensemble d'un seul coup d'œil, rien ne vaut une carte mentale : elle montre comment les notions sont reliées entre elles autour d'une idée centrale.

\begin{popfigure}
\begin{tikzpicture}[mindmap, grow cyclic,
every node/.style={concept, font=\small},
root concept/.append style={concept color=popblue, font=\bfseries\small, minimum size=2.6cm},
level 1 concept/.append style={level distance=3.4cm, sibling angle=72, font=\footnotesize}]
\node[root concept]{Éthique et empathie}
child[concept color=poporange] { node {Éthique : distinguer le bien du mal} }
child[concept color=popgreen] { node {Empathie : se mettre à la place d'autrui} }
child[concept color=poppurple] { node {Citoyenneté numérique : agir en citoyen responsable} }
child[concept color=poppink] { node {ODARJ : Observer, Décrire, Analyser, Réagir, Justifier} }
child[concept color=popgold] { node {TAEC : Thèse, Argument, Exemple, Conclusion} };
\end{tikzpicture}
\legende{Carte mentale de synthèse de l'unité : cinq notions reliées autour d'une idée centrale.}
\end{popfigure}

Remarque la logique d'ensemble : l'\emph{éthique} et l'\emph{empathie} sont des attitudes intérieures ; la \emph{citoyenneté numérique} est leur traduction en actes ; \emph{ODARJ} et \emph{TAEC} sont les outils qui te permettent d'analyser une situation puis de rédiger ta réponse, notamment au BEPC.

\courssub{Grille d'auto-évaluation de ma conduite numérique}

Pourquoi s'auto-évaluer ? Parce qu'on ne peut progresser que si l'on sait où l'on en est. Un footballeur qui ignore ses points faibles ne s'améliore pas ; un internaute non plus. L'auto-évaluation consiste à se regarder honnêtement dans un miroir, sans se juger sévèrement, mais sans se mentir.

\begin{definition}[Auto-évaluation]
L'\cle{auto-évaluation} est une démarche par laquelle on examine soi-même ses propres comportements à l'aide de critères précis, afin d'identifier ses forces et ses points à améliorer.
\end{definition}

Voici une grille de dix critères couvrant les grandes situations de la vie numérique d'un élève de 3ème. Pour chaque ligne, coche la case qui correspond à ta conduite réelle : \textbf{toujours}, \textbf{parfois} ou \textbf{jamais}. Sois honnête : personne ne te notera, cette grille est pour toi.

\begin{center}
\begin{tabularx}{\linewidth}{|c|X|c|c|c|}
\hline
\rowcolor{popblueL} \textbf{N°} & \textbf{Critère de conduite} & \textbf{Toujours} & \textbf{Parfois} & \textbf{Jamais} \\
\hline
1 & Je demande le \textbf{consentement} d'une personne avant de publier sa photo. & $\square$ & $\square$ & $\square$ \\
\hline
2 & J'utilise un \textbf{langage respectueux} dans mes messages, même quand je suis en colère. & $\square$ & $\square$ & $\square$ \\
\hline
3 & Je vérifie les \textbf{sources} d'une information avant de la croire. & $\square$ & $\square$ & $\square$ \\
\hline
4 & Je réfléchis avant de \textbf{relayer} (partager) un message ou une vidéo. & $\square$ & $\square$ & $\square$ \\
\hline
5 & Je \textbf{soutiens les victimes} de moqueries en ligne au lieu de rire ou de partager. & $\square$ & $\square$ & $\square$ \\
\hline
6 & Je respecte la \textbf{vie privée} des autres (je ne lis pas leurs messages, je ne diffuse pas leurs secrets). & $\square$ & $\square$ & $\square$ \\
\hline
7 & Je ne crée pas de \textbf{faux comptes} et je n'usurpe pas l'identité d'autrui. & $\square$ & $\square$ & $\square$ \\
\hline
8 & Je respecte les \textbf{auteurs} : je ne copie pas un travail en le présentant comme le mien. & $\square$ & $\square$ & $\square$ \\
\hline
9 & Je \textbf{signale} les contenus dangereux ou haineux à un adulte ou à la plateforme. & $\square$ & $\square$ & $\square$ \\
\hline
10 & Je maîtrise mon \textbf{temps d'écran} pour préserver mes études et ma santé. & $\square$ & $\square$ & $\square$ \\
\hline
\end{tabularx}
\end{center}

\begin{methode}[Comment interpréter ma grille]
\begin{enumerate}
\item Compte tes cases « toujours » : ce sont tes \textbf{forces}. Félicite-toi sincèrement.
\item Repère tes cases « parfois » : ce sont tes \textbf{axes de progrès} prioritaires, car tu es déjà sur le bon chemin.
\item Regarde tes cases « jamais » sans paniquer : choisis-en \textbf{une seule} sur laquelle travailler ce mois-ci. Vouloir tout changer d'un coup mène à l'échec.
\item Refais la grille dans un mois et compare : la progression est la seule vraie mesure du succès.
\end{enumerate}
\end{methode}

\begin{attention}[L'éthique ne concerne pas que les « grands problèmes »]
Erreur fréquente : croire que l'éthique numérique ne s'applique qu'aux grands délits (piratage, escroquerie, cyberharcèlement grave). Faux ! L'éthique se joue dans \textbf{chaque message quotidien} : transférer la photo d'un camarade sans son accord, répondre par une moquerie dans le groupe de classe, partager une rumeur « juste pour informer »... Ce sont ces petits gestes répétés qui font de toi un bon ou un mauvais citoyen numérique. Au BEPC, un sujet peut très bien porter sur une situation ordinaire de la vie de classe.
\end{attention}

\courssub{Élaboration d'une charte de classe pour le groupe WhatsApp}

La plupart des classes de 3ème possèdent un groupe WhatsApp pour s'échanger les devoirs, les leçons et les informations. Mais sans règles, ces groupes dérivent vite : moqueries, photos gênantes, rumeurs, messages à toute heure. La solution n'est pas de supprimer le groupe, mais de le doter d'une \emph{charte}.

\begin{definition}[Charte d'usage]
Une \cle{charte} est un document écrit, accepté par tous les membres d'un groupe, qui fixe les règles de bonne conduite à respecter dans un cadre commun (ici, un groupe de discussion en ligne).
\end{definition}

\begin{methode}[Rédiger une charte efficace]
\begin{enumerate}
\item Limiter la charte à \textbf{5 règles maximum} : au-delà, personne ne les retient.
\item Formuler chaque règle \textbf{positivement} : on écrit « je vérifie avant de partager » plutôt que « ne pas partager de fausses informations ». Une règle positive dit quoi faire ; une interdiction dit seulement quoi éviter.
\item Rédiger à la \textbf{première personne} (« je ») : chacun s'engage personnellement.
\item Faire \textbf{adopter la charte par toute la classe} (vote ou signature) : une règle imposée sans accord est une règle ignorée.
\item \textbf{Afficher la charte} en tête du groupe (description du groupe) et la relire régulièrement.
\end{enumerate}
\end{methode}

\begin{exemplebox}[Charte-type du groupe WhatsApp de la classe de 3ème — Collège de Ngaoundéré]
\textbf{Notre charte — adoptée par toute la classe le 12 janvier :}
\begin{enumerate}
\item Je respecte chaque membre : mes messages restent polis et bienveillants, même en cas de désaccord.
\item Je vérifie une information avant de la partager, et je cite ma source.
\item Je demande l'accord d'une personne avant de publier sa photo ou ses propos.
\item J'utilise le groupe d'abord pour les études : devoirs, leçons, entraide et informations utiles.
\item Je soutiens tout camarade moqué ou attaqué, et je signale les dérives à l'administrateur ou à un enseignant.
\end{enumerate}
\emph{Signée : les 42 élèves de la classe et le professeur principal.}
\end{exemplebox}

\begin{experience}[TP guidé : construire la charte de notre classe]
\textbf{Objectif :} rédiger collectivement la charte du groupe de classe.\\
\textbf{Matériel :} tableau, craie ou marqueurs, cahier de textes, téléphone du délégué.\\
\textbf{Manipulation :}
\begin{enumerate}
\item Par groupes de quatre, listez trois problèmes déjà observés dans des groupes WhatsApp (moquerie, fausse information, photo publiée sans accord...).
\item Pour chaque problème, proposez une règle \emph{positive} qui le résout.
\item En classe entière, votez pour retenir les cinq meilleures règles.
\item Rédigez la charte finale au propre, faites-la signer par tous, puis le délégué la copie dans la description du groupe.
\end{enumerate}
\textbf{Observations :} notez les désaccords éventuels lors du vote et la façon dont ils ont été tranchés.\\
\textbf{Interprétation :} une règle décidée ensemble est mieux respectée qu'une règle imposée : c'est le principe même de la citoyenneté.
\end{experience}

\courssub{Engagement personnel : une résolution concrète et mesurable}

Une bonne résolution ne ressemble pas à « je serai plus gentil sur Internet » : c'est vague, invérifiable, donc vite oublié. Pour être efficace, un engagement doit être \emph{concret} (on sait exactement quoi faire) et \emph{mesurable} (on peut vérifier si on l'a tenu).

\begin{propriete}[Critères d'un bon engagement]
Un engagement personnel est efficace s'il est :
\begin{itemize}
\item \textbf{concret} : il décrit une action précise (« je demande l'accord avant de publier une photo ») ;
\item \textbf{mesurable} : il contient une quantité ou une fréquence vérifiable (« à chaque fois », « une fois par semaine ») ;
\item \textbf{daté} : il fixe un moment de bilan (« je ferai le point dans un mois »).
\end{itemize}
\end{propriete}

\begin{exoresolu}[Formuler un engagement personnel]
\textbf{Énoncé.} Amina, élève de 3ème, a coché « parfois » au critère n°4 de la grille (réfléchir avant de relayer un message). Aide-la à transformer cette constatation en engagement personnel efficace.
\tcblower
\textbf{Solution détaillée.} Une mauvaise résolution serait : « je ferai attention aux fausses informations » (ni concrète, ni mesurable). Appliquons les trois critères :
\begin{itemize}
\item \emph{Concret} : préciser l'action — vérifier la source et l'auteur avant tout partage.
\item \emph{Mesurable} : préciser la fréquence — avant \textbf{chaque} partage, sans exception.
\item \emph{Daté} : fixer un bilan — dans un mois, refaire la grille d'auto-évaluation.
\end{itemize}
\textbf{Engagement rédigé :} « Pendant un mois, avant de partager le moindre message dans un groupe, je vérifierai qui en est l'auteur et je chercherai une source fiable. Si je ne trouve pas de source, je ne partage pas. Je referai ma grille d'auto-évaluation le 15 du mois prochain pour vérifier que je suis passée de "parfois" à "toujours". »
\end{exoresolu}

\begin{exercice}[À toi de jouer]
Choisis le critère de ta grille où tu as obtenu ton moins bon résultat. Rédige ton engagement personnel en respectant les trois critères (concret, mesurable, daté), puis présente-le à ton voisin qui vérifiera qu'il est bien mesurable.
\end{exercice}

\courssub{Complément utile au BEPC : éthique numérique et éducation aux médias}

\begin{savaistu}[Les 3 réflexes de la vérification de l'information]
La \cle{vérification des sources} (en anglais \emph{fact-checking}) est le pont entre l'éthique numérique et l'éducation aux médias. Elle repose sur trois réflexes simples :
\begin{enumerate}
\item \textbf{Identifier l'auteur} : qui a écrit ou publié cette information ? Un site connu, un journaliste, ou un compte anonyme ?
\item \textbf{Croiser avec une source fiable} : est-ce qu'un média sérieux ou un site officiel (ministère, organisme reconnu) dit la même chose ?
\item \textbf{Vérifier la date} : une vraie information ancienne ressortie de son contexte peut devenir une fausse information aujourd'hui.
\end{enumerate}
Au Cameroun comme ailleurs, les rumeurs circulent vite sur les réseaux sociaux (santé, examens, sécurité). Avant de relayer une « information » sur le BEPC ou sur un événement local, applique ces trois réflexes : c'est déjà de la citoyenneté numérique en action.
\end{savaistu}

Ce complément est directement utile au BEPC : un sujet d'analyse de situation peut te présenter une rumeur propagée dans un groupe de classe et te demander de réagir. Avec ODARJ pour analyser, TAEC pour rédiger, et les trois réflexes de vérification pour argumenter, tu disposes de tout l'attirail nécessaire.

\begin{aretenir}[L'essentiel de la section]
\begin{itemize}
\item L'unité repose sur cinq piliers reliés : éthique, empathie, citoyenneté numérique, démarche ODARJ, rédaction TAEC.
\item La grille d'auto-évaluation (10 critères) permet de mesurer honnêtement sa conduite et de choisir un axe de progrès.
\item Une charte de groupe efficace compte au plus 5 règles, formulées positivement et adoptées par tous.
\item Un engagement personnel doit être concret, mesurable et daté.
\item Vérifier les sources (auteur, recoupement, date) est un réflexe citoyen indispensable, au quotidien comme au BEPC.
\end{itemize}
\end{aretenir}

Tu as maintenant terminé cette unité. Souviens-toi : derrière chaque écran se trouve une personne réelle, avec ses émotions et sa dignité. La technologie change vite, mais le respect d'autrui reste la règle d'or. Sois, en ligne comme dans la cour du collège, le citoyen que tu aimerais rencontrer.

\newpage

\coursec{Activité d'intégration}

\coursec{Leçon 11 : Application cumulée de l'éthique et de l'empathie}

\courssub{Objectifs de l'activité}

À l'issue de cette activité d'intégration, tu seras capable de :

\begin{itemize}
\item mobiliser les définitions structurantes de l'unité (\cle{éthique}, \cle{empathie}, \cle{citoyenneté numérique}) pour analyser une situation réaliste de la vie courante ;
\item identifier, à l'aide des quatre critères étudiés, les conduites non éthiques présentes dans un cas concret de \cle{cyberharcèlement} ;
\item appliquer de manière autonome la méthode \cle{ODARJ} (Observer, Décrire, Analyser, Réagir, Justifier) pour proposer et comparer plusieurs réactions possibles ;
\item adopter et défendre une \cle{attitude citoyenne} vis-à-vis d'autrui, y compris lorsque cela demande du courage ;
\item rédiger et justifier une réponse argumentée selon la structure \cle{TAEC} (Thèse, Argument, Exemple, Conclusion) ;
\item coopérer au sein d'un groupe pour élaborer des règles communes de vie numérique (charte).
\end{itemize}

\courssub{Situation}

\textbf{Affaire n° 2026-03 : le groupe WhatsApp de la 3\textsuperscript{ème} B.}

Au collège Bilingue de Bonabéri, à Douala, la classe de 3\textsuperscript{ème} B possède un groupe WhatsApp intitulé « 3e B – Solidarité », créé en septembre pour s'échanger les devoirs, les horaires et les annonces. Le groupe compte \textbf{42 membres}, tous élèves de la classe, et \textbf{2 administrateurs} (le délégué et son adjoint).

Mardi soir, vers 20 h 15, \textbf{Junior}, élève de la classe, utilise une application de retouche d'images pour modifier une photo de \textbf{Amina}, sa camarade, prise pendant la fête de la jeunesse. Il grossit exagérément son visage et ajoute un texte moqueur. À 20 h 18, il publie l'image dans le groupe avec le message : « Regardez la star de la classe ! »

En 25 minutes, voici ce qui se passe :

\begin{itemize}
\item \textbf{17 élèves} ont vu le message (marque « lu ») ;
\item \textbf{Moussa} et \textbf{Carine} relaient l'image : Moussa la transfère dans le groupe des garçons de la classe (12 membres supplémentaires) et Carine répond par trois emojis qui rient aux larmes, suivis du commentaire « elle ressemble à un ballon » ;
\item \textbf{Kevin} voit le message, se dit que « ce n'est pas ses affaires » et ne fait rien ;
\item \textbf{Sandrine}, amie d'Amina, est choquée. Elle hésite : répondre dans le groupe risque de la faire passer pour celle qui « gâche l'ambiance », mais elle sait qu'Amina, qui n'est pas connectée ce soir-là, découvrira tout au réveil ;
\item à 20 h 43, \textbf{12 commentaires} moqueurs au total ont été publiés sous l'image.
\end{itemize}

Le lendemain matin, Amina arrive au collège en pleurs : une camarade d'un autre établissement lui a envoyé une capture d'écran. L'image a quitté le groupe d'origine. La vie scolaire est saisie de l'affaire, et la classe entière est invitée à une séance de réflexion sur la citoyenneté numérique.

\textbf{La classe est transformée en tribunal citoyen du numérique.} Par groupes de 4, vous êtes les « assesseurs » chargés d'analyser l'affaire, de proposer la réaction la plus citoyenne que Sandrine aurait pu adopter, et de rédiger une charte pour que cela ne se reproduise plus.

\courssub{Consignes}

\textbf{Travail en groupe de 4 (étapes 1 et 2), puis travail individuel (étape 3), puis travail collectif (étape 4).}

\begin{enumerate}
\item \textbf{Décrire la situation.} Résumez les faits en 4 ou 5 phrases, en distinguant clairement : les victimes, les auteurs, les relais, les témoins. Qui sont les protagonistes ? Quel outil numérique est en cause ?

\item \textbf{Identifier les conduites non éthiques.} Pour chacun des quatre élèves suivants (Junior, Moussa, Carine, Kevin), indiquez si sa conduite est éthique ou non, en vous appuyant sur les \textbf{quatre critères} vus en leçon : (a) le respect de la dignité d'autrui, (b) l'intention (nuire ou non), (c) l'effet produit sur la victime, (d) la conformité aux règles et à la loi. Présentez votre analyse dans un tableau.

\item \textbf{Appliquer la méthode ODARJ au cas de Sandrine.}
\begin{itemize}
\item \textbf{O}bservez et \textbf{D}écrivez la situation du point de vue de Sandrine ;
\item \textbf{A}nalysez : proposez \textbf{deux réactions possibles} et différentes que Sandrine aurait pu adopter ce soir-là (par exemple : répondre publiquement dans le groupe, alerter en privé un administrateur ou un adulte, etc.) ;
\item \textbf{R}éagissez : choisissez, entre ces deux réactions, celle qui vous semble \textbf{la plus citoyenne} ;
\item \textbf{J}ustifiez votre choix en mobilisant l'\cle{empathie} (se mettre à la place d'Amina) et les conséquences de chaque réaction pour la victime, pour le groupe et pour Sandrine elle-même.
\end{itemize}

\item \textbf{Rédiger individuellement un paragraphe TAEC.} En 8 à 12 lignes, chaque élève rédige sa réponse à la question : « Quelle était, selon toi, la conduite la plus citoyenne à adopter face à cette publication, et pourquoi ? » Structure obligatoire : \textbf{T}hèse (ta position), \textbf{A}rgument (la raison principale), \textbf{E}xemple (tiré du dossier), \textbf{C}onclusion (ouverture ou règle générale).

\item \textbf{Rédiger collectivement une charte.} En groupe, rédigez \textbf{5 règles précises} pour la charte du groupe WhatsApp « 3e B – Solidarité ». Chaque règle doit être formulée clairement, interdire ou encourager un comportement identifiable, et prévoir ce qu'il faut faire en cas de problème.

\item \textbf{Le tribunal citoyen.} Chaque groupe désigne un « avocat » qui plaide sa solution en 2 minutes devant la classe. La classe vote ensuite la démarche la mieux justifiée (critères du vote : respect des étapes ODARJ, qualité de la justification, réalisme de la charte).
\end{enumerate}

\courssub{Ressources à mobiliser}

\begin{aretenir}[Les outils de la leçon à réutiliser]
\begin{itemize}
\item \textbf{Définitions} : l'\cle{éthique} est l'ensemble des règles de conduite qui distinguent le bien du mal ; l'\cle{empathie} est la capacité à se mettre à la place d'autrui pour comprendre ce qu'il ressent ; la \cle{citoyenneté numérique} est l'ensemble des comportements responsables et respectueux sur les outils numériques.
\item \textbf{Les 4 critères d'une conduite non éthique} : atteinte à la dignité d'autrui ; intention de nuire ou légèreté coupable ; effet réel sur la victime ; violation des règles (règlement scolaire, conditions d'utilisation de l'application, loi).
\item \textbf{La méthode ODARJ} : Observer $\rightarrow$ Décrire $\rightarrow$ Analyser $\rightarrow$ Réagir $\rightarrow$ Justifier.
\item \textbf{La structure TAEC} : Thèse $\rightarrow$ Argument $\rightarrow$ Exemple $\rightarrow$ Conclusion.
\item \textbf{Rappel de sécurité} : au Cameroun, la loi n° 2010/012 du 21 décembre 2010 relative à la cybersécurité et à la cybercriminalité réprime notamment les atteintes à l'honneur et à la vie privée commises par voie électronique. Relayer une image humiliante peut donc engager la responsabilité de celui qui la transfère.
\end{itemize}
\end{aretenir}

\courssub{Démarche et corrigé commenté}

\begin{exoresolu}[Corrigé commenté de l'affaire n° 2026-03]
\textbf{Énoncé.} Analyser l'affaire du groupe WhatsApp de la 3\textsuperscript{ème} B : identifier les conduites non éthiques, appliquer ODARJ au cas de Sandrine, rédiger un paragraphe TAEC et proposer une charte de 5 règles.

\tcblower

\textbf{Étape 1 — Description des faits.}

Dans un groupe WhatsApp scolaire de 42 membres, Junior publie une photo modifiée et humiliante d'Amina. Moussa transfère l'image dans un autre groupe (12 membres de plus) et Carine la commente avec moquerie. Kevin voit tout et se tait. Sandrine, témoin choqué, hésite à intervenir. En 25 minutes, 12 commentaires moqueurs s'accumulent, et l'image finit par sortir du groupe et à atteindre la victime.

\emph{Commentaire :} cette étape oblige à distinguer les rôles : Amina est la \textbf{victime} ; Junior est l'\textbf{auteur} ; Moussa et Carine sont des \textbf{relais actifs} ; Kevin est un \textbf{témoin passif} ; Sandrine est un \textbf{témoin qui peut devenir acteur citoyen}.

\textbf{Étape 2 — Identification des conduites non éthiques (les 4 critères).}

\begin{center}
\begin{tabularx}{\linewidth}{|l|X|X|}
\hline
\rowcolor{popblueL}
\textbf{Élève} & \textbf{Conduite} & \textbf{Analyse par les 4 critères} \\
\hline
Junior & Modifie et publie la photo & Non éthique : atteinte à la dignité (a), intention de moquer (b), humiliation réelle (c), atteinte au droit à l'image et à l'honneur, réprimée par la loi (d). \\
\hline
Moussa & Transfère l'image dans un autre groupe & Non éthique : il amplifie la diffusion (a, c) ; même sans être l'auteur, relayer une image humiliante engage sa responsabilité (d). \\
\hline
Carine & Commentaires et emojis moqueurs & Non éthique : participation active à la moquerie collective (a, b, c) ; les emojis ne rendent pas la conduite anodine. \\
\hline
Kevin & Voit et se tait & Conduite critiquable : le silence laisse faire et encourage les auteurs (c) ; un témoin passif contribue au sentiment d'isolement de la victime, même s'il n'a pas d'intention de nuire (b). \\
\hline
\end{tabularx}
\end{center}

\emph{Commentaire :} l'erreur fréquente est de croire que « ne rien faire » est neutre. En citoyenneté numérique, le silence du témoin est une demi-faute : il ne punit personne, mais il ne protège personne non plus.

\textbf{Étape 3 — Méthode ODARJ appliquée à Sandrine.}

\begin{itemize}
\item \textbf{Observer} : une image humiliante circule ; la victime est absente et ne peut pas se défendre ; les moqueries s'accumulent (12 commentaires en 25 minutes).
\item \textbf{Décrire} : Sandrine est témoin d'un cyberharcèlement en cours. Elle dispose de plusieurs moyens d'action : le groupe lui-même, les administrateurs, les adultes (parents, vie scolaire, censeur).
\item \textbf{Analyser} — deux réactions possibles :
\begin{itemize}
\item \emph{Réaction 1} : répondre publiquement dans le groupe pour dénoncer la publication et demander son retrait. Avantage : réaction immédiate et visible. Risque : s'exposer aux moqueries, envenimer le débat à 42 membres.
\item \emph{Réaction 2} : écrire en privé à un administrateur pour demander la suppression du message, faire des captures d'écran comme preuves, soutenir Amina en privé et alerter un adulte (parent ou vie scolaire) dès le lendemain. Avantage : action efficace, mesurée, qui protège la victime et conserve les preuves.
\end{itemize}
\item \textbf{Réagir} : la réaction 2 est la plus citoyenne, éventuellement complétée par un message public bref et calme (« Ce n'est pas drôle, supprimez cette photo ») si Sandrine s'en sent la force.
\item \textbf{Justifier} : par empathie, Amina a besoin d'être défendue et soutenue, pas d'un spectacle supplémentaire. La réaction 2 arrête la diffusion (suppression par l'administrateur), préserve les preuves pour la vie scolaire, associe un adulte — ce qui est indispensable face à un harcèlement — et protège aussi Sandrine. Elle transforme un témoin en citoyen numérique.
\end{itemize}

\textbf{Étape 4 — Exemple de paragraphe TAEC (production individuelle attendue).}

\emph{(T)} Face à la publication de la photo modifiée d'Amina, la conduite la plus citoyenne était d'intervenir pour faire supprimer l'image et alerter un adulte, plutôt que de se taire ou de répondre par la moquerie. \emph{(A)} En effet, un citoyen numérique ne se contente pas de ne pas nuire : il protège autrui, car le silence des témoins permet au harcèlement de continuer et isole la victime. \emph{(E)} Dans le dossier, Kevin a vu le message sans réagir : en 25 minutes, 12 commentaires moqueurs se sont accumulés et l'image a quitté le groupe, alors qu'un simple signalement à un administrateur aurait pu l'effacer avant qu'Amina ne la découvre. \emph{(C)} Sur Internet comme dans la cour du collège, la vraie solidarité consiste donc à agir pour celui qui est attaqué : signaler, soutenir la victime et prévenir un adulte sont les trois réflexes du citoyen numérique.

\textbf{Étape 5 — Exemple de charte du groupe (5 règles).}

\begin{enumerate}
\item On ne publie dans le groupe que des contenus utiles à la classe (devoirs, informations, entraide) et respectueux de chacun.
\item Il est interdit de publier, modifier ou transférer la photo d'un camarade sans son accord.
\item Aucune moquerie, insulte ou humiliation n'est tolérée, même « pour rire » ou avec des emojis.
\item Celui qui voit un message blessant ne le transfère pas : il le signale à un administrateur et soutient la personne visée en privé.
\item En cas de problème grave, les administrateurs suppriment le message, conservent une capture comme preuve et en informent un adulte (parent, professeur principal ou vie scolaire).
\end{enumerate}

\emph{Commentaire :} une bonne règle de charte est courte, observable et prévoit une conduite de secours (« que faire si... »). La règle 5 est essentielle : un groupe de mineurs ne se gère jamais sans adulte référent.
\end{exoresolu}

\courssub{Bilan de l'activité}

Cette activité a permis de mettre en évidence plusieurs acquis essentiels de la leçon :

\begin{itemize}
\item \textbf{L'éthique s'applique à des gestes quotidiens.} Une publication, un transfert, un emoji ou un silence dans un groupe WhatsApp ne sont jamais neutres : chacun peut être analysé avec les quatre critères (dignité, intention, effet, règles et loi).
\item \textbf{Le témoin est un acteur décisif.} Dans un cyberharcèlement, la différence entre un drame et un incident maîtrisé tient souvent à la réaction d'un seul témoin : signaler, soutenir, alerter. L'attitude citoyenne vis-à-vis d'autrui commence par le refus de laisser faire.
\item \textbf{La méthode ODARJ structure la réflexion.} Observer, décrire, analyser, réagir, justifier : cette démarche en cinq étapes transforme une émotion (« ce n'est pas bien ») en décision raisonnée et défendable devant les autres.
\item \textbf{La rédaction justifiée est un savoir-faire à part entière.} Le paragraphe TAEC et la plaidoirie du tribunal citoyen ont montré qu'une opinion ne vaut que par la qualité de sa justification : thèse claire, argument solide, exemple précis tiré des faits, conclusion qui dégage une règle générale.
\item \textbf{La prévention passe par des règles communes.} La charte rédigée collectivement illustre qu'une communauté numérique (classe, groupe WhatsApp, réseau social) a besoin de règles explicites et partagées pour rester un espace sûr.
\end{itemize}

\begin{attention}[Erreurs fréquentes à éviter]
\begin{itemize}
\item Confondre « relayer » et « ne rien faire » : transférer une image humiliante est une faute active, pas une simple maladresse.
\item Croire que l'humour excuse tout : une moquerie reste une atteinte à la dignité si la victime la vit comme une humiliation.
\item Rédiger un paragraphe TAEC sans exemple concret tiré du dossier : l'argumentation doit toujours s'appuyer sur les faits.
\item Oublier d'associer un adulte : face au harcèlement, un élève ne doit jamais rester seul, ni laisser la victime seule.
\end{itemize}
\end{attention}

Tu peux maintenant t'auto-évaluer : sais-tu décrire une situation, repérer les conduites non éthiques avec les quatre critères, appliquer ODARJ sans aide, rédiger un paragraphe TAEC complet et proposer des règles de vie numérique ? Si oui, tu es prêt pour l'évaluation de l'unité — et surtout, tu es devenu un meilleur citoyen du numérique.

\newpage

\coursec{Méthodes et exercices résolus}

\courssub{Rappels structurants : éthique, empathie, citoyenneté numérique}

\begin{definition}[Éthique numérique]
Ensemble de principes moraux qui guident le comportement des individus et des organisations dans l'utilisation des technologies de l'information et de la communication (TIC). Elle dépasse le cadre légal : ce qui est \emph{technique}ment possible ou \emph{légal} n'est pas forcément \emph{éthique}.
\end{definition}

\begin{definition}[Empathie numérique]
Capacité à se mettre à la place de l'autre (autrui) dans un environnement dématérialisé pour anticiper l'impact de ses actes (publication, commentaire, partage, code) sur sa dignité, sa vie privée, sa sécurité et son bien-être psychologique. C'est le fondement du \cle{respect} en ligne.
\end{definition}

\begin{definition}[Citoyenneté numérique]
Exercice responsable, critique et solidaire de ses droits et devoirs dans le cyberespace. Elle implique : la connaissance du cadre juridique (lois camerounaises), la protection de soi et des autres (vie privée, sécurité), la lutte contre la désinformation, le respect de la propriété intellectuelle et l'engagement positif (création, entraide, signalement).
\end{definition}

\begin{propriete}[Cadre juridique camerounais de référence (Niveau 3ème)]
\begin{itemize}
    \item \textbf{Loi n°2010/012 du 21 décembre 2010} relative à la cybercriminalité : Sanctionne l'accès frauduleux, l'interception, l'atteinte à l'intimité (Art. 8), le harcèlement électronique (Art. 7), la pédopornographie (Art. 10), l'apologie du terrorisme.
    \item \textbf{Loi n°2024/006} relative à la protection des données à caractère personnel : Consentement explicite, droit d'accès, de rectification, d'effacement, d'opposition, portabilité. Création de l'ANTIC comme autorité de protection.
    \item \textbf{Code Pénal} : Atteinte au secret de la correspondance (Art. 302-1), menaces (Art. 299-300), diffamation, injure publique, pédopornographie (Art. 294 : 10 à 20 ans de prison).
    \item \textbf{Droit à l'image (Art. 9 Code Civil / Jurisprudence)} : Toute personne a droit au respect de sa vie privée ; publication d'image sans consentement = faute civile (réparation) et souvent pénale (cybercriminalité).
\end{itemize}
\end{propriete}

\begin{aretenir}[Les 4 piliers de l'attitude citoyenne en 3ème]
\begin{enumerate}
    \item \textbf{Respect} : Dignité, vie privée, image, opinions, cultures, lois.
    \item \textbf{Responsabilité} : Assumer ses clics (publication, partage, code), réparer ses erreurs, protéger son identité numérique.
    \item \textbf{Légalité} : Connaître les lois (2010/012, 2024/006, Code Pénal), signaler les infractions graves (ANTIC 1510, Police 117, Allô 119).
    \item \textbf{Solidarité / Vigilance} : Ne pas être spectateur passif (complicité), soutenir les victimes, vérifier avant de partager (S.T.O.P.), éduquer son entourage.
\end{enumerate}
\end{aretenir}

\courssub{Identifier les conduites non éthiques dans la vie courante}

\begin{exemplebox}[Typologie des conduites non éthiques observées au Cameroun (Contexte 3ème)]
\begin{center}
\begin{tabularx}{\linewidth}{|l|X|X|}
\hline
\textbf{Catégorie} & \textbf{Conduite non éthique (Exemple concret)} & \textbf{Valeur bafouée / Risque} \\
\hline
\cle{Cyberharcèlement} & Créer un groupe « Anti-Paul » pour se moquer ; diffuser un montage humiliant (deepfake) ; insulter en commentaires. & Dignité, sécurité psychologique, Art. 7 Loi 2010/012. Risque : dépression, phobie scolaire, prison. \\
\hline
\cle{Vie privée / Image} & Filmer une bagarre/accident et poster ; taguer sans accord ; géolocaliser un tiers à son insu ; fouiller le téléphone d'un ami. & Droit à l'image, Art. 8 Loi 2010/012, Loi 2024/006. Risque : plainte, dommages-intérêts, trahison amicale. \\
\hline
\cle{Désinformation (Fake news)} & Partager « BEPC annulé », « Kidnapping à Carrefour X », « Remède miracle Covid/Ebola » sans vérifier (S.T.O.P.). & Vérité, cohésion sociale, sécurité publique. Risque : panique, lynchage, santé publique, Art. 7 (trouble à l'ordre public). \\
\hline
\cle{Propriété intellectuelle} & Copier-coller un devoir/rapport sans citer (plagiat) ; télécharger film/logiciel piraté ; utiliser code IA sans le déclarer. & Honnêteté, droit d'auteur, éthique scolaire/pro. Risque : 0 au devoir, exclusion, virus, perte de crédibilité. \\
\hline
\cle{Arnaque / Broutage} & Se faire passer pour un ami en détresse (WhatsApp piraté) ; vendre faux billets/visas ; phishing (faux lien MINSEC/Banque). & Confiance, légalité (Art. 311 Code Pénal escroquerie), Art. 7 Loi 2010/012. Risque : prison, ruine victimes. \\
\hline
\cle{Contenus illicites graves} & Détention/diffusion pédopornographie ; apologie terrorisme ; incitation haine tribale/religieuse. & Humanité, sécurité de l'État. \textbf{CRIMES} (Art. 294 Code Pénal, Loi antiterroriste). Tolérance zéro. \\
\hline
\cle{E-réputation négative} & Profil public : photos alcool/tabac (mineur), insultes, fautes massives, geoloc domicile, conflits publics. & Avenir professionnel (80\% recruteurs googolent), honneur familial. Risque : refus stage/bourse/emploi. \\
\hline
\cle{Conflit non géré} & Exclure d'un groupe sans explication ; vengeance (piratage, diffusion secrets) ; escalade publique au lieu dialogue privé. & Paix, dialogue, réparation. Risque : rupture définitive, sanctions scolaires, judiciarisation. \\
\hline
\cle{Écologie numérique} & Laisser 5000 mails non lus ; streaming 4K en 4G ; changer tel chaque an ; ne jamais vider cloud/corbeille. & Responsabilité planétaire (Green IT), sobriété. Risque : empreinte carbone, saturation stockage, gaspillage. \\
\hline
\cle{Usage non éthique IA} & Rendre devoir 100\% ChatGPT sans lecture ; générer deepfake nuisible ; croire aveuglément (hallucinations) ; tricher au concours. & Intégrité intellectuelle, esprit critique, vérité. Risque : incompétence réelle, sanction, propagation erreurs. \\
\hline
\end{tabularx}
\end{center}
\end{exemplebox}

\courssub{Adopter une attitude citoyenne vis-à-vis d'autrui}

\begin{methode}[Réflexe « 3 Questions » avant tout clic (Publication, Partage, Envoi, Code)]
\textbf{Objectif :} Ancrer l'empathie préventive comme automatisme (System 1 : rapide, intuitif).
\begin{enumerate}
    \item \textbf{Est-ce VRAI ?} (Vérification S.T.O.P. : Source, Temps, Origine, Recoupement). Si doute $\rightarrow$ Ne pas partager.
    \item \textbf{Est-ce UTILE / BIENVEILLANT ?} Apporte-t-il de la valeur (info, soutien, joie partagée) ou du bruit/du mal (rage, moquerie, rumeur, plainte stérile) ?
    \item \textbf{Est-ce RESPECTUEUX / LÉGAL ?} Consentement (photo/data) ? Droit d'auteur ? Pas d'insulte/menace/haine ? Conforme loi 2010/012 / 2024/006 ?
\end{enumerate}
\textbf{Règle d'or :} Si \textbf{un seul NON} $\rightarrow$ \textbf{On ne clique pas.} On respire 10 secondes. On reformule, on floute, on met en privé, on supprime le brouillon, on signale.
\end{methode}

\begin{propriete}[Principes directeurs de l'action citoyenne (Subsidiarité, Proportionnalité, Réparation)]
\begin{itemize}
    \item \textbf{Subsidiarité :} On règle au plus bas niveau possible. 1. Dialogue direct (OSBD) / Suppression / Signalement plateforme. 2. Médiation adulte (prof, parent, pair formé) / Sanction éducative. 3. Signalement autorités (ANTIC, Police, Justice) \emph{uniquement si infraction grave ou échec niveau 1-2}.
    \item \textbf{Proportionnalité :} La réponse est adaptée à la gravité. Une blague de mauvais goût $\neq$ pédopornographie. Exclusion définitive = dernier recours (récidive, danger).
    \item \textbf{Réparation :} Privilégier la restauration du lien (excuses, suppression, service, médiation) à la seule punition. L'auteur doit \emph{comprendre} le mal fait (empathie cognitive).
\end{itemize}
\end{propriete}

\courssub{La démarche d'analyse d'une situation d'éthique (Méthode ODARJ)}

\begin{methode}[Analyser une situation éthique avec la démarche ODARJ]
\textbf{Objectif :} Structurer son raisonnement pour juger une situation numérique et proposer une action citoyenne.
\begin{enumerate}
    \item \textbf{O -- Observer :} Décrire les faits \emph{objectifs} (qui, quoi, où, quand, comment) sans interprétation ni jugement.
    \item \textbf{D -- Discerner :} Identifier les valeurs en conflit (vie privée, propriété intellectuelle, respect, sécurité, vérité) et les personnes impactées (victimes, auteurs, témoins).
    \item \textbf{A -- Analyser :} Évaluer les conséquences à court et long terme (juridiques, psychologiques, sociales, techniques). Se demander : « Est-ce légal ? Est-ce respectueux ? Est-ce que j'aimerais que l'on me le fasse ? ».
    \item \textbf{R -- Résoudre :} Choisir la meilleure option parmi : signaler, supprimer, ne pas partager, dialoguer, alerter un adulte de confiance, consulter la charte d'usage.
    \item \textbf{J -- Justifier :} Rédiger une réponse argumentée en reliant la décision choisie aux valeurs citoyennes (empathie, responsabilité, légalité) et aux conséquences évitées.
\end{enumerate}
\emph{Réflexe :} Toujours distinguer le \textbf{fait} (ce qui s'est passé) de l'\textbf{opinion} (ce que j'en pense).
\end{methode}

\begin{exoresolu}[Partage non consenti d'une photo dans un groupe WhatsApp de la classe]
\textbf{Énoncé :} Lors d'une sortie scolaire au Musée National de Yaoundé, Kevin prend en photo sa camarade Amina qui a trébuché et a une expression ridicule. Sans lui demander son avis, il poste la photo sur le groupe WhatsApp « 3ème A 2026 » avec la légende « La chute de l'année [emoji rire] ». La photo fait le tour de l'établissement en une heure. Amina est humiliée, elle ne veut plus venir en cours. Le professeur principal saisit l'administration. En utilisant la méthode ODARJ, analysez la situation et proposez une résolution argumentée.

\tcblower
\textbf{Solution détaillée selon la méthode ODARJ :}

\textbf{1. O -- Observer (Les faits bruts) :}
\begin{itemize}
    \item \textbf{Acteurs :} Kevin (auteur), Amina (victime), élèves du groupe (témoins/relais), administration (autorité).
    \item \textbf{Action :} Prise de photo dans un moment de vulnérabilité (chute) + publication sur un groupe privé mais large (classe/établissement) sans consentement + légende moqueuse.
    \item \textbf{Contexte :} Sortie scolaire encadrée, espace semi-public, mineurs de 14-15 ans.
    \item \textbf{Conséquence immédiate :} Viralité intra-établissement, humiliation publique d'Amina, refus de scolarité (absentéisme).
\end{itemize}

\textbf{2. D -- Discerner (Valeurs et personnes impactées) :}
\begin{itemize}
    \item \textbf{Valeurs bafouées :}
    \begin{itemize}
        \item \cle{Respect de la dignité humaine} et de l'image (Droit à l'image, Art. 9 Code Civil camerounais / Loi 2010/012 sur la cybercriminalité).
        \item \cle{Vie privée} : moment de faiblesse exposé sans consentement.
        \item \cle{Empathie} : absence totale de mise à la place d'autrui (Kevin rit, Amina souffre).
        \item \cle{Responsabilité numérique} : tout partage engage son auteur.
    \end{itemize}
    \item \textbf{Personnes impactées :}
    \begin{itemize}
        \item \textbf{Amina :} Traumatisme psychologique, perte de confiance, décrochage scolaire (risque grave).
        \item \textbf{Kevin :} Risque disciplinaire (conseil de discipline, exclusion), risque pénal (harcèlement, atteinte à l'image), réputation ternie.
        \item \textbf{Témoins (élèves) :} Banalisation de la moquerie, effet de meute, complicité passive par like/partage.
        \item \textbf{École :} Climat scolaire dégradé, obligation de protection des mineurs.
    \end{itemize}
\end{itemize}

\textbf{3. A -- Analyser (Conséquences et légalité) :}
\begin{itemize}
    \item \textbf{Juridique :} Au Cameroun, la Loi n°2010/012 du 21 décembre 2010 relative à la cybercriminalité (Art. 7 et suivants) sanctionne le harcèlement par voie électronique et l'atteinte à l'intimité/vie privée. La publication d'image sans consentement est une infraction. Kevin (mineur > 13 ans) engage sa responsabilité pénale.
    \item \textbf{Psychologique :} Cyberharcèlement avéré (répétition via partages, déséquilibre de pouvoir, intention de nuire ou imprudence caractérisée). Risque de dépression, phobie scolaire.
    \item \textbf{Social :} Destruction du lien de confiance dans la classe. Normalisation de la violence symbolique.
    \item \textbf{Test de l'empathie :} Kevin aurait-il aimé qu'on filme sa propre chute et qu'on la diffuse ? Non. \textbf{Règle d'or :} Ne pas faire à autrui ce qu'on ne veut pas qu'on nous fasse.
\end{itemize}

\textbf{4. R -- Résoudre (Actions concrètes et graduées) :}
\begin{enumerate}
    \item \textbf{Immédiat (Kevin) :} Supprimer le message \emph{partout} (groupe, statut, téléphone), présenter des excuses sincères \emph{en privé} à Amina (pas en public pour ne pas raviver l'humiliation).
    \item \textbf{Immédiat (Témoins) :} Cesser tout partage, supprimer la photo de leurs appareils, ne pas commenter. Signaler le contenu à l'admin du groupe / professeur.
    \item \textbf{Institutionnel (Administration/Professeur) :} Retrait de la photo si encore visible. Médiation éducative (pas seulement punitive) : Kevin doit comprendre le mal fait. Sanction éducative (ex: exposé sur le droit à l'image, service à la communauté scolaire). Signalement au procureur si gravité avérée (procédure mineure).
    \item \textbf{Structurel (Classe) :} Séance de vie de classe sur le cyberharcèlement, charte numérique de la classe signée par tous.
\end{enumerate}

\textbf{5. J -- Justifier (Réponse argumentée type) :}
\begin{quote}
\textit{« Le partage de la photo d'Amina sans son consentement constitue une violation grave de son droit à l'image et de sa dignité, aggravée par la visée moqueuse (cyberharcèlement). Au regard de la loi camerounaise (Loi 2010/012) et des valeurs citoyennes d'empathie et de responsabilité, Kevin doit assumer les conséquences de son acte (sanction disciplinaire, excuses réparatrices). La classe doit collectivement refuser la viralité de l'humiliation en supprimant le fichier et en soutenant la victime. L'école a le devoir de protéger Amina (droit à l'éducation dans un climat serein) et d'éduquer Kevin à la citoyenneté numérique plutôt que de l'exclure définitivement, sauf récidive. »}
\end{quote}

\textbf{Conclusion :} L'éthique numérique n'est pas optionnelle : un clic peut briser une scolarité. L'empathie impose de \textbf{réfléchir avant de publier} (règle des 3 questions : Est-ce vrai ? Est-ce utile ? Est-ce bienveillant ?).
\end{exoresolu}

\courssub{Rédiger et justifier une réponse argumentée}

\begin{methode}[Rédiger une charte de bonne conduite numérique (Classe / Cybercafé / Famille)]
\textbf{Objectif :} Transformer des principes éthiques abstraits en règles concrètes, partagées et applicables.
\begin{enumerate}
    \item \textbf{Concertation :} Réunir les usagers (élèves, gérant, parents). Chacun propose 2 règles indispensables.
    \item \textbf{Regroupement :} Fusionner les doublons. Classer par thèmes : \textbf{Respect (Empathie)}, \textbf{Sécurité (Protection)}, \textbf{Légalité (Droit)}, \textbf{Écologie (Usage raisonné)}.
    \item \textbf{Formulation positive :} Écrire des \textbf{« Je m'engage à... »} plutôt que des « Interdiction de... ». Ex: « Je demande l'autorisation avant de poster une photo d'autrui » au lieu de « Interdit de poster sans accord ».
    \item \textbf{Sanctions graduées et réparatrices :} Prévoir 3 niveaux : 1. Rappel à la règle (oral) -- 2. Mesure éducative (exposé, aide à la victime, suspension d'accès temporaire) -- 3. Sanction institutionnelle (conseil de discipline, plainte).
    \item \textbf{Signature et affichage :} Chaque usager signe. Affichage visible (mur classe, écran d'accueil cybercafé). Révision annuelle.
\end{enumerate}
\emph{Reflexe :} Une charte non signée et non affichée n'existe pas. La \cle{co-construction} garantit l'adhésion (savoir-être : appropriation).
\end{methode}

\begin{exoresolu}[Rédaction d'une clause de charte pour un cybercafé scolaire à Douala]
\textbf{Énoncé :} Le gérant du cybercafé « École Numérique » au quartier Bonamoussadi souhaite ajouter une clause « Respect de la vie privée et lutte contre le harcèlement » à sa charte d'usage affichée sur chaque poste. Il veut une formulation claire, positive, juridiquement solide (référence loi camerounaise) et avec une sanction précise. Rédigez cette clause complète.

\tcblower
\textbf{Clause proposée pour la Charte d'Usage -- Cybercafé « École Numérique »}

\begin{center}
\begin{tabularx}{\linewidth}{|>{\bfseries}l|X|}
\hline
\textbf{Thème} & \textbf{Respect de la vie privée, de l'image et lutte contre le cyberharcèlement} \\
\hline
\textbf{Engagement (Formulation positive)} & \textbf{« Je m'engage à :} \\
& \quad 1. Ne jamais photographier, filmer, enregistrer ou diffuser l'image, la voix ou les données personnelles d'un tiers sans son consentement libre, éclairé et explicite. \\
& \quad 2. Ne jamais créer, transmettre ou relayer des contenus moqueurs, humiliants, menaçants ou haineux à l'encontre d'une personne (messages, montages, \emph{deepfakes}, rumeurs). \\
& \quad 3. Signaler immédiatement au gérant tout contenu harcelant ou illicite dont j'ai connaissance sur les postes du cybercafé. \\
& \quad 4. Respecter la vie privée des autres usagers : ne pas regarder leur écran, ne pas accéder à leurs sessions ouvertes, me déconnecter de mes comptes à mon départ. » \\
\hline
\textbf{Fondement légal (Cameroun)} & Loi n°2010/012 du 21 décembre 2010 relative à la cybercriminalité (Art. 7 : harcèlement électronique ; Art. 8 : atteinte à l'intimité ; Art. 9 : interception). Code Pénal (Art. 302-1 : violation du secret de la correspondance). Loi n°2024/006 sur la protection des données à caractère personnel. \\
\hline
\textbf{Sanctions graduées appliquées par le gérant} & \textbf{Niveau 1 (Imprudence) :} Avertissement oral + Rappel de la clause + Suppression du contenu sous contrôle. \\
& \textbf{Niveau 2 (Récidive / Gravité modérée) :} Exclusion temporaire du cybercafé (7 jours) + Information aux parents/tuteurs si mineur + Travail de réflexion écrit (1 page) sur « Les conséquences du harcèlement ». \\
& \textbf{Niveau 3 (Infraction caractérisée / Danger) :} Exclusion définitive + Signalement aux autorités compétentes (Police judiciaire, ANTIC) + Transmission des logs de connexion (adresse MAC, heure, poste) conformément à la réquisition judiciaire. \\
\hline
\textbf{Signature usager / Date} & \hrulefill \quad / \quad \hrulefill \\
\hline
\end{tabularx}
\end{center}

\textbf{Justification pédagogique :}
\begin{itemize}
    \item \textbf{Clarté :} Verbes d'action (photographier, diffuser, signaler), pas de jargon.
    \item \textbf{Empathie :} Le point 3 transforme le témoin en acteur protecteur (savoir-être citoyen).
    \item \textbf{Ancrage légal :} Citation précise des articles de la loi 2010/012 et de la nouvelle loi 2024 sur les données personnelles : montre que ce n'est pas une « règle de la maison » mais la \textbf{Loi de la République}.
    \item \textbf{Sanction réparatrice (Niv 2) :} L'écrit de réflexion force l'empathie cognitive (se mettre à la place de la victime).
    \item \textbf{Preuve technique (Niv 3) :} Mention des logs (traçabilité) : dissuasion technique réelle.
\end{itemize}
\end{exoresolu}

\begin{methode}[Réagir face à une fake news / rumeur virale (Réflexe « S.T.O.P. »)]
\textbf{Objectif :} Ne pas propager la désinformation et protéger la cohésion sociale (empathie collective).
\begin{enumerate}
    \item \textbf{S -- Source :} Qui a émis l'info ? Compte officiel (Gouvernement, École, Média reconnu) ou compte anonyme / « ami d'un ami » / capture d'écran sans lien ?
    \item \textbf{T -- Temps :} L'info est-elle datée ? Une vieille info (ex: grève 2021) ressortie aujourd'hui crée la panique. Vérifier la date de publication originale.
    \item \textbf{O -- Origine / Recoupement :} L'info est-elle reprise par \textbf{au moins 2 sources fiables indépendantes} (ex: CRTV + Cameroon Tribune + Site officiel Minesec) ? Attention aux sites « miroirs » qui copient-collent la même erreur.
    \item \textbf{P -- Partager / Pas partager :}
    \begin{itemize}
        \item \textbf{Vérifié VRAI + Utile :} Je partage en citant la source officielle.
        \item \textbf{Vérifié FAUX :} Je \textbf{ne partage pas}. Je signale (bouton « Signaler » sur WhatsApp/FB/X). J'informe poliment l'expéditeur : « Attention, cette info est fausse, voici la source officielle : [lien] ».
        \item \textbf{Incertain :} Je \textbf{ne partage pas}. J'attends la confirmation officielle. « Mieux vaut perdre une info que propager une rumeur. »
    \end{itemize}
\end{enumerate}
\emph{Règle d'or :} L'empathie numérique, c'est \textbf{protéger les autres de l'angoisse inutile} (ex: fausse rumeur de kidnapping, de changement de programme d'examen, de décès d'une personnalité).
\end{methode}

\begin{exoresolu}[Gestion d'une rumeur d'annulation du BEPC sur les réseaux sociaux]
\textbf{Énoncé :} Vendredi 20 juin, 14h30. Une capture d'écran circule massivement sur les groupes WhatsApp « Parents 3ème » et « Élèves 3ème » : « COMMUNIQUÉ OFFICIEL MINSEC -- BEPC 2026 ANNULÉ SUITE AUX FUITES. NOUVELLES DATES EN SEPTEMBRE. » L'image porte un en-tête imitant le logo du Ministère, mais pas de numéro de communiqué, pas de signature du Ministre, pas de lien vers le site \texttt{minesec.gov.cm}. La panique s'installe. En tant que délégué de classe responsable, vous devez rédiger un message officiel pour calmer vos camarades et les parents. Appliquez la méthode S.T.O.P.

\tcblower
\textbf{Application rigoureuse de la méthode S.T.O.P. :}

\textbf{1. S -- Source :} \textbf{Inconnue / Non officielle.} Capture d'écran (image), pas de lien cliquable vers le site source. Compte émetteur : « Transféré plusieurs fois » (double flèche WhatsApp). Aucun compte officiel (\texttt{@Minesec\_Cameroun} sur X, page Facebook « Ministère des Enseignements Secondaires ») n'a relayé.

\textbf{2. T -- Temps :} \textbf{Incohérence temporelle.} Le BEPC écrit démarre généralement mi-juin. Une annulation « suite à des fuites » à J-2 ou pendant les épreuves est une décision majeure qui aurait été annoncée en conférence de presse télévisée (CRTV) jours avant, pas par une image floue un vendredi après-midi.

\textbf{3. O -- Origine / Recoupement :}
\begin{itemize}
    \item Site \texttt{minesec.gov.cm} : Aucune actualité ce jour.
    \item Site \texttt{crtv.cm} / \texttt{cameroon-tribune.cm} : Silence total.
    \item Comptes officiels vérifiés (badge bleu) : Aucun communiqué.
    \item \textbf{Conclusion :} Zéro source fiable. C'est une \textbf{fake news} (infox) classique de fin d'année.
\end{itemize}

\textbf{4. P -- Partager / Pas partager :} \textbf{Décision : NE PAS PARTAGER. SIGNALER. INFORMER.}

\textbf{Message type à diffuser par le délégué (Modèle rigoureux) :}

\begin{lstlisting}[language=text]
[ALERTE] FAKE NEWS / INFOX -- CLASSE 3EME A -- 20/06/2026 14:45

Chers camarades, chers parents,

Une capture d'écran annonçant l'ANNULATION DU BEPC 2026 circule actuellement.
[FAUX] C'EST FAUX.

[OK] VÉRIFICATION FAITE (Méthode S.T.O.P.) :
* Source : Image anonyme, sans numéro, sans signature, sans lien officiel.
* Temps : Incohérent (en plein épreuves).
* Recoupement : AUCUNE trace sur minesec.gov.cm, CRTV, Cameroon Tribune, comptes officiels MINSEC.

[BOUCLIER] CONSIGNES :
1. NE TRANSFÉREZ PAS ce message (vous participez à la panique).
2. SIGNALEZ le message sur WhatsApp (tenir appuyé > Signaler).
3. CONCENTREZ-VOUS sur vos révisions / repos. Les épreuves se poursuivent normalement.
4. SEULES sources valides : Site MINSEC, CRTV, Communiqués signés du Ministre.

Restons calmes, solidaires et lucides. L'empathie, c'est ne pas angoisser les autres avec des rumeurs.

[Votre Nom] -- Délégué 3ème A -- Tél: 6XX XX XX XX (pour vérif)
\end{lstlisting}

\textbf{Justification de la rédaction :}
\begin{itemize}
    \item \textbf{Ton :} Calme, autoritaire (délégué), bienveillant (empathie pour les angoissés).
    \item \textbf{Structure :} Alerte visuelle ([ALERTE]), Verdict immédiat ([FAUX] C'EST FAUX), Preuve méthodique (S.T.O.P.), Action concrète (1.2,3.4), Rappel valeurs (solidarité, lucidité).
    \item \textbf{Traçabilité :} Signature + téléphone = responsabilité assumée.
\end{itemize}
\end{exoresolu}

\begin{methode}[Gérer un conflit numérique entre pairs (Médiation par l'empathie -- OSBD / CNV)]
\textbf{Objectif :} Résoudre un différend (insultes, malentendu, exclusion de groupe) sans violence ni escalade, en restaurant le lien.
\begin{enumerate}
    \item \textbf{Cadre sécurisé :} Proposer un échange \textbf{en privé} (pas sur le groupe public), de vive voix si possible (appel vocal/visio) ou messages vocaux (ton = 90\% du sens). Éviter le texte froid qui amplifie les malentendus.
    \item \textbf{Écoute active (Technique OSBD - CNV) :}
    \begin{itemize}
        \item \textbf{O -- Observation :} « J'ai vu/reçu ce message : "..." » (Fait pur, sans « tu es méchant »).
        \item \textbf{S -- Sentiment :} « Moi, je me sens blessé / inquiet / en colère... » (Responsabilité de son émotion).
        \item \textbf{B -- Besoin :} « J'ai besoin de respect / de clarté / de sécurité... » (Besoin universel).
        \item \textbf{D -- Demande :} « Est-ce que tu acceptes de... (supprimer / t'excuser / m'expliquer en privé) ? » (Action concrète, négociable, positive).
    \end{itemize}
    \item \textbf{Reformulation :} L'autre répète ce qu'il a compris : « Si je comprends bien, tu te sens X parce que tu as besoin de Y, et tu me demandes Z. » Valide l'écoute.
    \item \textbf{Recherche de solution gagnant-gagnant :} « Qu'est-ce qui nous permet à tous les deux de sortir grandis ? » (Excuses publiques ? Clarification ? Règle de groupe modifiée ? Médiation adulte ?).
    \item \textbf{Clôture :} Accord oral/écrit. Suppression des contenus blessants. Engagement de non-récidive. Merci mutuel.
\end{enumerate}
\emph{Attitude :} Ne pas juger l'intention (« Tu as fait exprès »), juger l'\textbf{impact} (« Ce message a eu pour effet de m'exclure »).
\end{methode}

\begin{exoresolu}[Conflit : Exclusion d'un élève d'un groupe de travail WhatsApp]
\textbf{Énoncé :} Pour un devoir de groupe (Exposé Histoire-Géo), le chef de groupe « LeaderMax » retire « TimidePaul » du groupe WhatsApp « Expo 3ème B » sans explication, en écrivant : « Paul t'es trop lent, on avance pas. Dégage. » Paul, humilié, ne dit rien mais ne rend pas sa partie. Le professeur note l'inachèvement. En tant que médiateur pair (élève formé), vous intervenez. Montrez le dialogue type OSBD entre Paul et LeaderMax.

\tcblower
\textbf{Mise en situation :} Médiation en visio (15 min), animée par le médiateur pair. Règle : On se coupe pas, on parle en « Je ».

\textbf{Étape 1 : Paul s'exprime (OSBD)}
\begin{quote}
\textbf{Paul :} « \textbf{O -- Observation :} Hier à 20h, tu as écrit "Paul t'es trop lent, on avance pas. Dégage" et tu m'as retiré du groupe "Expo 3ème B". \textbf{S -- Sentiment :} Moi, je me suis senti humilié, injustement traité et angoissé pour ma note. \textbf{B -- Besoin :} J'ai besoin de respect dans le travail d'équipe, j'ai besoin qu'on me donne une chance de m'expliquer si je suis en retard, j'ai besoin de sécurité pour bien travailler. \textbf{D -- Demande :} Est-ce que tu acceptes de me réintégrer le groupe, de t'excuser pour le mot "Dégage", et qu'on fixe ensemble une deadline claire pour ma partie ? »
\end{quote}

\textbf{Étape 2 : LeaderMax reformule (Validation de l'écoute)}
\begin{quote}
\textbf{LeaderMax :} « Si je comprends bien, tu te sens humilié et angoissé parce que je t'ai viré sans te parler, avec un mot violent. Tu as besoin de respect et de clarté sur les délais. Tu me demandes de te remettre dans le groupe, de m'excuser pour "Dégage" et de fixer une date butoir. C'est ça ? »
\end{quote}
\textbf{Paul :} « Oui, c'est exactement ça. »

\textbf{Étape 3 : LeaderMax s'exprime (OSBD)}
\begin{quote}
\textbf{LeaderMax :} « \textbf{O -- Observation :} Paul, tu n'as pas envoyé ta partie (introduction) mercredi soir comme convenu oralement lundi. Je t'ai relancé mardi, pas de réponse. \textbf{S -- Sentiment :} J'ai eu peur qu'on rate la date de rendu (vendredi), je me suis senti seul à porter le groupe, j'ai paniqué. \textbf{B -- Besoin :} J'ai besoin de fiabilité, de communication (répondre aux messages), de réussite collective. \textbf{D -- Demande :} Est-ce que tu t'engages à me répondre sous 2h max sur WhatsApp pour le reste du projet, et à m'envoyer ton intro ce soir 22h max ? »
\end{quote}

\textbf{Étape 4 : Paul reformule}
\begin{quote}
\textbf{Paul :} « Tu as eu peur pour le délai, tu te sentais seul, tu as besoin de fiabilité et de réponses rapides. Tu me demandes de répondre sous 2h et d'envoyer l'intro ce soir 22h. C'est noté, j'accepte. »
\end{quote}

\textbf{Étape 5 : Accord gagnant-gagnant (Écrit sur papier / note partagée)}
\begin{center}
\begin{tabularx}{\linewidth}{|X|}
\hline
\textbf{ACCORD DE MÉDIATION -- EXPO HISTOIRE 3ÈME B -- 15/05/2026} \\
\hline
1. LeaderMax réintègre Paul dans le groupe WhatsApp \textbf{immédiatement}. \\
2. LeaderMax présente ses excuses pour le terme "Dégage" \textbf{dans le groupe} (message : « Paul, désolé pour mon agressivité hier. On bosse ensemble. »). \\
3. Paul envoie sa partie (Introduction) \textbf{ce soir avant 22h00}. \\
4. Communication : Paul répond aux relances sous \textbf{2h max} (sinon appel vocal). LeaderMax ne prend plus de décision unilatérale d'exclusion sans en parler d'abord. \\
5. Si blocage : On appelle le prof mardi 8h (médiation adulte). \\
\hline
\textbf{Signatures :} \hrulefill (Paul) \quad \hrulefill (LeaderMax) \quad \hrulefill (Médiateur) \\
\hline
\end{tabularx}
\end{center}

\textbf{Analyse :} L'empathie a permis de passer du \textbf{jugement moral} (« Tu es méchant / Tu es nul ») à la \textbf{révélation des besoins} (Sécurité vs Fiabilité). La sanction (exclusion) est remplacée par la \cle{réparation} (excuses + engagement concret). Le professeur n'est saisi qu'en dernier recours (subsidiarité).
\end{exoresolu}

\begin{methode}[Évaluer l'empreinte numérique avant de publier (Test du « Grand-Père / Recruteur / Juge »)]
\textbf{Objectif :} Anticiper la pérennité et la viralité d'une publication (texte, photo, commentaire, like) pour protéger sa réputation future (e-réputation).
\begin{enumerate}
    \item \textbf{Projection temporelle :} Imaginez cette publication dans \textbf{5 ans, 10 ans, 20 ans}. Elle est \textbf{indexée par Google}, archivée par \texttt{archive.org}, stockée sur les serveurs de Meta/Google/TikTok (hors de votre contrôle).
    \item \textbf{Test des 3 regards (Filtres) :}
    \begin{itemize}
        \item \textbf{[GRAND-PERE] Regard du Grand-Père (Valeurs morales / Famille) :} « Est-ce que je suis fier que ma grand-mère / mon imam / mon pasteur / mes futurs enfants voient ça ? » (Respect, pudeur, honneur).
        \item \textbf{[RECRUTEUR] Regard du Recruteur (Employabilité) :} « Est-ce que ça prouve mon professionnalisme, ma fiabilité, mon esprit d'équipe ? Ou ça montre de l'irrespect, de l'alcool, des propos haineux, des fautes d'orthographe massives ? » (80\% des recruteurs googolent les candidats au Cameroun).
        \item \textbf{[JUGE] Regard du Juge (Légalité / Preuve) :} « Ce contenu peut-il être utilisé \textbf{contre moi} dans une procédure (harcèlement, diffamation, apologie du terrorisme, atteinte à la sûreté de l'État, vol de propriété intellectuelle) ? » (Loi 2010/012, Code Pénal, Loi antiterroriste).
    \end{itemize}
    \item \textbf{Décision :}
    \begin{itemize}
        \item \textbf{3 feux verts :} Publier sereinement.
        \item \textbf{1 feu orange :} Modifier (flouter visage, reformuler, mettre en privé « Amis proches »).
        \item \textbf{1 feu rouge :} \textbf{NE PAS PUBLIER.} Supprimer le brouillon.
    \end{itemize}
    \item \textbf{Nettoyage régulier :} Tous les 6 mois : Google son nom + prénom + ville. Supprimer les vieux comptes inutilisés (vieux Skyblog, ancien FB). Paramétrer la confidentialité (Profil privé, « Qui peut me taguer ? » = « Moi seulement »).
\end{enumerate}
\emph{Mantra :} « Internet n'oublie jamais. Le clic est plus rapide que la réflexion. \textbf{Respirez 10 secondes avant d'appuyer sur Entrée / Envoyer / Publier.} »
\end{methode}

\begin{exoresolu}[Cas pratique : « Je veux poster ma photo en tenue de fête (alcool visible) pour mon anniversaire »]
\textbf{Énoncé :} Stéphanie (3ème, 15 ans) veut poster sur son Instagram public (profil non verrouillé) une photo de son anniversaire : elle tient une bouteille de bière (marque visible), fait un doigt d'honneur « pour rigoler », légende : « 15 ans, la vie est belle, les haters peuvent crever [emoji rire][emoji bière] ». Elle vise 500 likes. Appliquez le test des 3 regards et donnez la décision finale avec modifications concrètes si feu orange.

\tcblower
\textbf{Application du Test des 3 Regards :}

\textbf{1. [GRAND-PERE] Regard du Grand-Père (Valeurs / Famille / Culture) :}
\begin{itemize}
    \item \textbf{Alcool :} Mineure (15 ans) + alcool visible = Transgression forte (Loi camerounaise : vente interdite < 18 ans, consommation déconseillée). Image « débauche » pour aînés.
    \item \textbf{Doigt d'honneur + Insulte (« haters peuvent crever ») :} Violence symbolique, irrespect total. Inacceptable dans la culture du respect des aînés (valeur cardinale au Cameroun).
    \item \textbf{Verdict :} \textbf{[ROUGE] FEU ROUGE.}
\end{itemize}

\textbf{2. [RECRUTEUR] Regard du Recruteur (Employabilité future -- dans 3-5 ans, BTS/Université/Stage) :}
\begin{itemize}
    \item \textbf{Profil public :} Accessible à tous. Première image = alcool + insultes + provocation.
    \item \textbf{Interprétation :} « Manque de maturité », « Problème d'alcool potentiel », « Gestion des conflits par l'agressivité », « Orthographe approximative (haters) », « Mauvaise e-réputation ».
    \item \textbf{Risque :} Refus de stage, refus de bourse, refus d'embauche. Le recruteur ne demandera pas d'explication, il passera au CV suivant.
    \item \textbf{Verdict :} \textbf{[ROUGE] FEU ROUGE.}
\end{itemize}

\textbf{3. [JUGE] Regard du Juge (Légalité / Preuve numérique) :}
\begin{itemize}
    \item \textbf{Incitation à la consommation d'alcool par mineur :} Photo = preuve visuelle. Risque pour l'établissement (vente/consommation en milieu scolaire ?) et pour les parents (abandon moral ?).
    \item \textbf{Menace / Insulte publique :} « Les haters peuvent crever » = Menace de mort / Violence verbale (Code Pénal Art. 299-300, Loi 2010/012 Art. 7 harcèlement). Capture d'écran = preuve irréfutable.
    \item \textbf{Droit à l'image :} Si d'autres mineurs sont sur la photo sans floutage = infraction (Loi 2024/006 données personnelles mineurs).
    \item \textbf{Verdict :} \textbf{[ROUGE] FEU ROUGE (Risque pénal réel).}
\end{itemize}

\textbf{Décision finale :} \textbf{INTERDICTION FORMELLE DE PUBLIER CETTE PHOTO TELLE QUELLE.}

\textbf{Proposition alternative (Si Stéphanie TIENT à marquer l'événement) -- Transformation en FEU VERT :}
\begin{enumerate}
    \item \textbf{Photo :} Choisir une photo \textbf{sans alcool visible} (bouteille cachée/jetée), \textbf{sans geste obscène}, sourire naturel avec amis/gâteau.
    \item \textbf{Légende :} « 15 ans aujourd'hui ! Merci la vie, merci mes parents, mes amis, mes profs pour le chemin parcouru. Hâte de construire la suite [GATEAU][PRIERE][CM] \#Anniversaire \#Gratitude \#BEPC2026 ».
    \item \textbf{Confidentialité :} \textbf{Profil PRIVÉ} (paramètre : « Abonnés approuvés uniquement »). \textbf{Taguer} uniquement les personnes sur la photo (après accord oral).
    \item \textbf{Vérification :} Relire 3 fois. Demander à un adulte de confiance (grand frère, prof) : « Tu trouves ça correct pour mon futur ? »
\end{enumerate}

\textbf{Conclusion pédagogique :} L'empathie envers \textbf{son futur soi} (le soi de 20 ans qui cherche un stage) est la forme la plus élevée de responsabilité numérique. Publier l'original = s'auto-saboter. Publier l'alternative = construire une \cle{e-réputation positive}.
\end{exoresolu}

\begin{methode}[Signaler un contenu illicite ou dangereux (Procédure officielle Cameroun -- ANTIC / Police Judiciaire)]
\textbf{Objectif :} Savoir agir concrètement face à un contenu pédopornographique, terroriste, harcelant, arnaque (brouteur), ou usurpation d'identité. Ne pas rester spectateur passif.
\begin{enumerate}
    \item \textbf{Ne pas interagir :} Ne pas liker, commenter, partager, télécharger (surtout pédopornographie = détention d'images pédopornographiques = CRIME, Art. 294 Code Pénal / Loi 2010/012). Ne pas répondre aux brouteurs (arnaqueurs).
    \item \textbf{Preuve (Capture d'écran horodatée) :} \texttt{Ctrl+Maj+S} (Windows) / \texttt{Cmd+Maj+4} (Mac) / Power+Vol- (Android) / Power+Home (iOS). \textbf{Doit apparaître :} L'URL complète (barre d'adresse), le pseudo/numéro de l'auteur, le contenu litigieux, la DATE et l'HEURE système visibles.
    \item \textbf{Signaler sur la plateforme (1er niveau) :} Bouton « Signaler » (3 points ...) > Choisir motif exact (Harcèlement, Violence, Terrorisme, Enfant, Arnaque, Usurpation). Bloquer l'utilisateur.
    \item \textbf{Signaler aux autorités camerounaises (2ème niveau - OBLIGATOIRE pour crimes graves) :}
    \begin{itemize}
        \item \textbf{ANTIC (Agence Nationale des Technologies de l'Information et de la Communication) :} Plateforme officielle \texttt{https://signalement.antic.cm} (accessible 24/7). Remplir le formulaire : Type d'incident, URL, Captures, Description, Vos coordonnées (anonyme possible mais identifié = suivi plus efficace). \textbf{Numéros d'urgence :} \texttt{1510} (Appel gratuit) / \texttt{+237 6XX XX XX XX} (WhatsApp ANTIC).
        \item \textbf{Police Judiciaire / Gendarmerie :} Commissariat / Brigade le plus proche. Déposer plainte (PV). Apporter captures + téléphone (pour expertise si besoin). Numéro d'urgence : \texttt{117} (Police) / \texttt{113} (Gendarmerie) / \texttt{119} (Enfance en danger - Allô 119).
    \end{itemize}
    \item \textbf{Protection de la victime (Si mineur) :} Ne pas laisser l'enfant seul face à l'écran. Soutien psychologique (CMPJ - Centre Médico-Psychologique Jeunesse, associations comme \texttt{ALDEPA}, \texttt{RENATA}). Changement de mots de passe, activation 2FA (Double authentification) sur tous comptes.
\end{enumerate}
\emph{Devoir citoyen :} Signaler, c'est \textbf{protéger la prochaine victime potentielle}. Le silence = complicité passive.
\end{methode}

\begin{exoresolu}[Scénario : Réception d'une vidéo pédopornographique sur un groupe WhatsApp « Anciens 3ème »]
\textbf{Énoncé :} Mardi 22h, vous recevez sur un groupe WhatsApp d'anciens camarades (une 30aine de membres) une vidéo courte (15s) montrant clairement un acte sexuel impliquant un enfant (visage visible, estimation 8-10 ans). L'expéditeur est « Inconnu +237 6XX XX XX XX » (numéro non enregistré). Deux membres ont déjà réagi par « [emoji choc] » et « C'est quoi ça ? ». Vous êtes le seul à avoir suivi la formation éthique 3ème. Détaillez \textbf{pas à pas} vos actions immédiates et vos consignes au groupe, en citant les bases légales camerounaises.

\tcblower
\textbf{Plan d'action immédiat (Chronologie critique -- Temps visé < 5 minutes) :}

\textbf{1. ACTION IMMÉDIATE -- NE PAS OUVRIR / NE PAS TÉLÉCHARGER (Si pas déjà fait).}
\begin{itemize}
    \item Si la vidéo s'est auto-téléchargée (paramètre par défaut WhatsApp) : \textbf{Ne pas cliquer pour lire.} Aller dans \texttt{Paramètres > Stockage et données > Gérer le stockage > WhatsApp > Vidéos} et \textbf{Supprimer le fichier} du stockage local du téléphone. Vider la corbeille Galerie.
    \item \textbf{Base légale :} Détention d'images pédopornographiques = \textbf{Crime} (Art. 294 Code Pénal Camerounais : 10 à 20 ans de prison). La curiosité n'est pas une défense. \textbf{La possession même involontaire est risquée juridiquement.}
\end{itemize}

\textbf{2. PREUVE -- CAPTURE D'ÉCRAN DE L'ENVOI (PAS DE LA VIDÉO).}
\begin{itemize}
    \item Faire une capture de l'\textbf{interface WhatsApp} montrant : Le numéro expéditeur (+237 6XX XX XX XX), la vignette floutée/icône vidéo, la date/heure (22h), le nom du groupe, les réactions « [emoji choc] ».
    \item \textbf{Ne JAMAIS} faire de capture du contenu vidéo pédopornographique lui-même (c'est recréer une image illicite).
\end{itemize}

\textbf{3. SIGNALEMENT PLATEFORME (WhatsApp / Meta).}
\begin{itemize}
    \item Maintenir appuyé sur le message > \texttt{Signaler} > \texttt{Exploitation sexuelle d'enfants} (Child Sexual Abuse Material).
    \item Cocher « Bloquer le contact » et « Supprimer la discussion ».
\end{itemize}

\textbf{4. SIGNALEMENT AUTORITÉS CAMEROUNAISES (ANTIC -- PRIORITAIRE).}
\begin{itemize}
    \item Ouvrir navigateur : \texttt{https://signalement.antic.cm} (ou App « ANTIC Signalement » si installée).
    \item Catégorie : \textbf{« Contenu pédopornographique / Exploitation sexuelle des enfants en ligne »}.
    \item Champs obligatoires : URL du groupe (lien d'invitation si public) ou Description précise (Groupe « Anciens 3ème », 30 membres, numéro expéditeur +237 6XX XX XX XX, Heure 22h00, Vidéo 15s enfant ~8-10 ans).
    \item Joindre : Capture d'écran de l'interface (Étape 2).
    \item Coordonnées : Donner son identité réelle (Nom, Téléphone, École) -- \textbf{crucial} pour la traçabilité de l'expéditeur par la police judiciaire (réquisition opérateur MTN/Orange/Nexttel sur le numéro +237 6XX XX XX XX).
    \item \textbf{Appeler le 1510 (ANTIC) en parallèle} pour signaler l'urgence (vidéo en circulation active).
\end{itemize}

\textbf{5. GESTION DU GROUPE WHATSAPP (Protection des autres / Preuve collective).}
\begin{itemize}
    \item \textbf{Message type à poster IMMÉDIATEMENT dans le groupe :}
\begin{lstlisting}[language=text]
[STOP] ALERTE SÉCURITÉ MAJEURE -- ACTION REQUISE IMMÉDIATE

Une vidéo PÉDOPORNOGRAPHIQUE (enfant ~8-10 ans) vient d'être envoyée par le numéro +237 6XX XX XX XX.

[SIRENE] CONSIGNES STRICTES (LOI CAMEROUNAISE) :
1. NE REGARDEZ PAS LA VIDÉO. NE LA TRANSFÉREZ JAMAIS.
2. SUPPRIMEZ-LA DE VOTRE TÉLÉPHONE (Galerie > WhatsApp Vidéos > Corbeille).
3. NE RÉPONDEZ PAS À L'EXPÉDITEUR. NE L'INSULTEZ PAS (alertes les réseaux).
4. SIGNALEZ SUR WHATSAPP : Message > Signaler > Exploitation enfant.
5. SIGNALEZ À L'ANTIC : https://signalement.antic.cm ou appelez le 1510 (gratuit).

[JUSTICE] RAPPEL LÉGAL : Détention / Diffusion = CRIME (10-20 ans prison, Art. 294 Code Pénal + Loi 2010/012). Vous êtes témoins. Protégez-vous : signalez.

[Votre Nom] -- Référent Numérique 3ème -- 22h05
\end{lstlisting}
    \item \textbf{Quitter le groupe} immédiatement après le message (pour ne plus recevoir de contenu similaire).
    \item \textbf{Contacter l'admin du groupe} en privé : Exiger la suppression du message par l'admin (fonction « Supprimer pour tout le monde » -- si < 2h) et le bannissement de l'expéditeur + restriction « Seuls les admins peuvent envoyer des messages » le temps de l'enquête.
\end{itemize}

\textbf{6. SUIVI PERSONNEL :}
\begin{itemize}
    \item Changer mot de passe WhatsApp (vérification en 2 étapes : \texttt{Paramètres > Compte > Validation en deux étapes}).
    \item Vérifier appareils liés (\texttt{Appareils liés} -- déconnecter les inconnus).
    \item En parler à un adulte de confiance (Parent, Professeur principal, CPE) pour décharge émotionnelle.
\end{itemize}

\textbf{Conclusion :} L'éthique numérique face à l'horreur, c'est \textbf{l'action froide, légale, technique et solidaire}. Vous avez protégé 29 autres personnes d'un traumatisme et d'un risque pénal, et vous avez donné à la police la piste (numéro de téléphone) pour arrêter un criminel. C'est ça, la citoyenneté numérique en 3ème.
\end{exoresolu}

\courssub{Synthèse et auto-évaluation citoyenne}

\begin{methode}[Auto-évaluation citoyenne numérique (Grille de positionnement trimestrielle)]
\textbf{Objectif :} Mesurer sa progression sur les savoir-être (Empathie, Responsabilité, Légalité, Vigilance) pour le bulletin / portfolio.
\begin{enumerate}
    \item \textbf{Grille à 4 niveaux (Non acquis / Fragile / Acquis / Expert) sur 12 items clés :}
    \begin{center}
    \begin{tabularx}{\linewidth}{|>{\bfseries}l|X|X|X|X|}
    \hline
    \textbf{Compétence (Savoir-être)} & \textbf{Non acquis (1)} & \textbf{Fragile (2)} & \textbf{Acquis (3)} & \textbf{Expert (4)} \\
    \hline
    1. \textbf{Empathie préventive} : Je réfléchis à l'impact sur autrui \emph{avant} de publier/partager. & Je publie sans réfléchir. & Je réfléchis parfois. & \textbf{Je me mets systématiquement à la place de l'autre (règle 3 questions).} & J'anticipe les impacts indirects (viralité, contexte futur) et j'alerte mes pairs. \\
    \hline
    2. \textbf{Respect vie privée / Droit à l'image} : Je demande l'accord (photo, data, localisation). & Je poste/tague sans demander. & Je demande pour les photos « sensibles ». & \textbf{Je demande toujours (explicite), je floute les tiers, je refuse qu'on me tague sans accord.} & J'éduque mon entourage (famille, amis) sur le consentement numérique. \\
    \hline
    3. \textbf{Lutte contre cyberharcèlement} : Je ne participe pas, je signale, je soutiens la victime. & Je like/partage les moqueries / Je suis spectateur passif. & Je ne participe pas mais je n'ose pas signaler. & \textbf{Je signale (plateforme + adulte/ANTIC), j'envoie un message de soutien privé à la victime.} & J'initie une médiation / J'anime une sensibilisation en classe. \\
    \hline
    4. \textbf{Vérification info (Fake news)} : J'applique S.T.O.P. avant de partager. & Je partage tout ce qui me révolte/amuse. & Je vérifie seulement si ça a l'air bizarre. & \textbf{J'applique S.T.O.P. systématiquement (Source, Temps, Origine, Partage).} & Je produis des « Fact-checks » pour mes groupes (débunkage constructif). \\
    \hline
    5. \textbf{Protection données / Sécurité} : Mots de passe forts uniques, 2FA, MAJ, phishing. & Mdp « 123456 » partout, pas de 2FA, clique sur tout lien. & Mdp corrects, 2FA sur mail principal. & \textbf{Gestionnaire de mots de passe (Bitwarden), 2FA partout, MAJ auto, détection phishing experte.} & J'aide ma famille/voisins à sécuriser leurs comptes (aidant numérique). \\
    \hline
    6. \textbf{Propriété intellectuelle / Citation} : Je cite mes sources, j'utilise licences libres (CC). & Copier-coller sans source (plagiat). & Je mets le lien parfois. & \textbf{Citation norme APA/MLA systématique, recherche images CC0/Unsplash, respect logiciels.} & Je crée mes propres ressources (photos, code, schémas) et les partage en licence libre. \\
    \hline
    7. \textbf{Gestion conflit numérique} : Dialogue privé, OSBD, réparation, non-escalade. & Insultes publiques, vengeance, blocage sans explication. & Évite le conflit ou explose. & \textbf{Dialogue privé (vocal), OSBD, accord écrit, réparation.} & Je sers de médiateur pair formé pour mes camarades. \\
    \hline
    8. \textbf{Empreinte / E-réputation} : Test « Grand-Père/Recruteur/Juge » avant publication. & Profil public, contenu risqué (alcool, insultes, geoloc). & Profil privé, contenu neutre. & \textbf{Profil privé, contenu positif/constructif, nettoyage semestriel, Google Alert sur son nom.} & Portfolio numérique pro (LinkedIn/Behance/GitHub) commencé, veille sur son nom. \\
    \hline
    9. \textbf{Signalement contenus illicites} : Connaît et utilise ANTIC (1510), Police (117), Allô 119. & Ne sait pas signaler / Peur de se mêler des affaires des autres. & Signale sur la plateforme seulement. & \textbf{Signale plateforme + ANTIC/Police (avec captures horodatées), suit la procédure.} & Forme d'autres élèves à la procédure de signalement officiel. \\
    \hline
    10. \textbf{Écologie numérique} : Nettoyage mails/cloud, extinction veille, réparation/recyclage. & Boîte mail 5000 non lus, tel changé tous les ans, streaming 4K 4G. & Supprime spams parfois, éteint tel la nuit. & \textbf{Nettoyage mensuel (Cleanfox), stockage local > cloud, WiFi > 4G, tel 4+ ans, réparation.} & Sensibilise l'école au « Numérique Responsable » (Green IT), organise collecte DEEE. \\
    \hline
    11. \textbf{Utilisation éthique de l'IA} : Déclare l'usage, vérifie hallucinations, ne triche pas (devoirs). & Copie-colle ChatGPT sans lire pour les devoirs. & Utilise IA pour idées, recopie à la main. & \textbf{Utilise IA comme tuteur (explications), cite « Assisté par IA », vérifie sources, ne rend pas du code IA non compris.} & Crée des prompts éthiques, teste les biais de l'IA, comprend les limites (data training). \\
    \hline
    12. \textbf{Engagement citoyen positif} : Utilise le numérique pour aider, créer, s'engager (bénévolat, open source). & Usage purement consommateur (TikTok, Jeux, Scroll infini). & Partage parfois des infos utiles. & \textbf{Crée du contenu utile (tutos, alertes météo, entraide devoirs), participe à des projets solidaires en ligne.} & Porte un projet numérique à impact social (ex: Appli signalement insalubrité quartier, Wiki école). \\
    \hline
    \end{tabularx}
    \end{center}
    \item \textbf{Calcul du score :} Somme des 12 items (Min 12 -- Max 48).
    \begin{itemize}
        \item 12-23 : \textbf{Niveau 1 -- Novice} (Besoin d'accompagnement renforcé).
        \item 24-35 : \textbf{Niveau 2 -- Intermédiaire} (Bases acquises, automatismes à fixer).
        \item 36-41 : \textbf{Niveau 3 -- Compétent / Acquis BEPC} (Profil citoyen numérique solide).
        \item 42-48 : \textbf{Niveau 4 -- Expert / Ambassadeur} (Référent pair, porteur de projets).
    \end{itemize}
    \item \textbf{Plan d'action personnel :} Choisir \textbf{2 items « Fragiles »} à transformer en « Acquis » d'ici le prochain trimestre. Noter action concrète + date échéance + personne ressource.
\end{enumerate}
\end{methode}

\begin{exoresolu}[Auto-évaluation de fin de séquence -- Cas « Kevin » (Exercice 1) et « Stéphanie » (Exercice 5)]
\textbf{Énoncé :} Kevin (Exo 1 : photo humiliante d'Amina) et Stéphanie (Exo 5 : photo anniversaire alcool/insultes) passent l'auto-évaluation ci-dessus \textbf{avant} les faits, puis \textbf{après} les faits et la médiation/le conseil. Remplissez leur grille (niveau 1 à 4) pour les items 1, 2, 3, 8, 12 et calculez leur score. Commentez l'évolution.

\tcblower
\textbf{Grille Kevin (Auteur de l'humiliation)}

\begin{center}
\begin{tabularx}{\linewidth}{|l|c|c|X|}
\hline
\textbf{Item} & \textbf{Avant (Inconscient)} & \textbf{Après (Post-médiation Exo 1)} & \textbf{Commentaire} \\
\hline
1. Empathie préventive & 1 (Ne pense pas à Amina) & 3 (Comprend la douleur, règle 3 questions intégrée) & \textbf{+2} : Prise de conscience majeure via médiation OSBD. \\
\hline
2. Droit à l'image & 1 (Ignore le consentement) & 3 (Sait qu'il faut accord explicite, loi 2010/012) & \textbf{+2} : Apport juridique concret (sanction risquée). \\
\hline
3. Cyberharcèlement & 1 (Auteur actif) & 3 (Ne refera plus, sait signaler, soutient victime) & \textbf{+2} : Passage d'acteur nuisible à protecteur potentiel. \\
\hline
8. E-réputation & 2 (Profil public, contenu « drôle » risqué) & 3 (Nettoyage fait, profil privé, contenu positif) & \textbf{+1} : Action corrective immédiate (suppression, excuses). \\
\hline
12. Engagement positif & 1 (Consommateur / Blagueur) & 2 (S'excuse publiquement = 1er acte réparateur) & \textbf{+1} : Début d'engagement (réparation). Objectif : devenir médiateur pair. \\
\hline
\textbf{SOUS-TOTAL (5 items)} & \textbf{6 / 20} & \textbf{14 / 20} & \textbf{Progression +8 pts} \\
\hline
\textbf{Projection Totale (12 items est.)} & \textbf{~18/48 (Niv 1 Novice)} & \textbf{~32/48 (Niv 2 Intermédiaire $\rightarrow$ Acquis)} & \textbf{Trajectoire positive forte si accompagnement maintenu.} \\
\hline
\end{tabularx}
\end{center}

\textbf{Grille Stéphanie (Projet publication risquée -- Exo 5)}

\begin{center}
\begin{tabularx}{\linewidth}{|l|c|c|X|}
\hline
\textbf{Item} & \textbf{Avant (Projet initial)} & \textbf{Après (Conseil Test 3 Regards)} & \textbf{Commentaire} \\
\hline
1. Empathie préventive & 1 (Pense aux likes, pas aux aînés/futur soi) & 4 (Intègre le regard Grand-Père/Recruteur/Juge comme boussole) & \textbf{+3} : Internalisation de l'outil méta-cognitif. \\
\hline
2. Droit à l'image & 2 (Taguerait amis sans demander) & 4 (Taguera seulement après accord oral, floutage tiers) & \textbf{+2} : Application rigoureuse du consentement. \\
\hline
3. Cyberharcèlement & 2 (Insulte « haters » = harcèlement latent) & 4 (Supprime l'insulte, message bienveillant « Gratitude ») & \textbf{+2} : Transformation langage haine $\rightarrow$ langage constructif. \\
\hline
8. E-réputation & 1 (Profil public, contenu à haut risque pénal/pro) & 4 (Profil privé, contenu « portfolio positif », nettoyage planifié) & \textbf{+3} : Basculement stratégique majeur (protection avenir). \\
\hline
12. Engagement positif & 1 (Consommatrice / Exhibitionnisme) & 3 (Crée un contenu positif, inspire pairs par l'exemple) & \textbf{+2} : Publication devient acte citoyen (modèle). \\
\hline
\textbf{SOUS-TOTAL (5 items)} & \textbf{7 / 20} & \textbf{19 / 20} & \textbf{Progression +12 pts (Spectaculaire)} \\
\hline
\textbf{Projection Totale (12 items est.)} & \textbf{~20/48 (Niv 1 Novice)} & \textbf{~40/48 (Niv 3 Compétent $\rightarrow$ Expert)} & \textbf{Le « Test 3 Regards » agit comme un levier méta-cognitif puissant.} \\
\hline
\end{tabularx}
\end{center}

\textbf{Synthèse comparative :}
\begin{itemize}
    \item \textbf{Kevin} : Progression par \textbf{réparation} (choc de la conséquence réelle sur victime + cadre légal + médiation). Besoin de consolidation dans la durée (risque de récidive si non suivi).
    \item \textbf{Stéphanie} : Progression par \textbf{anticipation} (outil intellectuel « Test 3 Regards » appliqué \emph{avant} l'acte). Intégration plus rapide, plus stable, transférable à toutes situations futures.
    \item \textbf{Enseignement :} L'éthique \textbf{préventive} (outils de décision comme S.T.O.P., Test 3 Regards, ODARJ) est plus efficace et moins coûteuse que l'éthique \textbf{curative} (gestion de crise, médiation, sanctions). L'objectif du cours de 3ème est de donner ces outils \textbf{avant} le passage au lycée.
\end{itemize}
\end{exoresolu}

\begin{attention}[Les 5 erreurs qui coûtent des points au BEPC (Éthique / Empathie / Citoyenneté)]
\begin{enumerate}
    \item \textbf{Confondre « Éthique » et « Légal » (ou « Technique »).}
    \begin{quote}
    \textit{Erreur :} « C'est éthique car c'est pas interdit par la loi » OU « C'est éthique car j'ai le droit technique de le faire (copier-coller, screen) ». \\
    \textit{Correction :} L'éthique \textbf{dépasse} la loi (ex: la loi ne t'interdit pas de ne pas tenir la porte, l'éthique te dit de la tenir). L'éthique guide \textbf{là où la loi est silencieuse ou en retard} (ex: IA générative, deepfakes). Au BEPC, justifie toujours par une \textbf{valeur} (respect, dignité, justice) + une \textbf{conséquence} (sur autrui, sur soi, sur la société), pas seulement par un article de loi.
    \end{quote}
    \item \textbf{Réponse « Morale » sans argumentation technique/juridique (Trop vague).}
    \begin{quote}
    \textit{Erreur :} « C'est pas bien », « C'est méchant », « Il faut être gentil », « Les parents vont gronder ». \\
    \textit{Correction :} Utilise le \textbf{vocabulaire précis du programme} : \cle{Cyberharcèlement}, \cle{Droit à l'image}, \cle{Données personnelles}, \cle{Propriété intellectuelle}, \cle{Empreinte numérique}, \cle{E-réputation}, \cle{Loi 2010/012}, \cle{ANTIC}, \cle{Phishing}, \cle{Deepfake}, \cle{Biais algorithmique}. Cite la \textbf{méthode} utilisée (ODARJ, S.T.O.P., Test 3 Regards, OSBD). Structure : \textbf{Constat -> Qualification juridique/éthique -> Conséquences -> Action citoyenne}.
    \end{quote}
    \item \textbf{Oublier la dimension « Empathie / Victime » (Réponse centrée sur l'auteur ou la technique).}
    \begin{quote}
    \textit{Erreur :} Analyser longuement la technique du deepfake ou la faille de sécurité, mais ne pas mentionner la souffrance de la personne visée, l'atteinte à sa dignité, son droit à l'oubli. \\
    \textit{Correction :} Commencez TOUJOURS par l'\textbf{humain} : « La victime subit... », « Cela porte atteinte à sa vie privée car... », « L'empathie impose de... ». Le savoir-être « Adopter une attitude citoyenne vis-à-vis d'autrui » est le cœur de la leçon.
    \end{quote}
    \item \textbf{Proposer des solutions irréalistes, disproportionnées ou illégales.}
    \begin{quote}
    \textit{Erreurs classiques :} « Il faut bannir Internet au Cameroun », « Il faut hacker le compte du harceleur pour supprimer la photo », « Il faut le frapper / l'exclure à vie de l'école », « Je ne fais rien car c'est pas mon problème », « Je signale à la police pour un simple malentendu entre amis ». \\
    \textit{Correction :} Solutions \textbf{graduées, proportionnées, légales, réparatrices}. Niveau 1 : Dialogue / Suppression / Signalement plateforme. Niveau 2 : Médiation adulte / Sanction éducative / Signalement ANTIC (si infraction). Niveau 3 : Justice (si crime/délit grave). Respecte la \cle{subsidiarité} (on règle au plus bas niveau possible).
    \end{quote}
    \item \textbf{Négliger la rédaction / l'orthographe / la structure dans l'exercice de « Rédaction justifiée ».}
    \begin{quote}
    \textit{Erreur :} Écrire en langage SMS (« pk tu fé sa ? »), phrases sans verbes, pas de paragraphes, pas de connecteurs logiques (donc, cependant, par conséquent), conclusion absente. \\
    \textit{Correction :} Le savoir-faire « Rédiger et justifier une démarche, une réponse ou une conclusion » est \textbf{évalué sur la forme AUSSI}. Utilisez : \textbf{Introduction} (Rappel faits + Problématique), \textbf{Développement structuré} (Arguments numérotés 1/2/3 avec mots-liens), \textbf{Conclusion} (Décision + Valeur citoyenne). Relisez-vous. Un argumentaire illisible = 0 point même si l'idée est bonne.
    \end{quote}
\end{enumerate}
\end{attention}

\newpage

\coursec{Exercices}

\courssub{Je vérifie mes connaissances}

Avant de t'engager dans les situations complexes de cette leçon, assure-toi que les notions de base de l'unité (éthique, empathie, citoyenneté numérique, consentement, propriété intellectuelle, empreinte numérique) sont bien acquises. Les exercices qui suivent sont un contrôle rapide : réponds seul(e), puis compare avec le corrigé.

\begin{exercice}[QCM : Repérage de comportements numériques]
Pour chacun des huit comportements ci-dessous, coche la case correspondante (\textbf{Éthique} ou \textbf{Non éthique}) et justifie brièvement ton choix pour les items 3 et 6.

\begin{center}
\begin{tabularx}{\linewidth}{|c|X|c|c|}\hline
\rowcolor{popblueL}\textbf{N°} & \textbf{Comportement} & \textbf{Éthique} & \textbf{Non éthique} \\\hline
1 & Je demande l'autorisation avant de publier une photo de mon camarade. & $\square$ & $\square$ \\\hline
2 & J'utilise le mot de passe d'un ami pour accéder à sa messagerie « pour rigoler ». & $\square$ & $\square$ \\\hline
3 & Je partage une rumeur non vérifiée sur un groupe WhatsApp de la classe. & $\square$ & $\square$ \\\hline
4 & Je signale un contenu haineux vu sur un réseau social à la modération. & $\square$ & $\square$ \\\hline
5 & Je copie-colle un devoir trouvé sur Internet sans citer la source. & $\square$ & $\square$ \\\hline
6 & Je refuse de transmettre une capture d'écran d'une conversation privée. & $\square$ & $\square$ \\\hline
7 & J'installe un logiciel piraté sur l'ordinateur du cybercafé du quartier. & $\square$ & $\square$ \\\hline
8 & J'aide une personne âgée à configurer la visioconférence pour parler à sa famille. & $\square$ & $\square$ \\\hline
\end{tabularx}
\end{center}

\textbf{Justification item 3 :} \dotfill \\[1cm]
\textbf{Justification item 6 :} \dotfill
\end{exercice}

\begin{exercice}[Vrai ou Faux justifié]
Indique si chaque affirmation est \textbf{Vraie} ou \textbf{Fausse} et justifie ta réponse en une phrase en t'appuyant sur le vocabulaire du cours (consentement, propriété intellectuelle, empreinte numérique, empathie).

\begin{enumerate}
\item « Tout ce qui est publié sur Internet est libre de droit, je peux l'utiliser comme je veux. » \hfill [Vrai / Faux] \\
Justification : \dotfill
\item « L'empathie numérique consiste à se mettre à la place de l'autre avant de poster un commentaire. » \hfill [Vrai / Faux] \\
Justification : \dotfill
\item « Supprimer une photo de mon profil efface définitivement mon empreinte numérique. » \hfill [Vrai / Faux] \\
Justification : \dotfill
\item « Le consentement pour une photo de groupe est implicite si personne ne proteste. » \hfill [Vrai / Faux] \\
Justification : \dotfill
\end{enumerate}
\end{exercice}

\begin{exercice}[Définitions essentielles]
Donne une définition précise et complète (2 à 3 lignes) pour chacun des trois concepts suivants en utilisant les termes techniques vus en classe :

\begin{enumerate}
\item \textbf{Citoyenneté numérique} : \dotfill \dotfill \dotfill
\item \textbf{Empathie (dans un contexte numérique)} : \dotfill \dotfill \dotfill
\item \textbf{Consentement éclairé (pour une donnée personnelle)} : \dotfill \dotfill \dotfill
\end{enumerate}
\end{exercice}

\begin{exercice}[Association : Situation $\leftrightarrow$ Principe violé ou respecté]
Relie chaque situation de la colonne de gauche au principe éthique correspondant dans la colonne de droite (un principe peut être utilisé plusieurs fois).

\begin{center}
\begin{tabularx}{\linewidth}{|X|c|X|}\hline
\rowcolor{popblueL}\textbf{Situation} & \textbf{N°} & \textbf{Principe éthique} \\\hline
A. Je télécharge un film sur un site de streaming illégal. & 1. & Respect de la propriété intellectuelle \\\hline
B. Je floute le visage d'un inconnu sur ma photo de voyage avant de poster. & 2. & Protection de la vie privée / Droit à l'image \\\hline
C. J'insulte un camarade dans les commentaires d'une publication. & 3. & Respect d'autrui / Non-violence (Netiquette) \\\hline
D. Je cite l'auteur et le site source pour un paragraphe utilisé dans mon exposé. & 4. & Honnêteté intellectuelle / Citation des sources \\\hline
E. Je signale un compte qui usurpe l'identité de mon professeur. & 5. & Responsabilité / Signalement (Citoyenneté active) \\\hline
F. Je garde pour moi le secret confié par un ami sur la messagerie. & 6. & Confidentialité / Loyauté \\\hline
\end{tabularx}
\end{center}
\end{exercice}

\courssub{Je m'entraîne}

\begin{exercice}[Analyse de cas : Diffusion d'une photo sans consentement]
Léa a pris en photo son camarade Max endormi en cours. Elle a publié la photo sur son statut WhatsApp avec la légende « Le roi de la sieste » suivie d'un smiley moqueur, sans demander l'avis de Max. La photo a été vue par 80 contacts, dont des élèves d'autres classes. Max l'apprend le soir même et se sent humilié.

\begin{enumerate}
\item Identifie \textbf{deux conduites non éthiques} commises par Léa dans cette situation.
\item Pour chaque conduite, cite \textbf{une notion du cours} (vocabulaire précis) qui permet de la qualifier.
\item Explique en deux lignes en quoi l'\textbf{empathie} aurait dû modifier le comportement de Léa \emph{avant} la publication.
\end{enumerate}
\end{exercice}

\begin{exercice}[Rédaction argumentée : Le dilemme du devoir copié]
\textbf{Situation :} Ton ami Kevin t'envoie un lien vers un site qui propose le corrigé intégral du devoir de maths à rendre demain. Il t'écrit : « Personne ne saura, on aura tous la même note, c'est de l'entraide ! » Tu as déjà commencé ton devoir mais tu bloques sur l'exercice 3.

\textbf{Consigne :} Rédige une réponse argumentée à Kevin (8 à 10 lignes) en suivant la structure \textbf{TAEC} :
\begin{itemize}
\item \textbf{T}hèse (ta position claire : tu acceptes ou tu refuses),
\item \textbf{A}rgument (la raison éthique principale : honnêteté, équité, respect du travail...),
\item \textbf{E}xemple (conséquence concrète pour toi, pour lui, pour la classe),
\item \textbf{C}onclusion (proposition alternative constructive).
\end{itemize}
\end{exercice}

\begin{exercice}[Rédaction : Charte d'utilisation du téléphone en classe]
En tant que délégué(e) de ta classe de 3ème, tu dois proposer une \textbf{Charte de bonne conduite pour l'utilisation du téléphone portable en classe} (5 règles maximum). Pour chaque règle, indique la justification éthique (respect, attention, vie privée...) ou empathique (gêne pour l'autre, distraction...).

\begin{center}
\begin{tabularx}{\linewidth}{|c|X|X|}\hline
\rowcolor{popblueL}\textbf{N°} & \textbf{Règle formulée positivement} & \textbf{Justification (Éthique / Empathie)} \\\hline
1 & & \\\hline
2 & & \\\hline
3 & & \\\hline
4 & & \\\hline
5 & & \\\hline
\end{tabularx}
\end{center}
\end{exercice}

\begin{exercice}[Auto-évaluation citoyenne : Grille de conduite du mois]
Remplis honnêtement la grille ci-dessous pour ton comportement numérique du mois écoulé (Jamais / Parfois / Souvent / Toujours). Puis formule un \textbf{engagement personnel} justifié pour le mois prochain.

\begin{center}
\begin{tabularx}{\linewidth}{|X|c|c|c|c|}\hline
\rowcolor{popblueL}\textbf{Critère citoyen} & \textbf{Jamais} & \textbf{Parfois} & \textbf{Souvent} & \textbf{Toujours} \\\hline
1. Je vérifie l'info avant de la partager. & $\square$ & $\square$ & $\square$ & $\square$ \\\hline
2. Je demande l'avis avant de poster une photo d'autrui. & $\square$ & $\square$ & $\square$ & $\square$ \\\hline
3. J'utilise un langage respectueux en ligne (pas d'insultes). & $\square$ & $\square$ & $\square$ & $\square$ \\\hline
4. Je cite mes sources pour les devoirs/recherches. & $\square$ & $\square$ & $\square$ & $\square$ \\\hline
5. Je protège mes mots de passe et ceux des autres. & $\square$ & $\square$ & $\square$ & $\square$ \\\hline
6. Je signale les contenus haineux ou harcelants. & $\square$ & $\square$ & $\square$ & $\square$ \\\hline
7. J'éteins mes notifications en cours / en réunion. & $\square$ & $\square$ & $\square$ & $\square$ \\\hline
8. J'aide un camarade en difficulté numérique. & $\square$ & $\square$ & $\square$ & $\square$ \\\hline
9. Je refuse de participer à un groupe qui critique quelqu'un. & $\square$ & $\square$ & $\square$ & $\square$ \\\hline
10. Je réfléchis à l'impact de mes mots avant d'envoyer. & $\square$ & $\square$ & $\square$ & $\square$ \\\hline
\end{tabularx}
\end{center}

\textbf{Mon engagement personnel pour le mois prochain (1 action concrète) :} \dotfill \\
\textbf{Justification (lien avec un critère faible) :} \dotfill
\end{exercice}

\courssub{Je me prépare au BEPC}

\begin{exercice}[Cas complet : La rumeur sur l'enseignant — Méthode ODARJ]
\textbf{Contexte :} Dans le groupe WhatsApp « 3ème A Officiel » (120 élèves), un élève poste une capture d'écran tronquée d'une note de service du proviseur. Il écrit : « M. Kamga va nous faire redoubler exprès, il a dit en réunion qu'il veut baisser la moyenne de la classe ! » La rumeur enfle, des insultes fusent contre l'enseignant. Le proviseur convoque le délégué de classe.

\textbf{Consigne :} Applique la méthode \textbf{ODARJ} (Observer, Définir, Analyser, Réagir, Juger/Réparer) pour analyser la situation et proposer une issue constructive.

\begin{enumerate}
\item \textbf{Observer :} Liste trois faits objectifs (vérifiables) et trois éléments subjectifs (opinions, émotions) présents dans le message initial.
\item \textbf{Définir :} Qualifie le problème éthique principal (ex : diffamation, atteinte à la réputation, propagation de fausses informations...).
\item \textbf{Analyser :} Identifie \textbf{trois conséquences} potentielles (pour l'enseignant, pour le climat de classe, pour l'auteur du post) si rien n'est fait.
\item \textbf{Réagir :} Propose \textbf{deux actions immédiates} concrètes que le délégué peut faire \emph{sur le groupe} et \emph{auprès de l'administration}.
\item \textbf{Juger / Réparer :} Rédige le message exact (3 à 4 lignes) que le délégué doit poster dans le groupe pour \textbf{réparer} le tort (excuses, rappel à la règle, appel au calme) sans envenimer la situation.
\end{enumerate}
\end{exercice}

\begin{exercice}[Situation-problème : Cyberharcèlement et responsabilité du groupe]
\textbf{Contexte :} Depuis trois semaines, une élève de ta classe, Amina, reçoit des messages anonymes sur une application de messagerie éphémère : « Personne ne t'aime », « Tu es nulle », « Fais-nous plaisir, change d'école ». Elle ne vient plus en cours, elle a maigri. Certains élèves du groupe « 3ème A Blagues » (45 membres) savent qui sont les auteurs (deux élèves influents de la classe) mais n'osent rien dire. Le professeur principal remarque le changement et lance une enquête anonyme.

\begin{enumerate}
\item \textbf{Qualification juridique et éthique :} En te basant sur la loi camerounaise (loi n° 2010/012 du 21 décembre 2010 relative à la cybersécurité et à la cybercriminalité) et sur l'éthique, nomme l'infraction commise par les auteurs et la faute éthique commise par les membres du groupe qui savent et se taisent.
\item \textbf{Analyse des rôles :} Explique la différence de responsabilité éthique entre : \textbf{l'auteur}, \textbf{le suiveur (qui like ou rit)}, \textbf{le témoin silencieux}, \textbf{le lanceur d'alerte}.
\item \textbf{Empathie active :} Si tu étais témoin silencieux, décris la démarche concrète en \textbf{3 étapes} (inspirée d'ODARJ) pour intervenir sans te mettre en danger.
\item \textbf{Mesure de réparation collective :} Propose une action concrète à mener par la classe entière (encadrée par le professeur principal) pour rétablir un climat de confiance et réparer le préjudice subi par Amina.
\end{enumerate}
\end{exercice}

\begin{exercice}[Cas technique et éthique : Données personnelles et « petits jeux » en ligne]
\textbf{Contexte :} Un jeu en ligne très populaire au quartier, « Quiz Culture G », demande à l'inscription : Nom, Prénom, Date de naissance, Numéro de téléphone, Localisation GPS précise, Accès aux contacts du téléphone, Accès à la caméra. Les CGU (Conditions Générales d'Utilisation), longues de 15 pages en petits caractères, mentionnent à l'article 12.4 : « L'utilisateur cède gratuitement tous droits sur ses données à la société X, basée à l'étranger, pour exploitation commerciale illimitée. » Ton petit cousin de 13 ans veut s'inscrire pour gagner « 10 000 FCFA de crédit téléphonique ».

\begin{enumerate}
\item \textbf{Analyse des CGU :} Identifie \textbf{trois clauses abusives ou dangereuses} pour un mineur dans la description ci-dessus (réfère-toi aux principes : proportionnalité, consentement libre et éclairé, minimisation des données, droit à l'oubli).
\item \textbf{Risques concrets :} Décris deux risques réels pour ton cousin s'il accepte (ex : usurpation d'identité, démarchage agressif, géolocalisation par des inconnus mal intentionnés, revente du numéro...).
\item \textbf{Conseil technique et éthique :} Rédige le message (5 à 6 lignes) que tu envoies à ton cousin pour lui expliquer \textbf{pourquoi il ne doit pas s'inscrire} et \textbf{comment s'amuser autrement} en toute sécurité (alternatives : jeux hors ligne, plateformes reconnues, paramétrage de la vie privée).
\item \textbf{Droit à l'oubli :} Si ton cousin s'est déjà inscrit et a donné son numéro, explique la démarche à suivre (auprès de qui ? avec quelles précautions ?) pour demander la suppression de ses données.
\end{enumerate}
\end{exercice}

\begin{aretenir}[L'essentiel de la leçon]
\begin{itemize}
\item L'\textbf{éthique numérique} commande de respecter autrui en ligne comme en face : consentement, honnêteté, confidentialité, propriété intellectuelle.
\item L'\textbf{empathie} est le réflexe préalable à toute publication : « Comment l'autre vivrait-il ce message ? »
\item La méthode \textbf{ODARJ} (Observer, Définir, Analyser, Réagir, Juger/Réparer) permet d'analyser toute situation d'éthique et de construire une réponse justifiée — une compétence directement évaluée au BEPC.
\item Être un \textbf{citoyen numérique}, c'est aussi agir : vérifier avant de partager, signaler les abus, protéger ses données et celles des autres, réparer le tort causé.
\end{itemize}
\end{aretenir}

\newpage

\coursec{Corrigés des exercices}

\begin{corrige}[Exercice 1 — QCM : Repérage de comportements numériques]
\begin{enumerate}
\item \textbf{Éthique} : le \cle{consentement} de la personne photographiée est respecté (\cle{droit à l'image}).
\item \textbf{Non éthique} : accéder au compte d'autrui, même « pour rigoler », est une violation de la \cle{vie privée} et une infraction (accès frauduleux à un système d'information).
\item \textbf{Non éthique} : propager une information non vérifiée, c'est participer à la \cle{désinformation} ; la rumeur peut nuire gravement à la réputation d'une personne.
\item \textbf{Éthique} : signaler un contenu haineux est un acte de \cle{citoyenneté numérique} active.
\item \textbf{Non éthique} : c'est du \cle{plagiat}, une violation de la \cle{propriété intellectuelle} et de l'honnêteté intellectuelle.
\item \textbf{Éthique} : refuser de diffuser une conversation privée, c'est respecter la \cle{confidentialité} et la confiance d'autrui.
\item \textbf{Non éthique} : le piratage de logiciels viole la propriété intellectuelle et expose le poste à des \cle{virus}.
\item \textbf{Éthique} : aider autrui à maîtriser l'outil numérique est un comportement solidaire et \cle{empathique}.
\end{enumerate}
\end{corrige}

\begin{corrige}[Exercice 2 — Vrai ou Faux justifié]
\begin{enumerate}
\item \textbf{Faux.} Une œuvre publiée sur Internet reste protégée par la \cle{propriété intellectuelle} : il faut l'autorisation de l'auteur ou vérifier la licence, et toujours citer la source.
\item \textbf{Vrai.} L'\cle{empathie numérique}, c'est anticiper l'effet de ses mots sur autrui : « aimerais-je recevoir ce commentaire ? »
\item \textbf{Faux.} Des copies (captures d'écran, partages, sauvegardes des serveurs) peuvent subsister : l'\cle{empreinte numérique} est difficile à effacer totalement.
\item \textbf{Faux.} Le \cle{consentement éclairé} doit être exprès et éclairé : le silence n'est pas un accord, il faut demander explicitement l'autorisation.
\end{enumerate}
\end{corrige}

\begin{corrige}[Exercice 3 — Définitions essentielles]
\begin{enumerate}
\item \textbf{\cle{Citoyenneté numérique}} : ensemble des comportements responsables, respectueux et solidaires qu'un utilisateur adopte dans l'environnement numérique (respect d'autrui, honnêteté, respect des lois, signalement des abus).
\item \textbf{\cle{Empathie numérique}} : capacité à se mettre à la place d'autrui derrière l'écran, à imaginer ses émotions et à adapter ses paroles et ses actes en ligne pour ne pas lui nuire.
\item \textbf{\cle{Consentement éclairé}} : accord donné librement, de manière explicite, par une personne qui a été clairement informée de l'usage qui sera fait de sa donnée personnelle (photo, numéro, localisation).
\end{enumerate}
\end{corrige}

\begin{corrige}[Exercice 4 — Association : Situation $\leftrightarrow$ Principe]
\emph{(Rappel : les situations A à F et les principes 1 à 6 sont définis dans l'énoncé de l'exercice.)}
A $\rightarrow$ 1 (violé) ; B $\rightarrow$ 2 (respecté) ; C $\rightarrow$ 3 (violé) ; D $\rightarrow$ 4 (respecté) ; E $\rightarrow$ 5 (respecté) ; F $\rightarrow$ 6 (respecté).
\end{corrige}

\begin{corrige}[Exercice 5 — Analyse de cas : Diffusion d'une photo sans consentement]
\begin{enumerate}
\item Deux conduites non éthiques : (a) publier la photo de Max \textbf{sans son consentement} ; (b) se \textbf{moquer} de lui publiquement devant 80 personnes.
\item (a) Violation du \cle{droit à l'image} et de la \cle{vie privée} ; (b) atteinte à la \cle{dignité} d'autrui, contraire à la \cle{Netiquette} (cela peut constituer une forme de \cle{cyberharcèlement}).
\item L'empathie aurait conduit Léa à se demander : « Comment me sentirais-je si l'on diffusait une photo ridicule de moi à 80 personnes ? » En anticipant l'humiliation de Max, elle aurait renoncé à publier ou aurait d'abord demandé son accord.
\end{enumerate}
\end{corrige}

\begin{corrige}[Exercice 6 — Rédaction argumentée : Le dilemme du devoir copié]
\textbf{Thèse :} je refuse d'utiliser le corrigé trouvé en ligne. \textbf{Argument :} recopier un corrigé, c'est du \cle{plagiat} et de la triche : cela viole l'honnêteté intellectuelle et crée une inéquité envers les camarades qui travaillent seuls. \textbf{Exemple :} si le professeur découvre des copies identiques, nous risquons tous un zéro et une sanction disciplinaire ; surtout, au BEPC, sans entraînement réel, l'échec est presque certain. \textbf{Conclusion :} je propose à Kevin de travailler ensemble l'exercice 3 ce soir, en nous expliquant mutuellement nos idées : c'est cela, la vraie entraide.
\end{corrige}

\begin{corrige}[Exercice 7 — Rédaction : Charte d'utilisation du téléphone en classe]
\begin{enumerate}
\item \textbf{Règle :} je mets mon téléphone en mode silencieux et dans mon sac pendant les cours. \textbf{Justification :} respect de l'attention du professeur et des camarades (empathie : une sonnerie déconcentre tout le monde).
\item \textbf{Règle :} je n'utilise mon téléphone qu'avec l'autorisation de l'enseignant, pour une activité scolaire. \textbf{Justification :} honnêteté et respect du cadre de la classe.
\item \textbf{Règle :} je ne photographie ni ne filme personne sans son accord. \textbf{Justification :} \cle{droit à l'image} et \cle{vie privée}.
\item \textbf{Règle :} je n'envoie aucun message moqueur ou insultant, même en dehors de la classe. \textbf{Justification :} \cle{Netiquette} et non-violence (empathie envers la victime).
\item \textbf{Règle :} je signale à un adulte tout contenu harcelant dont je suis témoin. \textbf{Justification :} \cle{citoyenneté numérique} active et solidarité.
\end{enumerate}
\end{corrige}

\begin{corrige}[Exercice 8 — Auto-évaluation citoyenne : Grille de conduite du mois]
Cet exercice est un outil d'auto-évaluation personnelle ; il n'y a pas de réponse unique, mais l'honnêteté dans le remplissage est la condition de son utilité. Voici un exemple de profil vertueux (visant le \textbf{Toujours} sur les critères de sécurité, de respect et de solidarité) et une proposition d'engagement type.

\textbf{Grille remplie (exemple de profil responsable) :}
\begin{center}
\begin{tabularx}{\linewidth}{|X|c|c|c|c|}\hline
\rowcolor{popblueL}\textbf{Critère citoyen} & \textbf{Jamais} & \textbf{Parfois} & \textbf{Souvent} & \textbf{Toujours} \\\hline
1. Je vérifie l'info avant de la partager. & & & & $\boxtimes$ \\\hline
2. Je demande l'avis avant de poster une photo d'autrui. & & & & $\boxtimes$ \\\hline
3. J'utilise un langage respectueux en ligne. & & & & $\boxtimes$ \\\hline
4. Je cite mes sources pour les devoirs/recherches. & & & & $\boxtimes$ \\\hline
5. Je protège mes mots de passe et ceux des autres. & & & & $\boxtimes$ \\\hline
6. Je signale les contenus haineux ou harcelants. & & & & $\boxtimes$ \\\hline
7. J'éteins mes notifications en cours / en réunion. & & & & $\boxtimes$ \\\hline
8. J'aide un camarade en difficulté numérique. & & & $\boxtimes$ & \\\hline
9. Je refuse de participer à un groupe qui critique quelqu'un. & & & & $\boxtimes$ \\\hline
10. Je réfléchis à l'impact de mes mots avant d'envoyer. & & & & $\boxtimes$ \\\hline
\end{tabularx}
\end{center}

\textbf{Mon engagement personnel pour le mois prochain (exemple) :} « Je m'engage à vérifier systématiquement la source de toute information (site officiel, média reconnu) avant de la partager sur les groupes de classe WhatsApp. »

\textbf{Justification (lien avec un critère faible) :} Le critère 1 (vérification de l'info) est souvent négligé par rapidité ; or la \cle{désinformation} nuit au climat de classe et peut nuire à des personnes innocentes (critères 3 et 6). Cet engagement renforce ma \cle{citoyenneté numérique} active.
\end{corrige}

\begin{corrige}[Exercice 9 — Cas complet : La rumeur sur l'enseignant — Méthode \cle{ODARJ}]
\begin{enumerate}
\item \textbf{Observer.} Faits objectifs : une capture d'écran tronquée a été publiée ; le groupe compte 120 élèves ; le message accuse nommément M. Kamga. Éléments subjectifs : l'interprétation « il veut nous faire redoubler exprès » ; la colère et la peur des élèves ; les insultes qui suivent.
\item \textbf{Définir.} Il s'agit de propagation d'une information non vérifiée (rumeur) et d'une atteinte à la réputation d'une personne, ce qui peut constituer de la \cle{diffamation}.
\item \textbf{Analyser.} Conséquences : pour l'enseignant, atteinte à son honneur et perte de confiance des élèves ; pour la classe, dégradation du climat scolaire et sanctions collectives possibles ; pour l'auteur, sanctions disciplinaires, voire poursuites (la loi n° 2010/012 du 21 décembre 2010 relative à la \cle{cybersécurité} et à la \cle{cybercriminalité} au Cameroun réprime la diffusion de fausses informations).
\item \textbf{Réagir.} Sur le groupe : demander le retrait immédiat du message et rappeler la règle de vérification des informations. Auprès de l'administration : informer le proviseur des faits exacts et proposer une rencontre de clarification entre M. Kamga et les délégués.
\item \textbf{Juger / Réparer.} Message possible : « Chers camarades, la capture publiée hier était tronquée et l'accusation contre M. Kamga n'était fondée sur aucun fait vérifié. Nous présentons nos excuses à M. Kamga au nom de la classe. Merci de ne plus relayer ce message et de vérifier toute information avant de la partager. Restons unis et respectueux. »
\end{enumerate}
\end{corrige}

\begin{corrige}[Exercice 10 — Situation-problème : Cyberharcèlement et responsabilité du groupe]
\begin{enumerate}
\item Les auteurs commettent du \cle{cyberharcèlement} : des messages répétés, injurieux et dégradants, punis par la loi n° 2010/012 relative à la cybersécurité et à la cybercriminalité. Ceux qui savent et se taisent commettent une faute éthique de \textbf{non-assistance à personne en danger moral} : leur silence encourage les auteurs et aggrave la souffrance d'Amina.
\item \textbf{L'auteur} porte la responsabilité directe des actes. \textbf{Le suiveur} (qui like ou rit) alimente le harcèlement : il est complice moralement. \textbf{Le témoin silencieux} ne participe pas mais laisse faire : sa responsabilité est moindre mais réelle. \textbf{Le lanceur d'alerte} agit pour protéger la victime : il accomplit un devoir de \cle{citoyenneté numérique}.
\item Démarche en 3 étapes : (1) \textbf{Observer et conserver les preuves} (captures d'écran datées) sans répondre aux harceleurs ; (2) \textbf{Soutenir la victime} en privé (parler à Amina, lui dire qu'elle n'est pas seule) ; (3) \textbf{Alerter un adulte de confiance} (professeur principal, parent, conseiller d'orientation) en remettant les preuves, éventuellement de façon anonyme.
\item Action collective : organiser, avec le professeur principal, une \textbf{séance de parole encadrée} où la classe adopte une charte anti-harcèlement, présente des excuses collectives à Amina et met en place un référent « écoute » (élève relais) pour que tout élève victime puisse signaler un abus en confiance.
\end{enumerate}
\end{corrige}

\begin{corrige}[Exercice 11 — Cas technique et éthique : Données personnelles et « petits jeux » en ligne]
\begin{enumerate}
\item Trois clauses abusives : (a) la collecte de la \textbf{localisation GPS précise} et de l'\textbf{accès aux contacts} est disproportionnée pour un simple quiz (violation du principe de \cle{minimisation des données}) ; (b) la \textbf{cession gratuite et illimitée} de toutes les données à une société étrangère prive l'utilisateur de tout contrôle (consentement ni libre ni éclairé, surtout pour un mineur) ; (c) des CGU de 15 pages en petits caractères rendent l'information \textbf{incompréhensible}, donc le consentement n'est pas éclairé.
\item Deux risques réels : (a) la \textbf{revente du numéro} entraîne démarchage publicitaire, arnaques par SMS ou appels frauduleux ; (b) la \textbf{géolocalisation précise} d'un enfant peut être exploitée par des personnes mal intentionnées, et ses contacts peuvent être démarchés à leur tour.
\item Message possible : « Salut petit frère, ne t'inscris pas à Quiz Culture G : le jeu réclame ta position GPS, tes contacts et ton numéro, puis revend tout à une société étrangère. Les 10 000 FCFA promis sont un piège pour voler tes données. Choisis plutôt des jeux hors ligne ou des applications connues qui ne demandent que le strict nécessaire, et règle toujours les autorisations dans les paramètres du téléphone. Si tu as un doute, montre-moi l'application avant de t'inscrire. »
\item Démarche de suppression : avec l'aide d'un parent (il est mineur), écrire au \textbf{service responsable du jeu} (adresse indiquée dans les CGU) pour demander la \textbf{suppression du compte et de toutes les données} en invoquant le \cle{droit à l'oubli} ; désinstaller l'application et révoquer ses autorisations dans les paramètres du téléphone ; conserver une preuve de la demande ; si l'entreprise ne répond pas, saisir l'autorité de protection des données compétente (au Cameroun, l'\cle{ANTIC} peut être sollicitée pour les questions de cybersécurité).
\end{enumerate}
\end{corrige}

\newpage

\coursec{Fiche bilan}

\begin{popfigure}
\begin{tikzpicture}[mindmap, grow cyclic, every node/.style={concept}, text width=2.6cm, align=center,
  root concept/.append style={concept color=popblue, font=\bfseries, text=white},
  level 1/.append style={level distance=3.4cm, sibling angle=60, font=\small},
  level 1 concept/.append style={text width=2.4cm}]
\node[root concept]{Éthique et empathie en action}
  child[concept color=poporange] { node {Rappels : éthique, empathie, citoyenneté numérique} }
  child[concept color=poppink] { node {Conduites non éthiques à repérer} }
  child[concept color=popgreen] { node {Attitude citoyenne vis-à-vis d'autrui} }
  child[concept color=poppurple] { node {Méthode ODARJ d'analyse} }
  child[concept color=popgold] { node {Rédiger et justifier une réponse} }
  child[concept color=popblue!70] { node {Synthèse et auto-évaluation} };
\end{tikzpicture}
\legende{Carte mentale de la leçon : appliquer l'éthique et l'empathie dans la vie numérique quotidienne.}
\end{popfigure}

\begin{bilan}
\textbf{Définitions clés}
\begin{itemize}
  \item \cle{Éthique} : ensemble de règles de conduite qui distinguent le bien du mal, y compris dans l'usage du numérique.
  \item \cle{Empathie} : capacité à se mettre à la place d'autrui, à comprendre ses émotions avant d'agir ou de publier.
  \item \cle{Citoyenneté numérique} : comportement responsable, respectueux et solidaire sur Internet et les réseaux sociaux.
  \item \cle{Conduite non éthique} : acte qui nuit à autrui ou enfreint les règles (cyberharcèlement, plagiat, vol de mot de passe, diffusion de fausses informations, piratage).
\end{itemize}

\textbf{Repères à connaître par c\oe ur}
\begin{center}
\begin{tabularx}{\linewidth}{|l|X|}
\hline
\rowcolor{popblueL} \textbf{Notion} & \textbf{À retenir} \\
\hline
Avant de publier & Me sentirais-je bien si c'était dit de moi ? (test de l'empathie) \\
\hline
Face à une info & Vérifier la source avant de partager \\
\hline
Travail d'autrui & Citer l'auteur, ne jamais copier sans permission \\
\hline
Données d'autrui & Ne jamais utiliser le compte ou le téléphone de quelqu'un sans accord \\
\hline
Témoin d'une injustice & Ne pas rester spectateur : alerter, soutenir la victime \\
\hline
\end{tabularx}
\end{center}

\textbf{Méthode express ODARJ (analyser une situation d'éthique)}
\begin{enumerate}
  \item \textbf{O}bserver : décrire les faits sans juger.
  \item \textbf{D}éfinir : nommer le problème éthique posé.
  \item \textbf{A}nalyser : identifier les victimes, les conséquences, les règles enfreintes.
  \item \textbf{R}éagir : proposer une conduite citoyenne et empathique.
  \item \textbf{J}ustifier : expliquer sa réponse avec des arguments clairs.
\end{enumerate}

\textbf{Rédiger une réponse argumentée}
\begin{itemize}
  \item Annoncer clairement sa position (une phrase).
  \item Donner au moins deux arguments, chacun illustré par un exemple de la vie courante.
  \item Conclure en rappelant la règle éthique ou la valeur défendue (respect, honnêteté, solidarité).
\end{itemize}
\end{bilan}

\begin{aretenir}[Checklist avant le BEPC]
\begin{itemize}
  \item $\square$ Je sais définir l'éthique, l'empathie et la citoyenneté numérique avec mes propres mots.
  \item $\square$ Je sais citer au moins quatre conduites non éthiques de la vie courante (cyberharcèlement, plagiat, fausses informations, utilisation du compte d'autrui).
  \item $\square$ Je sais expliquer pourquoi chacune de ces conduites nuit à autrui.
  \item $\square$ Je sais décrire l'attitude citoyenne à adopter face à chaque situation.
  \item $\square$ Je connais les cinq étapes de la méthode ODARJ dans l'ordre.
  \item $\square$ Je sais appliquer la méthode ODARJ à une situation concrète (classe, quartier, réseaux sociaux).
  \item $\square$ Je sais rédiger une réponse argumentée : position, deux arguments, exemples, conclusion.
  \item $\square$ Je pense au test de l'empathie avant toute publication ou tout partage.
  \item $\square$ Je sais qu'un témoin d'une injustice numérique doit réagir et non rester spectateur.
  \item $\square$ Je relis ma copie pour vérifier que chaque affirmation est justifiée.
\end{itemize}
\end{aretenir}

\findecours

\end{document}
